% Compile with pdflatex (three passes to settle references). No auxiliary source, custom style,
% external bibliography, image, or data file is needed.
\documentclass[11pt]{amsart}
\usepackage[a4paper,margin=1in]{geometry}
\usepackage[T1]{fontenc}
\usepackage{lmodern,microtype,amsmath,amssymb,amsthm,booktabs,array,mathtools,longtable}
\usepackage[hidelinks]{hyperref}
\newtheorem{theorem}{Theorem}[section]
\newtheorem{lemma}[theorem]{Lemma}
\newtheorem{proposition}[theorem]{Proposition}
\newtheorem{corollary}[theorem]{Corollary}
\theoremstyle{definition}\newtheorem{definition}[theorem]{Definition}
\theoremstyle{remark}\newtheorem{remark}[theorem]{Remark}
\numberwithin{equation}{section}
\newcommand{\R}{\mathbb R}
\newcommand{\diag}{\operatorname{diag}}
\newcommand{\tr}{\operatorname{tr}}
\newcommand{\rank}{\operatorname{rank}}
\newcommand{\one}{\mathbf1}
\newcommand{\eps}{\varepsilon}
\newcommand{\T}{\mathsf T}
\setcounter{tocdepth}{1}

\newcommand{\Sc}{\mathcal S}
\newcolumntype{L}[1]{>{\raggedright\arraybackslash}p{#1}}
\setlength{\emergencystretch}{2em}
\hypersetup{pdftitle={Constructive Komlos discrepancy: finite walks and fast implementations},
 pdfsubject={Complete signing algorithms, conditioned complexity, and structured cubic implementations}}
\title[Constructive Koml\'os discrepancy]{Constructive Koml\'os discrepancy:\\finite walks and fast implementations}
\author{}
\date{28 September 2026}
\begin{document}
\begin{abstract}
We develop finite arithmetic implementations of retained-row spectral
signing for matrices with Euclidean column norms at most one. Our
strongest deterministic construction returns discrepancy below $37.54$
in $O(mn)+\widetilde O(n^4)$ ordinary field operations and comparisons,
with $O(n\log^4(2+n))$ updates. Its normalized response estimate counts
protected rows, large rows, and curvature jointly. A screened below-$67$
variant retains the imported fast-arithmetic bound
$O(mn)+\widetilde O(m_{66}n^2+n^{\omega+1})$, where $m_{66}$ counts
rows of $\ell_1$ norm greater than 66; for $m\le n$ this gives
$\widetilde O(n^{\omega+1})$ work after the input scan.
These are distinct algorithmic and arithmetic-model guarantees.
The quadratic construction remains useful for a modular randomized
implementation whose non-sampling work is $\widetilde O(mn^2+n^3)$.
Conjugate gradients and polynomial sampling give complete bounds with
explicit conditioning or enclosure dependence; specified structural
classes admit cubic work. A two-sided secant controller has a cubic
bound conditional on one unproved expected-production hypothesis.
For practical execution, we retain long endpoint steps and add bounded
sign proposals checked directly against the $37.54$ target. Matched
experiments compare runtime, output discrepancy, filtering, and terminal
rules. Output checks certify returned signings of the represented input;
the numerical code does not yet implement the exact finite fallback.
The full arithmetic and finite-step proofs are given here and in the
accompanying optimized-profile and normalized-response supplements.
\end{abstract}
\maketitle
\tableofcontents
\clearpage

\section{Introduction}
The constructive Koml\'os problem asks for a signing
$\varepsilon\in\{-1,1\}^n$ with uniformly bounded $\|A\varepsilon\|_\infty$
whenever the columns of $A\in\mathbb R^{m\times n}$ have Euclidean norm
at most one. Guo, Fang, and Lu~\cite{GFL} give a polynomial-time spectral
construction, with arithmetic bound
$O((mn^9+n^{10})\log(2+m+n))$.
We accelerate that construction and give finite, verifiable
specifications for its operations.

Two costs govern the implementation. First, a favorable local direction
must permit a finite move whose length can be charged to progress. Once
the number of moves is nearly linear, the remaining task is to avoid
rebuilding dense spectral and constraint factorizations at every move. We resolve the
finite-step problem, give a complete deterministic implementation, and
then develop a randomized implementation whose remaining cost is an
explicit family of positive-definite systems.

The finite-step argument combines positivity with quadratic energy.
A comparison with the positive Perron vector controls the spectral
potential throughout a short segment. Spreading directions across the
coordinates keeps relative row derivatives small. The squared norm of
the full fractional signing measures progress through cleanup and
changes in dimension. These estimates yield $\widetilde O(n)$ accepted moves.

For the faster implementation, fixed supporting planes maintain exact
row feasibility. A bounded polynomial filter enforces the necessary
spectral leakage without constructing a high-eigenspace basis. Random
directions are then generated either by a certified linear solve or by
a fixed polynomial with constant relative accuracy. All data defining
the sampling operator are frozen before a random sample is drawn, so
applying it requires no further eigensolver calls. The total work outside
the sampler is cubic for square inputs.

The implementation notes~\cite{CertSteps,CertAnalysis} suggest a
complementary acceleration: test large candidate moves directly at their
endpoints, retaining the prescribed schedule as a bounded fallback.
We prove that tracked medium-row sums of magnitude at most $25$ certify
$L<0.84$ without spectral work. If boundedly many cheap trials produce a
coordinate-freezing endpoint at every state, exact projector maintenance
and row scans give $\widetilde O(mn^2+n^3)$ total work. This hypothesis
is checkable along an execution and does not involve inverse maintenance.

The general cubic bound is not established here. We give a complete
condition-number-sensitive algorithm, prove cubic bounds for stated
structural classes, and isolate one unproved amortization statement for a
fixed secant implementation. The distinction is summarized in
Table~\ref{i:table}. The quadratic implementations share the bound
$D_*$. Sections~\ref{opt:section} and~\ref{fin:section} give the separate
below-67 and below-37.54 constructions. The latter has the strongest
proved constant; the screened below-67 version retains the strongest
unconditional imported fast-arithmetic bound for square inputs.
Section~\ref{fin:practical} compares practical implementations under a
common output check. No one of these statements substitutes for the others.

The final construction gives
\[
 \|A\varepsilon\|_\infty<\frac{1877}{50}=37.54,
 \qquad O(mn)+\widetilde O(n^4)\text{ ordinary arithmetic work}.
\]
Theorem~\ref{fin:guarantee} specifies its profile, normalized response,
row screening, and terminal reserve. We first develop the quadratic
core, retaining its labels and independent randomized analysis, then
state precisely which parts transfer to the stronger deterministic variants.

\subsection{The complete signing guarantee}
For $n\ge1$, write
\[
 D_*=105+3\sqrt{99}=134.849623\ldots,\qquad
 J=1+\lceil\log_2n\rceil,\qquad
 q=\lceil\log_2(8(n+2m))\rceil.
\]
Empty instances are immediate. The notation $\widetilde O$ suppresses
fixed powers of $\log(2+m+n)$, except where a displayed statement
explicitly includes further accuracy or spectral logarithms. We use
$\omega>2$ for an achieved matrix-multiplication exponent. Bounds stated
with the infimum exponent instead require arbitrary fixed positive
exponent slack.

\begin{theorem}[Deterministic finite signing]\label{b:thm:main}
For every $A\in\mathbb R^{m\times n}$ with
$\sum_i a_{ij}^2\le1$, a deterministic algorithm returns
$\varepsilon\in\{-1,1\}^n$ such that
\[
 \|A\varepsilon\|_\infty<D_*.
\]
It uses $O(nJ^2q^2)$ updates. Every full update has Euclidean length at
least $c/(Jq)$ for an absolute $c>0$; every shorter update freezes a
coordinate. Its arithmetic costs are
\begin{enumerate}
\item $\widetilde O(mn^3+n^4)$ with field operations and comparisons.
Rational input produces rational states.
\item $\widetilde O(mn^\omega+n^{\omega+1})$ in the real-RAM model
with exact square roots, importing the algorithms specified in
Section~\ref{b:sec:implementation}.
\end{enumerate}
No polynomial bit-complexity bound is asserted.
\end{theorem}

The exact constant follows from retained-row tracking and terminal
rounding. The achieved exponent in~\cite{MM} gives the square-input bound
$\widetilde O(n^{3.371177})$ in the second arithmetic model.

The earlier optimized-profile alternative of Theorem~\ref{opt:guarantee}
gives $C_{67}=66.956814201324\ldots$ with the same two unscreened
arithmetic costs. Its screened extension has the sharper bound
\eqref{fin:fast67}; the final construction lowers the constant further
to $C_{37}=37.54$ with the ordinary-arithmetic cost stated above. It changes the row energy, deletion threshold,
finite-step parameters, and terminal rounding, rather than only
substituting a sharper number in Theorem~\ref{b:thm:main}.

\subsection{The modular algorithm and its cost}
In a retained epoch of dimension $s$, let $P_t$ project onto the exact
row-feasible subspace $H_t$. The sampler uses the finite arithmetic operator
\begin{equation}\label{i:operator}
 A_t=P_t(F_t+\rho_t E_t^TE_t)P_t\big|_{H_t},
 \qquad E_t\simeq f_t(M_t/\Lambda_t)T_t,
 \quad T_th=M_t'[h]v_t.
\end{equation}
Here $F_t$ is a clipped diagonal energy plus a ridge, $v_t$ is the Perron
vector used in the analysis, $f_t$ is an explicit bounded polynomial, and
$\Lambda_t$ is a certified upper spectral endpoint. The symbol $\simeq$
denotes the certified finite approximation of Section~\ref{r:filter},
not an unquantified approximation. The penalty $\rho_t$ is distinct from
the fixed row-scale ratio $\theta=101/100$.

\begin{theorem}[Complete modular signing algorithm]\label{i:modular}
The randomized algorithm in Section~\ref{a:algorithm} returns a signing
of discrepancy less than $D_*$ almost surely, with the bound valid on
every returned execution. It makes $\widetilde O(n)$ accepted moves
pathwise, and every sampling trial has conditional acceptance probability
at least $1/(4C)$, where $C=1200$.
With the residual-certified sampler its expected arithmetic cost is
\begin{equation}\label{i:cost}
 \widetilde O(mn^2+n^3)
 +\mathbb E\sum_{\text{all trials }t}
       \operatorname{SolveCost}(A_t,b_t,\mathrm{tol}_t).
\end{equation}
The solve term includes all preparation, maintenance, retries, and inverse
applications. In a large epoch,
$\rho_t=\widetilde O(s^2)$ and $\kappa(A_t)=\widetilde O(s^3)$.
Alternatively, the fixed-polynomial sampler of
Theorem~\ref{new:polysampler} may replace the solve and residual test;
it preserves the same acceptance and direction contracts.
Every operator application and every collection of acceptance tests costs
$\widetilde O(ms+s^2)$.
\end{theorem}

The two samplers give different forms of dependence on conditioning.
Conjugate gradients with a residual stopping rule needs no advance
knowledge of the condition number. If $\kappa$ uniformly bounds all systems that may be
visited, and $0<\epsilon\le1/2$ tightens the original residual tolerances,
Theorem~\ref{ifc:thm:kappa} gives
\begin{equation}\label{i:cg}
 \widetilde O\!\left((mn^2+n^3)
 \min\{n,\sqrt\kappa\log(2\kappa/\epsilon)\}\right).
\end{equation}
For polynomial sampling, supply certified spectral enclosures
$\alpha_t I\preceq A_t\preceq\beta_t I$ and set
$\bar\kappa=\sup_t\beta_t/\alpha_t$. Proposition~\ref{new:polycost}
gives instead
\begin{equation}\label{i:poly}
 \widetilde O\!\left((mn^2+n^3)\min\{n,\sqrt{\bar\kappa}\}\right).
\end{equation}
The approximation tolerance is the absolute constant $1/(100C)$.
Unlike $\kappa$, $\bar\kappa$ measures the supplied enclosure, which can
be looser than the actual spectrum. Both statements include the exact
constraint projector. A preconditioner's construction and application
costs are charged explicitly in Section~\ref{ifc:sec:kappa}.

\begin{table}[ht]
\centering\small
\caption{Arithmetic guarantees and their discrepancy constants. The final construction has bound $C_{37}=37.54$, the two below-67 rows have bound $C_{67}$, and the quadratic rows have bound $D_*$.}
\label{i:table}
\begin{tabular}{@{}L{0.24\textwidth}L{0.35\textwidth}L{0.33\textwidth}@{}}
\toprule
Implementation & Bound & Scope and reference\\
\midrule
Final, ordinary arithmetic & $O(mn)+\widetilde O(n^4)$ & Discrepancy $<C_{37}$; Theorem~\ref{fin:guarantee}\\[3pt]
Screened below 67, fast & $O(mn)+\widetilde O(m_{66}n^2+n^{\omega+1})$ & Imported fast primitives; \eqref{fin:fast67}\\[3pt]
Deterministic, rational & $\widetilde O(mn^3+n^4)$ & All inputs; Theorem~\ref{b:thm:main}\\[3pt]
Optimized profile & $\widetilde O(mn^3+n^4)$ & Discrepancy $<C_{67}$; Theorem~\ref{opt:guarantee}\\[3pt]
Fast deterministic & $\widetilde O(mn^\omega+n^{\omega+1})$ & Imported fast primitives; Theorem~\ref{b:thm:main}\\[3pt]
Conjugate gradients & Equation~\eqref{i:cg} & All inputs, conditioning explicit; Theorem~\ref{ifc:thm:kappa}\\[3pt]
Fixed polynomial & Equation~\eqref{i:poly} & Certified spectral enclosures; Proposition~\ref{new:polycost}\\[3pt]
Certified endpoints & $\widetilde O(mn^2+n^3)$ & Conditional on cheap boundary-endpoint success; no Perron data; Theorem~\ref{ep:cubic}\\[3pt]
Support separators & $\widetilde O(mn^2+n^3)$ & Supplied support tree with $(d+1)(c+1)=\widetilde O(\sqrt n)$; Corollary~\ref{ifc:cor:profilecubic}\\[3pt]
Two-sided secants & $\widetilde O(mn^2+n^3)$ & Conditional on~\eqref{d:conjecture}; Theorem~\ref{d:conditional}\\
\bottomrule
\end{tabular}
\end{table}

\subsection{What is proved about amortization}
Each of the three structural results includes the cost of constructing
its solver representation. Fixed joint profiles with the walk's nonmonotone
coefficient law admit lazy maintenance, even with arbitrary overlaps and
changing eliminated coordinate sets. A supplied signed update stream can
be processed in batches, reconstructing the representation after each
group of low-rank updates. Finally, small separators in the input row supports allow direct
construction of the actual polynomial filter, exact row projector, and
Gram interfaces. This last case gives the full signing bound in
Table~\ref{i:table}. The hypotheses and complete accounting appear in
Section~\ref{new:structured}.

For unrestricted inputs, Sections~\ref{c:section} and~\ref{d:section}
give a fixed two-sided secant controller. A verified correction decreases
a matrix divergence by a constant. The total costs of row changes, rebuilds, and numerical
errors fit the cubic budget. The remaining term is the expected signed
change in the divergence caused by changes of the filtered factor;
we call this term the expected signed production.
The unconditional accounting identity is Theorem~\ref{d:accounting};
the sole additional hypothesis for that controller's cubic bound is
\eqref{d:conjecture}. This hypothesis concerns the controller's specified probability law.
Substituting the polynomial sampler changes that law, so the hypothesis
would need to be established for the modified controller.

\subsection{Provenance, notation, and organization}
The retained-row state, concave row barriers, nonnegative Gram potential,
cancellation mechanism, and restricted-trace argument come from
Guo--Fang--Lu~\cite{GFL}. Their calculations are reproduced at the parameters
used here. The refinements developed here are the finite Perron comparison,
deterministic direction spreading, anchor-energy accounting, certified
covariance implementation, and the stated maintenance and construction
results. The constant-accuracy covariance argument refines our sampler using
classical Chebyshev approximation. Section~\ref{ep:section} integrates
the endpoint-verification proposal in~\cite{CertSteps} and its analysis
in~\cite{CertAnalysis}: the tracked-sum envelope, bounded fallback,
positive certificates, and a fully charged conditional cubic theorem.
All required arguments are given here.
Ground-state identities, Rayleigh--Temple estimates, conjugate gradients,
Chebyshev iteration, Schur complements, and hierarchical factorization are
standard tools. The faster arithmetic clauses import the algorithms of
Sobczyk~\cite{Sob}, Dumas--Pernet--Sultan~\cite{DPS}, and
Ahmadinejad--Jambulapati--Saberi--Sidford~\cite{AJSS}.
The S-divergence and quasi-Newton inverse updates are also established
tools~\cite{Sra,Gower,GR16}.

All matrices are real. Vector norms are Euclidean and matrix norms are
operator norms unless marked otherwise. The order $X\preceq Y$ means
$Y-X\succeq0$. A compressed inverse acts only on its specified subspace.
The vector $\one$ has all coordinates one, and $(t)_+=\max\{t,0\}$.
The analysis may use exact eigenvectors and exponentials; the algorithms
use the finite interfaces proved below. No bit-complexity claim is made.

Section~\ref{sec:deterministic} proves the deterministic walk.
Sections~\ref{sec:maintained}--\ref{a:algorithm} give exact constraint
maintenance, the two sampler implementations, and the full randomized
algorithm. Section~\ref{ep:section} develops certified endpoint moves,
the Perron-free cubic regime, and its existing implementation evidence.
Sections~\ref{ifc:sec:kappa} and~\ref{new:structured} establish
the complexity bounds in terms of conditioning and input structure.
Sections~\ref{c:section}--\ref{d:section} treat the conditional unrestricted
cubic bound. The appendices supply finite arithmetic, regularization,
Schur and deflation identities, and the limits of the proposed charging
arguments. The final appendix records component tests and their scope.
The complete endpoint runs in Section~\ref{ep:evidence} use a different
walk and do not test the production hypothesis~\eqref{d:conjecture}.

\section{The retained-row walk and a finite deterministic algorithm}\label{sec:deterministic}
The retained-row state provides a spectral barrier for every row.
We derive the differential estimates, construct a direction spread across
the coordinates, and prove that this direction permits a finite move. The phase budget
at the end of the section proves the geometric part of
Theorem~\ref{b:thm:main}; Appendix~\ref{app:arithmetic} supplies
the arithmetic implementations.

\subsection{State, barriers, and cleanup}\label{b:sec:state}
Start at $x=0$. Coordinates with $|x_j|=1$ are frozen; let
$V=\{j:|x_j|<1\}$ and $s=|V|$. Each row has a retained set $T_i\subseteq V$
and an offset $c_i$, initially $T_i=V$ and $c_i=0$. Write
\[
 b_{ij}=a_{ij}\one_{j\in T_i},\quad p_{ij}=b_{ij}^2,\quad
 d_i=c_i+\langle b_i,x_V\rangle,\quad q_i=\sum_jp_{ij},\quad
 G_i=\sum_jp_{ij}(1-x_j^2).
\]
Whenever a retained entry is removed, set
$c_i\leftarrow c_i+a_{ij}x_j$ and delete $j$ from $T_i$.
This preserves $d_i$. The inequalities $0\le G_i\le q_i$ and
$\sum_iq_i\le s$ always hold.

A row is large if $q_i>9/4$, medium if $0<q_i\le9/4$, and completed if
$q_i=0$. A large row imposes the exact constraint $b_i^\T h=0$ on every move.
\emph{Deliberate deletion occurs only in medium rows}: repeatedly remove an
entry with $p_{ij}>q_i/100$, recomputing $q_i$ after each removal. Eligible
pairs are taken in lexicographic order. Consequently every remaining medium
row satisfies $\max_jp_{ij}\le q_i/100$. Rows can only move to a lower class in this ordering.

Put $\theta=101/100$ and $R=3/2$. For a nonempty medium row let
$r_i^{\rm cand}=\theta^\ell\max_{j\in T_i}|a_{ij}|$, where $\ell\ge0$ is the
least integer with $(r_i^{\rm cand})^2\ge q_i$. On entry to the medium class
set $r_i=\min(R,r_i^{\rm cand})$; subsequently set
$r_i\leftarrow\min(r_i,r_i^{\rm cand})$. Thus
$q_i\le r_i^2<\theta^2q_i$, $r_i\le R$, and scales never increase.
Finding a candidate takes $O(\log(2+n))$ multiplications and comparisons,
because $q_i/\max_jp_{ij}\le n$.

For medium row-sign pairs $(i,\sigma)$, $\sigma\in\{-1,1\}$, define
\begin{equation}\label{b:eq:barriers}
 u_{i\sigma}=\frac{\sigma d_i-75}{r_i}
                +\frac72\frac{G_i}{r_i^2}+\frac{35}2,
 \qquad w_{i\sigma}=32r_i^{-4}\exp(2u_{i\sigma}/7).
\end{equation}
A pair is nearly tight if $u_{i\sigma}\ge-1/100$; it imposes
$g_{i\sigma}^\T h=0$, where $g_{i\sigma}=\nabla u_{i\sigma}$.
When taking derivatives, we hold all retained sets, offsets, and scales fixed.
Define
\begin{align}\label{b:eq:potential}
 W&=\sum_{i,\sigma}w_{i\sigma}p_ip_i^\T,&
 B&=\frac1{128}\sum_{i\text{ large}}\diag(p_i),\notag\\
 M&=W+B+\eta\one\one^\T,&
 \eta&>0,\notag\\
 L&=\lambda_{\max}(M),&
 \Phi&=(L-1)_+^3-\eta\sum_{j=1}^n x_j^2.
\end{align}
Set $L=0$ when $V$ is empty, with the matrices acting on the zero-dimensional
space. Throughout a phase, $\eta$ is fixed and satisfies $\eta s\le1/4096$.
The last sum includes frozen coordinates. The exponential and eigenvalue
expressions define the potential used in the analysis; the algorithm does
not evaluate them exactly.

Cleanup has the following order. First freeze boundary coordinates and absorb
their retained contributions into offsets. Introduce each formerly large row
that has become medium, assigning its initial medium scale. Then perform
medium-row deliberate deletions and decrease scales as specified above.
Completed rows have no matrix contribution. At initialization use the same
classification, initial scales, and deliberate-deletion procedure.

\begin{lemma}[Cleanup and initialization]\label{b:lem:cleanup}
Cleanup preserves $d_i$, cannot increase $L$ or $\Phi$, and preserves all
medium barriers $u_{i\sigma}\le0$. Initially
$L\le1/128+\eta n\le1/128+1/4096<1$ and $\Phi=0$. Every newly medium row has $d_i=0$ and
$u_{i\sigma}\le-29$.
\end{lemma}
\begin{proof}
No entry of a large row has been deliberately deleted. The exact tangency
constraint keeps
$d_i=0$, including after fixed entries are absorbed. When the row becomes medium,
$G_i\le r_i^2$ and $r_i\le3/2$, whence
$u_{i\sigma}\le-75/r_i+21\le-29$. Both sign weights coincide and
\[
 2r_i^2w_{i\sigma}\le\frac{64}{r_i^2}
       e^{(2/7)(-75/r_i+21)}
 \le\frac{256}{9}e^{-58/7}<\frac1{128}.
\]
The middle expression increases on $0<r_i\le3/2$, since its logarithmic
derivative is $(150/(7r_i)-2)/r_i>0$. The final inequality follows
from the positive Taylor sum for $e^{58/7}$ through degree $34$, which exceeds
$32768/9$. Since $pp^\T\preceq q\diag(p)$, the two new Gram terms are
bounded by the old contribution $(1/128)\diag(p)$ in positive semidefinite order.

At a fixed scale, deletion keeps $d$ fixed and decreases $G$, $p$, $u$, and
$w$. Write $t_0=7/2$, $c_0=35/2$, and $\beta=2/7$. During scale decrease,
\[
 \partial_r u=\frac{c_0-u-t_0G/r^2}{r},\qquad
 \partial_r\log w=\frac{\beta(c_0-u-t_0G/r^2)-4}{r}.
\]
Both are nonnegative whenever $u\le0$ and $G\le r^2$, because
$\beta(c_0-t_0)=4$. These conditions persist through the decrease, so existing
medium Gram terms decrease entrywise. The top eigenvalue of a nonnegative
symmetric matrix is monotone under entrywise increase. Large-row replacements
are PSD decreases, and freezing takes a principal submatrix. Cleanup leaves the full vector
$x$ unchanged, so the potential cannot increase. At initialization,
large and newly medium row contributions sum to at most $I/128$, and the ridge
has norm $\eta n\le1/4096$.
\end{proof}

\begin{lemma}[Tracking]\label{b:lem:tracking}
In each row the deliberately removed absolute coefficient mass is less than
$\frac32(10+\sqrt{99})<30$. The identity
\[
 (Ax)_i-d_i=\sum_{j\in D_i}a_{ij}(x_j-x_j^{\rm del})
\]
holds throughout, where $D_i$ records deliberate removals and $x_j^{\rm del}$
is the value at removal. In particular, the absolute tracking error is less than
$3(10+\sqrt{99})<60$.
For a completed row the same identity holds after any subsequent changes of
unfrozen coordinates within $[-1,1]$.
\end{lemma}
\begin{proof}
Deliberate removal starts at squared mass at most $R^2$. If $a^2>q/100$, then
\[
 \sqrt q-\sqrt{q-a^2}
 =\frac{|a|}{\sqrt{q/a^2}+\sqrt{q/a^2-1}}
 >\frac{|a|}{10+\sqrt{99}}.
\]
Other removals only decrease $q$. Summing the decreases in $\sqrt q$
therefore bounds the total deliberately removed mass by less than
$R(10+\sqrt{99})<30$. Fixed-coordinate removals create no tracking error;
each deliberately removed entry contributes exactly its change since removal,
of magnitude at most $2|a|$. The displayed identity follows by induction.
\end{proof}


\subsection{Curvature and the available dimensions}\label{b:sec:curvature}
Call a cleaned state admissible if $u_{i\sigma}\le0$, all large rows have
$d_i=0$, and $L\le9/8$. These invariants will be proved in
Section~\ref{b:sec:progress}. At such a state,
\[
 g_{i\sigma}=\frac{\sigma b_i}{r_i}
       -7\frac{p_i\circ x_V}{r_i^2},\qquad
 \nabla^2u_{i\sigma}=-7\frac{\diag(p_i)}{r_i^2}.
\]
The diffuseness bound $\max_jp_{ij}\le q_i/100$ implies
$\|g_{i\sigma}\|\le17/10$,
$g_{i\sigma,j}^2\le(289/100)p_{ij}/r_i^2$, and
$\|\nabla^2u_{i\sigma}\|\le7/100$.

The strictly positive matrix $M$ has a positive unit Perron vector $v$.
Its top eigenvalue is simple and has gap at least $\eta$. To see the
quantitative gap directly, $M-\eta vv^\T$ remains entrywise nonnegative,
since $M_{ij}\ge\eta$ and $v_iv_j\le1$. Its positive eigenvector $v$ has
eigenvalue $L-\eta$, while all eigenvalues on $v^\perp$ are unchanged.
The weighted maximum-norm characterization of the spectral radius of a
nonnegative matrix therefore gives $\lambda_2(M)\le L-\eta$.
Let $R_0$ be the reduced resolvent of $M$, extended by zero on $v$. Set
\begin{align}\label{b:eq:UD}
 \mathcal U&=\sum_{i,\sigma}w_{i\sigma}\beta^2(p_i^\T v)^2
                    g_{i\sigma}g_{i\sigma}^\T,\notag\\
 \mathcal D&=\sum_{i,\sigma}w_{i\sigma}\beta\frac7{r_i^2}
                    (p_i^\T v)^2\diag(p_i),\qquad\beta=2/7.
\end{align}
Then $\mathcal D$ is diagonal and nonnegative,
$\mathcal U_{jj}\le\rho_0\mathcal D_{jj}$ with
$\rho_0=289/2450$, and
\begin{equation}\label{b:eq:gram}
 \|M'[h]v\|^2\le Lh^\T\mathcal U h,\qquad
 v^\T M''[h,h]v=h^\T(\mathcal U-\mathcal D)h.
\end{equation}
For the first inequality, write $W=CC^\T$ with columns $\sqrt w\,p$.
The vector $M'[h]v$ is $C$ applied to the vector of entries
$\sqrt w\,\beta(p^\T v)g^\T h$; use $\|C\|^2=\|W\|\le L$.

\begin{lemma}[Restricted trace]\label{b:lem:trace}
If $S\succeq0$, $D\succeq0$ is diagonal, and $S_{jj}\le\rho D_{jj}$,
then every subspace of dimension greater than $\rho s$ contains a nonzero
$h$ with $h^\T(S-D)h\le0$.
\end{lemma}
\begin{proof}
Zero entries of $D$ give zero rows and columns of $S$. On the support of $D$,
$D^{-1/2}SD^{-1/2}$ has trace at most $\rho s$, so at most $\rho s$
eigenvalues can exceed one. Congruence preserves inertia. Thus $S-D$ has
at most $\rho s$ positive eigenvalues. The dimension assumption therefore gives the required vector.
\end{proof}

There are fewer than $4s/9$ large rows. A nearly tight pair contributes at least
$32\theta^{-4}e^{-1/350}$ to $\one^\T W\one\le Ls$.
Using $e^{1/350}\le350/349$, all row constraints together number at most $r_*s$,
where
\begin{equation}\label{b:eq:rowcount}
 r_*=\frac49+\frac98\frac{\theta^4}{32}\frac{350}{349}.
\end{equation}
Also $\|p_i\|^2\le q_i^2/100$ implies
\begin{equation}\label{b:eq:traceM}
 \tr M\le\left(\frac L{100}+\frac1{128}+\eta\right)s.
\end{equation}
These row and trace bounds leave enough dimensions to impose the
spectral cancellation equations below.


\subsection{A computable feasible negative subspace}\label{b:sec:subspace}
At a phase start with $S$ active coordinates, fix
\[
 \eta=\frac1{4096JS},\qquad \delta=\frac{\eta^5}{2^{60}}.
\]
The ridge parameter remains fixed until at most $S/2$ coordinates remain.
The phase schedule and its invariant are proved later. Throughout this section assume
$s\ge100$, $L\le9/8$, and all row invariants.

\subsubsection{A diagonalization interface}
For a symmetric matrix $H$ of norm at most two and a tolerance $\xi$, use a
finite routine returning a diagonal matrix $D$ and a matrix $V$ such that
there is an orthogonal analysis witness $U$ with
\begin{equation}\label{b:eq:diaginterface}
 \|V-U\|\le\xi,\qquad\|H-UDU^\T\|\le\xi.
\end{equation}
There are two implementations. Rational Jacobi rotations give
$V=U$ exactly, with cost $\widetilde O(s^3)$; the faster imported routine
has cost $\widetilde O(s^\omega)$. Both are specified in
Section~\ref{b:sec:implementation}. Neither evaluates an exact eigenvector.

Write $z_i=p_i/r_i^2$ and $\alpha_{i\sigma}=32e^{(2/7)u_{i\sigma}}$.
Approximate these nonnegative weights by $\bar\alpha_{i\sigma}\ge0$ with
sum of absolute errors at most $\delta/100$ and form $\bar M$ and its
analytic derivative map $\bar M'[h]$. Normalization gives $\|z_i\|\le1$
and uniform derivative multipliers, so the matrix and derivative errors are
at most $\delta/100$. Exponentials below the error threshold are discarded; the others are
approximated by positive clipped Taylor evaluations. The cost is
$O(m\log((m+1)/\delta))$ for all weights. Matrix entries do not require
separate exponential evaluations.

Apply \eqref{b:eq:diaginterface} to $\bar M$ with $\xi=\delta/100$ and clip
negative entries of $D$ to zero. The witness still gives
$\|M-UDU^\T\|<\delta$, and
$\tr D\le\tr M+\delta s/50$. Let $d_* =\max_jD_{jj}$ and
let $w$ be a unit normalization of the corresponding column of $V$, with
least-index tie breaking. In the rational Jacobi implementation this column
already has unit length. In the square-root implementation its normalization
uses one square root. In both implementations flip the sign, if necessary,
so that $\one^\T w\ge0$. The residual bound below makes this scalar
strictly positive, so the sign convention is unambiguous.

If $d_*\le1-\delta$, then $L\le1$. Use the common kernel of the row
constraints and $x_V^\T h=0$; no spectral cancellation is needed.
Otherwise $L>1-2\delta$. Let $V_{\bar F}$ contain the columns with
$D_{jj}\ge d_*/5$. In addition to the row constraints and $x_V^\T h=0$, impose
\begin{equation}\label{b:eq:cancel}
 V_{\bar F}^\T\bar M'[h]w=0.
\end{equation}
These are exact linear equations in the arithmetic model. Denote their
common kernel by $\mathcal H$ and its exact orthogonal projector by $P$.

\subsubsection{Counts, leakage, and the quadratic proxy}
The trace estimate~\eqref{b:eq:traceM} gives
$\dim\bar F\le(57/640)s+5/4096+3\delta s$.
Together with~\eqref{b:eq:rowcount} and
$\rho_*=(279/79)(289/2450)$, this yields
\begin{equation}\label{b:eq:dimslack}
 \dim\mathcal H-\rho_*s
 \ge\frac{19916394979897}{6225311232000000}s>\frac{s}{400}.
\end{equation}
The estimate uses only $\eta s\le1/4096$, so it remains valid for every phase.

The standard residual estimate gives $\|w-v\|\le3\delta/\eta$, with a
consistent sign. Let $P_H$ project onto true eigenvalues at least $21L/100$.
For witness columns $U_-$ outside $\bar F$, the Sylvester equation between
$M|_H$ and $D_-$ has ordered separation at least $L/100-\delta/5$.
Consequently
$\|P_HU_-\|\le\delta/(L/100-\delta/5)$.
Equation~\eqref{b:eq:cancel} additionally gives
$\|U_{\bar F}^\T\bar M'[h]w\|\le(\delta/100)\|\bar M'[h]w\|$.
Combining these bounds proves, for every $h\in\mathcal H$,
\begin{equation}\label{b:eq:leaknew}
 \|P_HM'[h]v\|\le\frac{32\delta}{\eta}\|h\|,
 \qquad |L'[h]|\le\frac{32\delta}{\eta}\|h\|.
\end{equation}
No gap inside either spectral cluster is assumed.

To construct a negative subspace, form the following quadratic proxy
for the potential's curvature:
\begin{align}\label{b:eq:proxy}
 \bar{\mathcal U}&=\sum_{i,\sigma}\bar\alpha_{i\sigma}\beta^2
              (z_i^\T w)^2g_{i\sigma}g_{i\sigma}^\T,\notag\\
 \bar{\mathcal D}&=\sum_{i,\sigma}\bar\alpha_{i\sigma}\beta7
              (z_i^\T w)^2\diag(z_i),\qquad
 \bar C=\frac{279}{79}\bar{\mathcal U}-\bar{\mathcal D}.
\end{align}
Its diagonal domination
$\bar{\mathcal U}_{jj}\le(289/2450)\bar{\mathcal D}_{jj}$ holds exactly,
regardless of whether $w$ is an eigenvector.
Lemma~\ref{b:lem:trace} and~\eqref{b:eq:dimslack} imply that $P\bar C P$ restricted
to $\mathcal H$ has at least $s/400$ nonpositive eigenvalues. Also
$\|\bar C\|<1$, as follows from $\|g\|\le17/10$,
$\max_jz_{ij}\le1/100$, and
$\sum\bar\alpha(z_i^\T w)^2=w^\T\bar Ww\le9/8+\delta$.

For the true matrices in~\eqref{b:eq:UD}, put
$C=(279/79)\mathcal U-\mathcal D$.
Cauchy--Schwarz gives
\[
 \sum\alpha\left|(z^\T v)^2-(z^\T w)^2\right|
 \le2L\|v-w\|,
 \qquad \|C-\bar C\|\le7\delta/\eta.
\]
The spectral second-derivative formula,~\eqref{b:eq:gram}, and~\eqref{b:eq:leaknew}
give
\[
 L''[h,h]\le h^\T Ch+
       \frac{2048\delta^2}{\eta^3}\|h\|^2.
\]
Thus every feasible $h$ with $h^\T\bar C h<\eta\|h\|^2$ satisfies
\begin{equation}\label{b:eq:curvaturecontract}
 \Phi''[h,h]<-1.9\eta\|h\|^2,
 \qquad |\Phi'[h]|\le\frac{3\delta}{2\eta}\|h\|.
\end{equation}
Indeed $f'(L)\le3/64$ and $f''(L)\le3/4$ for
$f(L)=(L-1)_+^3$; the error bound retains the term $f''(L)(L'[h])^2$.
The exact constraint $x_V^\T h=0$ removes the quadratic regularizer from the
first derivative. If $L\le1$, its Hessian is simply $-2\eta I$.

\subsubsection{An exactly feasible basis with controlled Gram matrix}
In the high branch form
$Q=P(\bar C-\eta I)P+(I-P)$.
In the low branch use $Q=-\eta P+(I-P)$.
In either case $\|Q\|\le2$, and more than $s/400$ eigenvalues are at most
$-\eta$. Apply~\eqref{b:eq:diaginterface} with $\xi=\eta/10000$ and retain
columns whose reported diagonal entries are at most $-3\eta/4$; call their
matrix $V_-$. There are $d>s/400$ such columns, and their span is a negative
subspace for $Q$ even after the approximation error.

Put $Z=PV_-$ and $T=Z/(1+\xi)$. Then $T$ is exactly feasible and
\begin{equation}\label{b:eq:basisgram}
 \tfrac34 I_d\preceq T^\T T\preceq I_d.
\end{equation}
To see the lower bound, let $U_-$ be the corresponding orthogonal witness
columns. Since $Q$ equals $I$ on $\ker P$,
\[
 (I-P)U_-(I-D_-)=(I-P)(Q-UDU^\T)U_-.
\]
All entries of $I-D_-$ exceed one, so
$\|(I-P)U_-\|\le\xi$. Hence $\|Z-U_-\|\le2\xi$, which gives the claimed
Gram bounds. In the high branch, projection also preserves negativity:
\[
 (Za)^\T(\bar C-\eta I)(Za)
 =(V_-a)^\T Q(V_-a)-\|(I-P)V_-a\|^2<0
\]
for every nonzero $a$. In the low branch, the restricted Hessian is
$-2\eta I$, so exact feasibility gives the same curvature conclusion
directly. Thus all vectors in the span of $T$ satisfy
\eqref{b:eq:curvaturecontract}. Exact orthonormalization is unnecessary.

\subsection{A deterministic direction spread across the coordinates}\label{b:sec:flat}
The feasible basis supplies many directions with negative curvature.
The next lemma selects a linear combination that also controls every
coordinate and row derivative, allowing a longer finite step.

\begin{lemma}\label{b:lem:deterministic-flat}
Let $T\in\mathbb R^{s\times d}$ have columns in a prescribed linear
subspace, and suppose
$\frac34 I_d\preceq T^{\mathsf T}T\preceq I_d$.
Let $g_1,\ldots,g_r\in\mathbb R^s$ satisfy $\|g_i\|\le17/10$.
Put $N=s+r$ and let $q$ be the least integer with $2^q\ge8N$.
There is a deterministic finite arithmetic procedure producing
$\varepsilon\in\{-1,1\}^d$ and a rational scalar $r_0>0$ such that
$h=r_0T\varepsilon$ belongs to the prescribed subspace and
\[
 \frac34\le\|h\|\le1,\qquad
 |h_j|\le3\sqrt{q/d},\qquad
 |g_i^{\mathsf T}h|\le\frac{51}{10}\sqrt{q/d}.
\]
After forming all products $T^{\mathsf T}g_i$, its cost is
$O(Nd\log(2Nd))$ arithmetic operations and comparisons.  No exponential,
square-root, logarithm, or probabilistic oracle is used.
\end{lemma}

\begin{proof}
Use the $N$ test vectors $a_j=e_j$ and $a_{s+i}=(10/17)g_i$.
Write $c_{kj}=a_k^{\mathsf T}T e_j$.  Then
$\sum_jc_{kj}^2\le1$ for every $k$.
Let $\lambda$ be the smallest positive integer with $\lambda^2\ge q$;
integer doubling and bisection find it, and $q\le\lambda^2\le4q$.
For independent uniform signs define
\[
 F(\varepsilon)=\sum_{k=1}^N
                 \cosh\!\left(\lambda\sum_{j=1}^dc_{kj}\varepsilon_j\right).
\]
The elementary inequality $\cosh u\le e^{u^2/2}$ gives
$\mathbb EF\le N e^{\lambda^2/2}$.

We implement conditional expectation using finite arithmetic as follows.
Set $\epsilon_0=1/(100d^2)$.
Precompute positive arithmetic relative-$\epsilon_0$ approximants to
$\cosh(\lambda c_{kj})$ and their suffix products.
At stage $j$, let $p_k=\sum_{\ell<j}c_{k\ell}\varepsilon_\ell$.
The exact two conditional expectations are
\[
 F_\pm=\sum_{k=1}^N
       \cosh\bigl(\lambda(p_k\pm c_{kj})\bigr)
       \prod_{\ell>j}\cosh(\lambda c_{k\ell}).
\]
Evaluate the remaining hyperbolic cosines to the same positive relative accuracy and
use the stored suffix products.  Each resulting sum approximates $F_\pm$
within relative error $\epsilon_1=1/(20d)$: at most $d$ relative errors
occur in a positive product, and all summands are positive.
Choose the sign with the smaller approximant, resolving ties toward $+1$.
The true chosen conditional expectation is at most
$(1+\epsilon_1)/(1-\epsilon_1)$ times the current expectation.  Since
$\epsilon_1\le1/20$,
\[
 F(\varepsilon)
 \le e^{3d\epsilon_1}\mathbb EF
 <2N e^{\lambda^2/2}.
\]
Thus for every $k$,
$\lambda|\sum_jc_{kj}\varepsilon_j|
\le\log(4N)+\lambda^2/2\le3q$, and hence
$|\sum_jc_{kj}\varepsilon_j|\le3\sqrt q$.

Find a positive dyadic rational $r_0$ with
$3/4\le r_0^2d\le1$ by binary bisection on $[0,1]$;
$O(\log(2d))$ comparisons suffice because the admissible interval has
length $(1-\sqrt3/2)/\sqrt d$. Here a dyadic rational means an integer
divided by a power of two, not necessarily a power of two itself.
The Gram bounds give $9/16\le\|r_0T\varepsilon\|^2\le1$.
Multiplication of the preceding coordinate bounds by
$r_0\le d^{-1/2}$ proves the conclusions.

For completeness, the relative approximants used above require only
arithmetic.  Every hyperbolic-cosine argument has absolute value at most
$\lambda\sqrt d\le2\sqrt{qd}$.
For $u\ge0$, repeatedly halve $u$ until it is at most one.  At that
scale, the positive Taylor series for $e^u$, truncated at
$O(\log(2Nd))$ terms, has relative error smaller than
$\epsilon_0/(10\cdot2^J)$, where $J$ is the number of halvings.
Repeated squaring, followed by
$(E+E^{-1})/2$, gives a positive relative-$\epsilon_0$ approximant to
$\cosh u$.  Here $J=O(\log(2Nd))$, so the work per scalar evaluation
is $O(\log(2Nd))$.  There are $O(Nd)$ evaluations; suffix products and
partial sums cost another $O(Nd)$.
All operations preserve rationality when $T$ and the test vectors are rational.
\end{proof}

\paragraph{Row energy and the step parameter.}
For every retained row $z_i=p_i/r_i^2\ge0$ with
$\sum_jz_{ij}\le1$, the constructed direction satisfies
\[
 \left|\tfrac27g_i^{\mathsf T}h\right|^2
       \le\frac{2601}{1225}\frac qd<3q/d,
 \qquad \sum_jz_{ij}h_j^2\le9q/d.
\]
The unit-direction bounds also give
$|\frac27g_i^{\mathsf T}h|\le17/35<1$ and
$\sum_jz_{ij}h_j^2\le1/100$.
Choose the smallest nonnegative integer $j$ such that the dyadic
$K=2^{-j}$ obeys $K^2\ge16q/d$ but $(K/2)^2<16q/d$;
if $16q/d\ge1$, simply use $K=1$.
Equivalently, halve from one while the next square remains at least
$16q/d$.  Then $\eta\le K\le1$, since $d\le s$, $q\ge1$, and $\eta\le1/(4096s)$. Moreover
\[
 |\tfrac27g_i^{\mathsf T}h|\le K,
 \qquad 0\le\sum_jz_{ij}h_j^2\le K^2,
 \qquad K^2\le\min\{1,64q/d\}.
\]
If $d>s/400$, this implies $K^2\le\min\{1,25600q/s\}$.
When the baseline choice $K=1$ applies, the greedy selector can be
omitted; any sign vector and the same normalization satisfy the
required derivative bounds with $K=1$.

The lemma also applies to the exactly feasible matrix obtained by
projecting an approximate spectral subspace, provided its Gram matrix is
between $\frac34I$ and $I$.  Exact orthonormalization is unnecessary.

We use the single global value $q=\lceil\log_2(8(n+2m))\rceil$ throughout;
it dominates the local value in the lemma. The weaker bounds
$1/2\le\|h\|\le1$ and $\Phi''[h,h]\le-\eta/8$ will suffice below.
For a row weight along this direction,
\[
 w_{i\sigma}(t)=w_{i\sigma}(0)e^{a_{i\sigma}t-b_it^2},\quad
 a_{i\sigma}=\tfrac27g_{i\sigma}^\T h,\quad
 b_i=\sum_jz_{ij}h_j^2.
\]
Thus $|a_{i\sigma}|\le K$, $0\le b_i\le K^2$, with
$K^2\le\min(1,25600q/s)$. These bounds control the relative derivatives of the entire matrix curve.


\subsection{A positive eigenvector comparison over a finite step}\label{b:sec:perron}
The next lemma converts local curvature into control over a finite step.
Its hypotheses use the nonnegative Gram structure of the signing walk.

\begin{lemma}[Perron preconditioning and a Rayleigh residual]\label{b:lem:perron}
Let
\[ M(t)=B+\sum_r w_r e^{a_rt-b_rt^2}p_rp_r^\T. \]
Suppose
$B\succeq0$ is entrywise nonnegative, $w_r\ge0$, $p_r\ge0$,
$|a_r|\le K\le1$, and $0\le b_r\le K^2$.
At zero suppose $1/2\le L\le9/8$, the top eigenvalue has gap at least
$\gamma\in(0,1]$, and its unit Perron vector $v$ is strictly positive.
Set $\nu=\min_i v_i$, $k_v=1+\log(1/\nu)$, $E=M'(0)$, and suppose
\begin{equation}\label{b:eq:perronleak}
 e:=\|P_{(L/4,L]}Ev\|\le K\gamma\nu/64.
\end{equation}
Put $\alpha=v^\T Ev$ and $z=R_0(Ev-\alpha v)$.
For $|t|\le\sqrt\gamma/(64Kk_v)$, the positive vector
$\psi(t)=v+tz$ has Rayleigh quotient $r(t)$ satisfying
\begin{align}\label{b:eq:temple}
 \|z\|&\le2K,& |z_i|&\le4Kk_vv_i,\notag\\
 \lambda_2(M(t))&\le L+\alpha t-\gamma/2,&
 |r(t)-L-\alpha t|&\le13K^2t^2,\notag\\
 0\le\lambda_1(M(t))-r(t)&\le512K^4t^4/\gamma.
\end{align}
For $\Phi(t)=(\lambda_1(M(t))-1)_+^3$ minus any quadratic energy,
\begin{equation}\label{b:eq:perronremainder}
 \Phi(t)\le\Phi(0)+\Phi'(0)t+\tfrac12\Phi''(0)t^2
                   +128K^3|t|^3+4096K^4t^4/\gamma.
\end{equation}
The derivatives on the right are those of the original potential.
\end{lemma}
\begin{proof}
Entrywise positivity gives $|E_{ij}|\le KM(0)_{ij}$, so
$|(Ev)_i|\le KLv_i$ and $\|Ev\|\le KL$.
Split the forcing into eigenvalues at most $L/4$ and the remaining modes.
On the low space the reduced resolvent is
$L^{-1}\sum_{j\ge0}(M(0)/L)^j$.
For its $i$th component, positivity bounds each summand relative to $v_i$
by $(33/32)K$, while spectral decay bounds it by $K4^{-j}/\nu$.
For each summand, take the smaller of these two bounds. Summing then
gives relative response at most
$(99/32)Kk_v$. The high response has norm at most $e/\gamma$ and relative
size at most $e/(\gamma\nu)\le K/64$.
These estimates prove the first line of~\eqref{b:eq:temple}; they also give
$|\alpha|\le e\le K/64$.

Set $A=LI-M(0)$, $Z=\diag(z_i/v_i)$, and $D=I+tZ$.
Since $Az=Ev-\alpha v$ and $z\perp v$, direct expansion gives
\begin{equation}\label{b:eq:congruence}
 D((L+\alpha t)I-M(t))D=A+tG+R(t),\quad
 G=ZA+AZ-E+\alpha I,
\end{equation}
with $Gv=0$ and $\|R(t)\|\le80K^2k_v^2t^2$.
Here $\|A\|\le2$, $\|E\|\le2K$, $\|Z\|\le4Kk_v$,
$|t|\|Z\|\le1/16$, and
$\|M(t)-M(0)-tE\|\le6K^2t^2$ suffice to check the expansion.
For example, the quadratic terms are $ZAZ-Z(E-\alpha I)-(E-\alpha I)Z$;
the remaining terms are bounded using the last Taylor estimate and
$\|D\|\le17/16$.

For any symmetric $G$ with $Gv=0$, the ground-state identity is
\[
 y^\T Gy=-\tfrac12\sum_{ij}G_{ij}v_iv_j
                (y_i/v_i-y_j/v_j)^2.
\]
Off the diagonal, $|G_{ij}|\le9Kk_vM(0)_{ij}$. Thus
$|y^\T Gy|\le9Kk_v\,y^\T Ay$.
If $u\perp Dv$ and $y=D^{-1}u$, then
$\operatorname{dist}(y,\R v)\ge\|u\|/\|D\|$.
Equation~\eqref{b:eq:congruence} therefore bounds the quadratic form of
$(L+\alpha t)I-M(t)$ on $(Dv)^\perp$ from below by
\[
 \left\{\frac{55}{64}\Big(\frac{16}{17}\Big)^2
       -\frac{80}{4096}\Big(\frac{16}{15}\Big)^2\right\}
       \gamma\|u\|^2>\frac\gamma2\|u\|^2.
\]
Min--max gives the bound on $\lambda_2(M(t))$.

The trial residual centered at $L+\alpha t$ is exactly
\[
 (M(t)-(L+\alpha t)I)\psi
 =t^2(E-\alpha I)z+(M(t)-M(0)-tE)\psi.
\]
After division by $\|\psi\|$, its norm is at most $13K^2t^2$.
This proves the Rayleigh estimate and ensures
$r(t)-\lambda_2(M(t))>\gamma/3$.
The residual about $r(t)$ is an orthogonal projection of the displayed
residual. The elementary Temple inequality gives
$0\le\lambda_1-r\le\|(M-rI)\psi\|^2/(\|\psi\|^2(r-\lambda_2))$,
proving~\eqref{b:eq:temple}. This inequality also follows directly by expanding
a unit trial vector in an orthonormal eigenbasis and comparing its variance
with $(\lambda_1-r)(r-\lambda_2)$.

Finally, expanding the Rayleigh quotient gives
\[
 r(t)=L+\alpha t+\tfrac12L''(0)t^2+c_3t^3+\mathcal R(t),\qquad
 |c_3|\le32K^3,\quad |\mathcal R(t)|\le256K^4t^4.
\]
The denominator is $1+t^2\|z\|^2$ and the quadratic coefficient in the
numerator, after division, is
$\tfrac12v^\T M''v+(Ev)^\T z=L''/2$.
For the constants, differentiation of the scalar exponential factors gives
$\|M''(0)\|\le6K^2$, $\|M'''(0)\|\le14K^3$, and
$\|M^{(4)}(t)\|\le128K^4$ on the stated interval.
Together with $\|z\|\le2K$ and $|\alpha|\le K/64$ these bound the numerator
and denominator remainders as displayed.

For $f(u)=(u-1)_+^3$, $f''$ is globally $6$-Lipschitz, giving a scalar
second-order remainder at most $|r-L|^3$, including crossings of one.
Use $|r-L-\alpha t|\le13K^2t^2$, $K|t|\le1/64$, and $|\alpha|\le K/64$.
The resulting cubic coefficient is at most $128K^3$; fourth-order Rayleigh
terms cost at most $768K^4t^4$. All relevant arguments lie below two, so
replacing $r$ by $\lambda_1$ costs at most $1536K^4t^4/\gamma$.
These bounds imply~\eqref{b:eq:perronremainder}, including the chain-rule term
$f''(L)\alpha^2$ in its exact second derivative. A quadratic regularizer has
no remainder.
\end{proof}

\paragraph{Application to the finite direction.}
Perron positivity for the ridge gives
$Lv_i\ge\eta\sum_jv_j\ge\eta$, hence $\nu\ge8\eta/9$.
The high space in~\eqref{b:eq:perronleak} lies within that in~\eqref{b:eq:leaknew}.
Because $K\ge\eta$ and $\delta\le\eta^4/2^{30}$, its leakage is at most
$32\delta/\eta\le K\eta\nu/64$.
Thus Lemma~\ref{b:lem:perron} applies with $\gamma=\eta$.
The vector $z$ and the trial vector $\psi$ appear only in the proof; the
algorithm never computes them or a reduced resolvent.

If the low branch certified $L\le1$, then the first two spectral derivatives
vanish and
$f(\lambda_1(M(t)))\le\|M(t)-M(0)\|^3\le64K^3|t|^3$
on the small interval used below. This branch needs no Perron comparison.


\subsection{Phase-dependent regularization and the full algorithm}\label{b:sec:progress}
Set $c_0=2^{-20}$. At phase start let $S$ be the active size, put
$\eta=1/(4096JS)$, $\delta=\eta^5/2^{60}$, and
$H=\lceil\log_2(2/\eta)\rceil$.
Choose the direction $h$ and its dyadic parameter $K$ as above, and choose
the largest dyadic $\tau\le1$ satisfying
\begin{equation}\label{b:eq:stepschedule}
 \tau\le c_0,\qquad K^2\tau\le c_0\eta,
 \qquad K^2H^2\tau^2\le c_0^2\eta.
\end{equation}
These tests require no square roots or logarithm primitive: find $H$ by
integer doublings, and find $\tau$ by halving. Throughout the phase,
$H\ge1+\log(1/\nu)$.

The precision choice can be made before the step size is known. Since
$K\le1$, the largest admissible
$\tau$ is at least $c_0\eta/(2H)$. Also $\eta H<1$. Consequently
\begin{equation}\label{b:eq:precisionschedule}
 \delta=\eta^5/2^{60}\le\eta^4/2^{30},\qquad
 \delta\le\eta^2\tau/4096.
\end{equation}
The first inequality ensures the Perron leakage contract, and the second
ensures $|\Phi'[h]|\le\eta\tau/128$ in~\eqref{b:eq:curvaturecontract}.

\begin{proposition}[A full or boundary update]\label{b:prop:onestep}
For either sign and any cube-contained move of length $0<t\le\tau$,
\[
 \Phi(x\pm th)-\Phi(x)
 \le(\eta\tau/128)t-\eta t^2/16
                     +128K^3t^3+4096K^4t^4/\eta.
\]
A full step decreases $\Phi$ by at least $\eta\tau^2/32$.
A shorter step increases it by at most $\eta\tau^2/128$.
All row barriers remain feasible.
\end{proposition}
\begin{proof}
The second derivative is at most $-\eta/8$, and the first derivative has the
stated bound. Equations~\eqref{b:eq:stepschedule} and
\eqref{b:eq:precisionschedule} make Lemma~\ref{b:lem:perron} applicable.
Its two remainder terms divided by $\eta t^2$ are at most
$128c_0K+4096c_0^2<1/128$. The same estimate covers the simpler low branch.
On a full step the total coefficient is therefore less than
$1/128-1/16+1/128=-3/64<-1/32$.
For shorter steps the quadratic and remainder terms are nonpositive.
Nearly tight row barriers are concave quadratics whose directional
derivatives vanish, so they cannot increase in either direction. For any other pair the positive
change is at most $\beta^{-1}K\tau\le(7/2)c_0<1/100$.
Large-row sums are preserved by their exact constraints. Cleanup can only
improve the spectral and row invariants.
\end{proof}

The complete deterministic procedure is:
\begin{enumerate}
\item If $n\le99$, return all positive signs. Otherwise set $x=0$ and begin
a phase with $S=n$ and $\eta=1/(4096Jn)$. Initialize the retained-row
state and perform cleanup.
\item Compute the finite feasible basis $T$ from Section~\ref{b:sec:subspace},
the deterministic direction $h$ and scale $K$ from Section~\ref{b:sec:flat},
and $\tau$ from~\eqref{b:eq:stepschedule}.
\item Put $t_* =\min_{h_j\ne0}(1-|x_j|)/|h_j|$ and $t=\min(\tau,t_*)$.
If $t=\tau$, use sign $+1$. Otherwise choose the least-index coordinate
attaining $t_*$ and select the sign $\xi$ so that $x_j\xi h_j\ge0$;
use $+1$ when $x_j=0$.
Move $x_V\leftarrow x_V+\xi t h$ and perform cleanup.
\item Stop when $s\le99$. Otherwise, if $s\le S/2$, start a new phase with
$S=s$ and increase $\eta$ to $1/(4096JS)$. Return to step 2.
\item At termination round remaining coordinates to their nearest signs,
with zero rounded to $+1$.
\end{enumerate}
All algebraic and greedy tie choices use least index or the positive sign.
No acceptance test computes the exact spectral potential.

\begin{proposition}[The accumulated spectral budget]\label{b:prop:budget}
All states, including phase resets, satisfy $L<9/8$.
\end{proposition}
\begin{proof}
Write $F=(L-1)_+^3$. Within a phase only $S$ coordinates can change, so
$\|x_{\rm end}\|^2-\|x_{\rm start}\|^2\le S$.
There are at most $S$ short steps. By Proposition~\ref{b:prop:onestep}, every
phase prefix increases $F$ by at most
$\eta S+S\eta/128=129/(128\cdot4096J)$.

At a boundary with new active size $S'$, cleanup has already been performed
with the old parameter. Increasing the ridge changes $L$ by at most
$(\eta'-\eta)S'\le1/(4096J)$.
If the preceding state has $L<9/8$, the tentative new value is below $5/4$,
where $F'(L)\le3/16$. Thus a reset increases $F$ by at most
$3/(16\cdot4096J)$. Initially $F=0$.
There are at most $J$ phases and fewer than $J$ resets, so every prefix obeys
\[
 F\le\frac{153}{128\cdot4096}<\frac1{512}=(1/8)^3.
\]
This proves the invariant inductively. No accumulated potential is discarded
at a reset; every ridge increase is charged explicitly.
\end{proof}

\begin{proposition}[Step length and update count]\label{b:prop:count}
Every full step has length at least $c/(Jq)$, and at most $O(nJ^2q^2)$
updates occur.
\end{proposition}
\begin{proof}
During a phase, $S/2<s\le S$ and $K^2\le25600q/s$.
Put $B_*=2\cdot4096\cdot25600Jq$. Then
$\eta/K^2\ge1/B_*$ and
$\sqrt\eta/(KH)\ge1/(\sqrt{B_*}H)$.
For $n\ge100$, $q\ge J$ and $H\le4J$, so $B_*\ge H^2$.
The largest dyadic choice in~\eqref{b:eq:stepschedule} is therefore at least
$c_0/(2B_*)=\Omega((Jq)^{-1})$.
Since $\|h\|\ge3/4$, this is also a lower bound up to a constant on the
Euclidean displacement.

The exact equation $x_V^\T h=0$ yields
$\|x+\xi th\|^2-\|x\|^2=t^2\|h\|^2$.
Consequently every full step increases squared norm by at least
$c/(J^2q^2)$, and the total available increase is at most $n$.
Every short step freezes a coordinate, so at most $n$ short steps occur.
This proves termination and the update bound. Within a single phase there are
$O(SJ^2q^2)$ updates, which is useful for arithmetic accounting.
\end{proof}

\begin{proposition}[The discrepancy constant]\label{b:prop:terminal}
The final signing satisfies $\|A\eps\|_\infty<D_*=105+3\sqrt{99}$.
\end{proposition}
\begin{proof}
At $s\le99$, no medium row survives cleanup: otherwise diffuseness would
give $q_i\le s\max_jp_{ij}\le(s/100)q_i<q_i$.
A completed row has fixed tracked sum $|d_i|<75$, inherited from its last
medium barrier or from the value zero in the large class. Its tracking error
from Lemma~\ref{b:lem:tracking} remains below $3(10+\sqrt{99})$ even after
final rounding, so its final discrepancy is below $D_*$.
A row that remains large has had no deliberate deletions and has actual
sum zero. Nearest-sign rounding changes at most $99$ coordinates by at most
one each. Since $|a_{ij}|\le1$, its discrepancy is at most $99$.
\end{proof}


\section{Maintained constraints and randomized boundary steps}\label{sec:maintained}
The randomized implementation retains feasible supporting planes
between cleanup events. Quadratic energy bounds both the cost of changing these planes and the
number of accepted moves.

\subsection{Fixed row planes and exact projector maintenance}\label{m:sec:rows}\label{m:maintenance}
We maintain the orthogonal projector onto the row-feasible space used in
Theorem~\ref{r:sampler}. This space is the common kernel of the large-row
constraints, normals to fixed supporting planes of medium-row barriers,
frozen-coordinate constraints, and two moving energy rows.
The regularized operator in Section~\ref{r:filter} controls the spectral
response without adding exact row constraints.

\subsubsection{Cleanup epochs inside regularization phases}\label{m:subsec:epochs}
Let $n$ be the original column count and set
$J=1+\lceil\log_2 n\rceil$. An outer phase starts with $N$ unfrozen
coordinates and uses
\[
                         \eta=\frac{1}{4096JN}.
\]
The phase ends when the number of unfrozen coordinates is at most $N/2$.
Within it, a cleanup epoch starts with $S$ unfrozen coordinates.
Keep the retained sets, scales, offsets, and the full $S$-dimensional
matrix index set fixed until the epoch ends. A coordinate that reaches
$\pm1$ stays in those arrays, and every subsequent direction satisfies
$h_j=0$ on that coordinate. Thus, before the final move of the phase,
\begin{equation}
 N/2<S\le N,\qquad \eta S\le1/(4096J).
                                                               \label{m:eq:epoch-size}
\end{equation}
For $S\ge10000$, end the epoch as soon as at least $S/100$ coordinates
have newly frozen, or when the outer phase ends. A move may freeze many
coordinates; process cleanup before choosing another direction.
Every direction-selection state therefore has fewer than $S/100$
held frozen coordinates. An epoch boundary alone does not change $\eta$.
At an outer boundary perform cleanup with the old $\eta$, and only then
start the new phase and increase the ridge.

This delayed cleanup preserves the state identities in
Section~\ref{b:sec:state}. In particular, with the retained quantities
$q_i=\sum_jp_{ij}$ and $G_i=\sum_jp_{ij}(1-x_j^2)$,
\[
 \sum_iq_i\le S,\qquad 0\le G_i\le q_i,
 \qquad \max_jp_{ij}\le q_i/100
\]
for each medium row. A held frozen coordinate contributes zero to $G_i$.
A row that is classified as large at epoch start keeps its exact equation
$b_i^{\mathsf T}h=0$ throughout the epoch. At its end, absorb frozen
coefficients into offsets, restrict the matrices to the surviving
coordinates, process class transitions and deliberate deletions, and
update scales in the order of Lemma~\ref{b:lem:cleanup}.
That lemma shows that cleanup cannot increase the spectral potential
or violate a row barrier, even when several frozen coordinates are
removed together. The tracking identity in Lemma~\ref{b:lem:tracking}
continues to hold.

Once fewer than $10000$ coordinates remain, clean up, discard the anchor
planes defined below, and continue the finite direction construction of
Sections~\ref{b:sec:subspace}--\ref{b:sec:progress}, with cleanup after each
move, from this current feasible state. In particular $x$ and all tracked
offsets are retained. This fixed-dimension continuation has only
polylogarithmically many remaining moves and costs
$\widetilde O(m+1)$, with an absolute constant depending on the threshold.
The terminal rounding and discrepancy proof in
Proposition~\ref{b:prop:terminal} apply.

\subsubsection{Anchors protected by quadratic energy}\label{m:subsec:anchors}
For a medium row-sign pair use the abbreviations
\[
 z_i=p_i/r_i^2,\qquad u_i=u_{i\sigma},\qquad
 g_i=\nabla u_i,\qquad a=1/100,\qquad b=1/350.
\]
Here and below an index $i$ in a sum over medium pairs includes the sign.
The exact Hessian is $\nabla^2u_i=-7\operatorname{diag}(z_i)$, and
$\sum_jz_{ij}\le1$.
Activate an inactive pair whenever $u_i(x)\ge-a$: store its anchor
$y_i=x$ and impose the fixed row
$g_i(y_i)^{\mathsf T}h=0$. While the pair is active, define
\[
 E_i(x)=\sum_jz_{ij}(x_j-y_{ij})^2.
\]
Every increment while the pair is active is orthogonal to its anchor
gradient. The quadratic barrier therefore satisfies the exact identity
\begin{equation}
 g_i(y_i)^{\mathsf T}(x-y_i)=0,
 \qquad u_i(x)=u_i(y_i)-\frac72E_i(x).                      \label{m:eq:anchor}
\end{equation}
Release an active pair after a move if $E_i>b$. Process releases before
new activations; least-index choices resolve ties. At each direction
selection, active pairs satisfy $E_i\le b$ and
\begin{equation}
                         -1/50\le u_i(x)\le0.             \label{m:eq:active-barrier}
\end{equation}
The upper bound holds throughout every active lifetime. On release,
$u_i(x)<-a$, since $u_i(y_i)\le0$ and $(7/2)b=a$.
An inactive pair starts a move below $-a$, and its positive increase is
at most $(7/2)K\tau\le(7/2)2^{-20}<a$ under the step contracts stated in
Section~\ref{m:sec:contracts}. Hence every medium barrier remains
feasible throughout the move. Overshoot of the release threshold
only decreases the protected barrier and is harmless.

In addition to $x^{\mathsf T}h=0$, impose the single aggregate equation
\begin{equation}
 r(x)^{\mathsf T}h=0,\qquad
 r(x)=\sum_{i\ {\rm active}}\operatorname{diag}(z_i)(x-y_i).
                                                               \label{m:eq:aggregate}
\end{equation}
Current row gradients $g_i(x)$ still enter the flatness tests of
Theorem~\ref{r:sampler}, even though an active pair's feasibility equation
uses its fixed anchor gradient.

\begin{lemma}[Dimension after anchoring]\label{m:lem:dimension}
Let $H_{\rm row}$ be the common kernel of the large-row equations, anchor
rows, held frozen-coordinate rows, $x^{\mathsf T}$, and
$r(x)^{\mathsf T}$. Put $\ell_0=\dim H_{\rm row}$ and
\[
 \rho=\frac{279}{79}\frac{289}{2450},\qquad
 \rho_{\rm eff}=\rho(1+3/4096).
\]
At a direction-selection state with $S\ge10000$ and $L\le9/8$, let $k$
be the dimension of the spectral subspace with eigenvalues greater than
$\Lambda/5$, where $\Lambda\ge L$ and $\Lambda>1$. Then
\begin{equation}
              \ell_0-k-\rho S>S/400,
 \qquad       \ell_0-k-\rho_{\rm eff}S>S/600.               \label{m:eq:dimension}
\end{equation}
Thus the effective-rank allowance used by Lemma~\ref{r:trace} is valid
for this row-feasible space, without imposing the high spectral equations
as exact rows.
\end{lemma}
\begin{proof}
Write $\vartheta=101/100$. Each active pair contributes at least
$32\vartheta^{-4}e^{-1/175}$ to the quantity
$\mathbf1^{\mathsf T}W\mathbf1\le LS$, since
$q_i\le r_i^2<\vartheta^2q_i$ and \eqref{m:eq:active-barrier} holds.
Using $e^{1/175}\le175/174$, their number is at most
\begin{equation}
 c_{\rm anc}S,\qquad
 c_{\rm anc}=\frac98\frac{\vartheta^4}{32}\frac{175}{174}
                         <\frac1{25}.                    \label{m:eq:anchor-count}
\end{equation}
There are fewer than $4S/9$ large rows and fewer than $S/100$ held frozen
rows. The two energy equations cost at most two dimensions.
Equation~\eqref{b:eq:traceM}, applied on the full epoch index set, gives
\[
 k\le\frac{5\operatorname{tr}M}{\Lambda}
 \le\left(\frac1{20}+\frac5{128}+5\eta\right)S
 \le\frac{57S}{640}+\frac5{4096}.
\]
Here $L/\Lambda\le1$, $\Lambda>1$, and $\eta S\le1/4096$
justify the bound, including when the true $L$ is just below one.
Consequently the first surplus divided by $S$ is at least
\[
 1-\frac49-c_{\rm anc}-\frac{57}{640}-\rho-\frac1{100}
   -\frac2{10000}-\frac5{4096\cdot10000}>0.0029>\frac1{400}.
\]
Finally $\rho<1/2$, so clipping uses less than $3S/8192$ of this
surplus. Since $1/400-3/8192>1/600$, the second assertion follows.
\end{proof}

\subsubsection{The exact row projector, including dependent rows}
\label{m:subsec:projector}
First maintain the rows that remain fixed for their entire lifetime:
large-row normals, anchor normals, and held frozen-coordinate equations. There are $O(S)$ current rows.
Store an independent row basis $B$, the inverse
$K_B=(BB^{\mathsf T})^{-1}$, the orthogonal projector
$P_0=I-B^{\mathsf T}K_BB$, and the coordinates $c_i$ of every current row
relative to $B$. These arrays have $O(S^2)$ entries.

For insertion of a row $a$, compute
$c=aB^{\mathsf T}K_B$ and $u=a-cB$. If $u=0$, record the dependent row.
Otherwise append $a$ to $B$, append a zero coordinate to all previous
records, give the new row its unit coordinate, and use
\begin{equation}
 K_{B,+}=\begin{pmatrix}
 K_B+c^{\mathsf T}c/(uu^{\mathsf T})&-c^{\mathsf T}/(uu^{\mathsf T})\\
 -c/(uu^{\mathsf T})&1/(uu^{\mathsf T})
 \end{pmatrix},\qquad
 P_{0,+}=P_0-\frac{u^{\mathsf T}u}{uu^{\mathsf T}}.           \label{m:eq:insert}
\end{equation}
Deletion of a nonbasis row only deletes its record. If basis row $j$ is
deleted and another current row has $c_{ij}\ne0$, promote the least such
row. The basis change is
\begin{equation}
 B_+=TB,\quad T=I+e_j(c_i-e_j^{\mathsf T}),\quad
 T^{-1}=I-\frac{e_j(c_i-e_j^{\mathsf T})}{c_{ij}}.
                                                               \label{m:eq:exchange}
\end{equation}
Update $K_B$ to $T^{-\mathsf T}K_BT^{-1}$ and each coordinate record to
$c_iT^{-1}$. The projector does not change. If no replacement exists, put
$v=B^{\mathsf T}K_Be_j$; removing the basis row gives
\begin{equation}
 P_{0,+}=P_0+\frac{vv^{\mathsf T}}{(K_B)_{jj}},\qquad
 (K_B)_+=(K_B)_{-j,-j}
  -\frac{(K_B)_{-j,j}(K_B)_{j,-j}}{(K_B)_{jj}}.             \label{m:eq:delete}
\end{equation}
Delete coordinate $j$ from the records. Each operation uses $O(S^2)$
arithmetic operations and comparisons. At rank zero, $P_0=I$ and the
basis arrays are empty. Exact zero comparisons handle arbitrarily small
nonzero pivots; there is no singular-value assumption in this arithmetic
cost.

At every direction selection add $x^{\mathsf T}$ and
$r(x)^{\mathsf T}$ temporarily. Given a projector $P$ and a new row $a$,
put $u=aP$; leave $P$ unchanged if $u=0$, and otherwise subtract
$u^{\mathsf T}u/(uu^{\mathsf T})$. Two such operations produce the exact
projector $P_{\rm row}$ onto $H_{\rm row}$ in $O(S^2)$ work. The stored
$P_0$ is retained for the next selection.

\begin{theorem}[Total row-maintenance cost]\label{m:thm:rowcost}
Suppose every selected direction satisfies the contracts in
Section~\ref{m:sec:contracts}, and put
$\mathcal B=\sup_t S_tK_t^2$. Over $T$ accepted moves, maintaining the exact row projector, evaluating
anchor energies, and making a bounded number of projector applications
per move costs
\[
 O\bigl(n^3(1+\mathcal B)+n^2T\bigr)
\]
arithmetic operations, apart from reading and scanning the input rows.
Each further projector application costs $O(S^2)$ and is charged to the
routine requesting it. With $\mathcal B=O(Jq)$ and $T=O(nJ^4q^2)$ this is
$\widetilde O(n^3)$. Row scans and cleanup can be included in
$\widetilde O(mn^2+n^3)$.
\end{theorem}
\begin{proof}
Let $\mathcal E=\sum_{i\ {\rm active}}E_i$. On a move $x\mapsto x+dh$,
\eqref{m:eq:aggregate} gives the exact increment
\begin{equation}
 \Delta\mathcal E=d^2\sum_{i\ {\rm active}}\sum_jz_{ij}h_j^2
       \le c_{\rm anc}SK^2d^2\le c_{\rm anc}\mathcal B d^2.
                                                               \label{m:eq:anchor-injection}
\end{equation}
Activation adds zero energy; release removes more than $1/350$;
epoch resets discard nonnegative energy. Since $x^{\mathsf T}h=0$ and
$\|h\|\ge1/2$, the full $n$-coordinate squared norm increases by at least
$d^2/4$. Hence $\sum d^2\le4n$, and ordinary releases number less than
$1400c_{\rm anc}\mathcal B n=O(n\mathcal B)$.

An ordinary epoch consumes at least $S/100$ newly frozen coordinates.
Their starting sizes therefore sum to at most $100n$. At most one epoch
per phase ends earlier because of the phase boundary; its $S$ is at most
the phase's $N$, and the $N$'s decrease geometrically. Including the final
transition to fixed dimension gives
\begin{equation}
                  \sum_{\rm large\ epochs} S^p=O(n^p)
                         \quad(p=1,2,3).                  \label{m:eq:epoch-sums}
\end{equation}
Only $O(S)$ anchors are present at a reset, so resets discard $O(n)$
anchors altogether. Every activation ends in an ordinary release, an epoch discard, or a
terminal discard. Thus activations and
releases together number $O(n(1+\mathcal B))$. Frozen-coordinate events
number at most $n$.

The explicit updates above charge $O(S^2)$ per event. Rebuilding row
bases at epoch boundaries costs $O(S^3)$ per epoch. Aggregate rows and
all anchor energies cost $O(S^2)$ per move. These observations and
\eqref{m:eq:epoch-sums} prove the displayed bound. Cleanup and row formation
can be implemented by scanning retained entries and sorting or maintaining
the largest entries for deliberate deletion, at
$O(mS\log(2+n))$ work per epoch. Scanning current row barriers and the
sampler's gradient tests costs $O(mS)$ per move. The stated soft bound
includes these costs and the fixed-dimension continuation.
\end{proof}


\subsection{Mean-zero steps and quadratic-variation accounting}
\label{m:sec:contracts}
The maintenance results use the following direction contracts.
Theorem~\ref{r:sampler} and Proposition~\ref{r:tolerance} establish them, with
an absolute constant $C_f$ and an integer
$q=O(\log(2+m+n))$ chosen at least $J$:
\begin{equation}
 \begin{gathered}
 h\in H_{\rm row},\qquad \tfrac12\le\|h\|\le1,\qquad K\le1,\\
 K^2\le C_fJq/S,\qquad \max_jh_j^2\le K^2,\\
 a_i=\tfrac27g_i(x)^{\mathsf T}h,\qquad |a_i|\le K,\\
 b_i=\sum_jz_{ij}h_j^2,\qquad 0\le b_i\le K^2.
 \end{gathered}                                             \label{m:eq:contracts}
\end{equation}
The direction may depend on every previous random choice. After the direction is selected, the boundary coin below makes the scalar
displacement conditionally mean zero. Write $c_0=2^{-20}$ and let $\tau$ be the largest dyadic
number at most one satisfying \eqref{b:eq:stepschedule}, with
$H_\eta=\lceil\log_2(2/\eta)\rceil$ in place of that equation's $H$.
In particular $K\tau\le c_0$. For either sign, the finite-step estimates
of Proposition~\ref{r:tolerance} give
\begin{equation}
 \begin{array}{ll}
 \Phi(x\pm\tau h)-\Phi(x)\le-\eta\tau^2/32,
       &\text{if the full step is in the cube},\\[1mm]
 \Phi(x\pm th)-\Phi(x)\le\eta\tau^2/128,
       &0<t<\tau\text{ and the move is in the cube}.
 \end{array}                                                \label{m:eq:two-sided}
\end{equation}
The anchor argument above establishes row feasibility for these same
moves. The potential uses all original coordinates in its energy term:
$\Phi=(L-1)_+^3-\eta\|x\|^2$.

\begin{proposition}[An unbiased boundary coin]\label{m:prop:coin}
Let $t_+,t_->0$ be the two distances to the cube boundary in directions
$h$ and $-h$, each capped at $\tau$. Choose
\begin{equation}
 d=\begin{cases}
 t_+,&\text{with probability }t_-/(t_++t_-),\\
 -t_-,&\text{with probability }t_+/(t_++t_-).
 \end{cases}                                               \label{m:eq:coin}
\end{equation}
Then, conditionally on the complete history and the selected direction,
$\mathbb E d=0$ and $\mathbb E d^2=t_+t_-$. Every realized short move
freezes a new coordinate; every realized full move has the decrease in
\eqref{m:eq:two-sided}. The resulting walk maintains $L<9/8$ and has
at most $O(nJ^4q^2)$ accepted moves, deterministically.
\end{proposition}
\begin{proof}
The conditional moments follow by substitution. A realized displacement
with $|d|<\tau$ is the actual one-sided boundary distance, so at most $n$
such outcomes occur. Full outcomes satisfy the first line of
\eqref{m:eq:two-sided}, regardless of the other coin outcome's length.
A coin of the stated real probability uses random bits and exact
comparisons: reveal a uniform binary expansion until its dyadic interval
lies entirely on one side of the probability. The expected number of
revealed bits is at most two.

We bound the spectral budget along every realized trajectory. Put $\mathcal F=(L-1)_+^3$.
A phase beginning with $N$ unfrozen coordinates has squared-norm growth
at most $N$ and at most $N$ short outcomes. Equation~\eqref{m:eq:two-sided}
therefore bounds every phase prefix increase of $\mathcal F$ by
\[
 \eta N+N\eta/128=\frac{129}{128\cdot4096J}.
\]
At an outer reset, after cleanup, the ridge increases $L$ by at most
$1/(4096J)$. Its tentative value is below $5/4$, where
$\mathcal F'(L)\le3/16$, so a reset costs at most $3/(16\cdot4096J)$.
There are at most $J$ phases and fewer than $J$ resets, and initialization
has $\mathcal F=0$. Thus
\begin{equation}
 \mathcal F\le\frac{153}{128\cdot4096}<\frac1{512},
 \qquad L<9/8.                                             \label{m:eq:spectral-budget}
\end{equation}
This closes the invariant inductively; epoch cleanup cannot increase it.

To count the moves, combine \eqref{m:eq:epoch-size} and
\eqref{m:eq:contracts} to obtain
$\eta/K^2\ge1/(8192C_fJ^2q)$. Also $H_\eta\le4J$ for $n\ge100$.
After increasing an absolute constant, the dyadic step schedule gives
$\tau\ge c/(J^2q)$. Exact tangency gives
\begin{equation}
 \|x+dh\|^2-\|x\|^2=d^2\|h\|^2\ge d^2/4.               \label{m:eq:norm-energy}
\end{equation}
Each full outcome consumes at least $c/(J^4q^2)$ of a total budget $n$;
the at most $n$ short outcomes complete the count. In a phase the same
argument gives $O(NJ^4q^2)$ moves. The fixed-dimension continuation has
no larger asymptotic cost.
\end{proof}

Within a phase, \eqref{m:eq:norm-energy} gives $\sum d^2\le4N$.
Combining this with $S>N/2$ at every selection gives the pathwise bound
\begin{equation}
 \mathcal V:=\sum_tK_t^2d_t^2\le8C_fJ^2q.                 \label{m:eq:variation}
\end{equation}
The sum may be restricted to the large cleanup epochs. This bound controls the squared increments of the coordinate path,
which suffice for the refresh argument below. It does not assert a small
total variation of the path.


\section{Sampling directions by a filtered covariance}\label{sec:sampling}
A covariance matrix can enforce the curvature and spectral-leakage
requirements without explicitly forming the high spectral space.
The sampler below specifies the required linear solve and certifies
its residual.
The finite Perron primitive used to evaluate the factor is specified
and justified in Appendix~\ref{app:arithmetic}.

\subsection{Regularized directions without a spectral basis}\label{r:regularized}
The regularized construction preserves the discrepancy and finite-step bounds
while replacing the explicit high-eigenspace equations by one structured
positive-definite system. Its solve cost will remain explicit.

In this section $s$ denotes the retained coordinate dimension of the
current metadata epoch (the variable $S$ in the algorithm specification).
The letters $S,D$ below denote matrices, not coordinate counts.
Let $P$ be the exact orthogonal projector onto the row-feasible space $H_0$,
of dimension $\ell_0$. The space $H_0$ enforces the frozen-coordinate, large-row, supporting-plane,
and two energy constraints of Section~\ref{m:sec:rows}. Let
$S,D\succeq0$, with $D$ diagonal, define the current curvature proxy, satisfying
$S_{jj}\le\rho D_{jj}$ and $\tr D\le3$, where
$\rho=(279/79)(289/2450)=80631/193550<1/2$.
The exact numerical proxies use common positive approximate coefficients
and common Perron data; their accuracy is specified below. Put
\[
 \widehat D_{jj}=\min(D_{jj},4096/s),\qquad
 F=\widehat D+\eta I,\qquad \mu=\eta+4096/s,
 \qquad \chi=\mu/\eta.
\]
In the large-epoch regime $s\ge10000$, we have $\mu\le1$ and
$\tr S\le2$. Write $\rho_{\rm eff}=\rho(1+3/4096)$.

\begin{lemma}[Effective rank can replace exact cancellation]\label{r:trace}
Let $E\in\mathbb R^{s\times s}$ have a decomposition
$E=E_{\rm hi}+E_{\rm lo}$ with $\rank E_{\rm hi}\le k$ and
$\|E_{\rm lo}\|\le d$. For $\rho_{\mathrm{ridge}}>0$ define, on $H_0$,
\begin{equation}\label{r:operator}
 \mathsf{A}=P(F+\rho_{\mathrm{ridge}} E^TE)P\big|_{H_0},\qquad
 X=P \mathsf{A}^{-1}P.
\end{equation}
Then $0\preceq X\preceq F^{-1}\preceq\eta^{-1}I$,
\[
 \rho_{\mathrm{ridge}}\tr(EXE^T)\le k+\frac{\rho_{\mathrm{ridge}} s d^2}{\eta},\qquad
 \tr((F-S)X)\ge
 \ell_0-k-\rho_{\rm eff}s-\frac{\rho_{\mathrm{ridge}} s d^2}{\eta}.
\]
\end{lemma}
\begin{proof}
The variational characterization of a compressed inverse gives
$X\preceq F^{-1}$, hence
$\tr(SX)\le\sum_j S_{jj}/F_{jj}\le\rho_{\rm eff}s$.
Indeed, $D_{jj}/F_{jj}\le1$ on the light coordinates and is at most
$sD_{jj}/4096$ on the remaining coordinates; sum and use $\tr D\le3$.
Also $\tr(FX)+\rho_{\mathrm{ridge}}\tr(EXE^T)=\ell_0$.
For an orthonormal basis $U$ of $H_0$, let $G=EU$. Since
$U^TFU\succeq\eta I$,
\[
 \rho_{\mathrm{ridge}}\tr(EXE^T)
 \le\tr\bigl(\rho_{\mathrm{ridge}} G^TG(\eta I+\rho_{\mathrm{ridge}} G^TG)^{-1}\bigr).
\]
The singular values of $G$ beyond the first $k$ are at most $d$, by its
rank-$k$ approximation $E_{\rm hi}U$. Each of the first $k$ summands is
at most one and every other summand is at most $\rho_{\mathrm{ridge}} d^2/\eta$.
\end{proof}


\subsubsection{One solve, with a finite residual certificate}
\begin{theorem}[Regularized direction sampler]\label{r:sampler}
Suppose the last lower bound in Lemma~\ref{r:trace} is at least $s/C$,
where $C\ge1$ is an absolute constant. Let $b_1,\ldots,b_N$ be all coordinate vectors and the normalized
row-gradient tests $(10/17)g_{i\sigma}$, each of norm at most one. Choose an integer $q$ with $2^q\ge16CN$, and set
\[
 \zeta_0=\min\{\eta/(1000C),1/(100s)\}.
\]
Choose independent random signs $\xi,\omega\in\{-1,1\}^s$ and form
\[
 b=P(F^{1/2}\xi+\sqrt{\rho_{\mathrm{ridge}}} E^T\omega).
\]
Compute $y\in H_0$ with the certified residual
\begin{equation}\label{r:residual}
 \|b-\mathsf{A}y\|^2\le\zeta_0^2\eta s.
\end{equation}
Accept precisely when
\begin{equation}\label{r:accept}
 y^T(F-S)y\ge\frac{s}{8C},\qquad
 |b_j^Ty|^2\le\frac{4q}{\eta}\quad(1\le j\le N).
\end{equation}
Resample otherwise. On acceptance choose a positive power of two $r$ such that
$1/4\le r^2\|y\|^2\le1$ and return $h=ry$.
The acceptance probability is at least $1/(4C)$, and
\[
 h\in H_0,\qquad \frac12\le\|h\|\le1,\qquad
 h^T(S-D)h\le\eta\|h\|^2,
\]
\begin{equation}\label{r:flat}
 |b_j^Th|^2\le\frac{32C\chi q}{s},\qquad
 \|Eh\|^2\le\frac{32C\mu}{\rho_{\mathrm{ridge}}}.
\end{equation}
Any procedure that obtains \eqref{r:residual} is admissible; its full cost
is charged separately from the cost of matrix-vector products.
\end{theorem}
\begin{proof}
Let $y_*=\mathsf{A}^{-1}b$, extended on $H_0$. Its covariance is exactly $X$.
If $B=P[F^{1/2},\sqrt{\rho_{\mathrm{ridge}}} E^T]$, then $BB^T=\mathsf{A}$ on $H_0$ and
$B^T\mathsf{A}^{-1}B$ is an orthogonal projector. Therefore, pathwise,
\[
 y_*^T\mathsf{A}y_*\le2s,\qquad
 Z_*:=y_*^T(F-S)y_*\le2s,\qquad
 \mathbb E Z_*\ge s/C.
\]
Consequently $\Pr\{Z_*\ge s/(4C)\}\ge3/(8C)$.
Each tested linear form of $y_*$ is a Rademacher sum whose squared
coefficient norm is $b_j^TXb_j\le1/\eta$. Hence the union probability
that any $|b_j^Ty_*|^2>2q/\eta$ is at most
$2N e^{-q}\le2N2^{-q}\le1/(8C)$.

Let $e=y-y_*$. The residual certificate gives
$\|e\|_{\mathsf{A}}\le\zeta_0\sqrt s$ and
$\|e\|\le\zeta_0\sqrt{s/\eta}$.
Since $\|F-S\|\le3$, $\|y_*\|\le\sqrt{2s/\eta}$, and $\zeta_0\le1$,
\[
 |y^T(F-S)y-Z_*|
 \le 3(2\sqrt2\zeta_0+\zeta_0^2)s/\eta
 \le12\zeta_0s/\eta< s/(8C).
\]
Also $|b_j^Te|\le1/(100\sqrt\eta)$, so the good exact-sample events
imply every test in \eqref{r:accept}. Their intersection has probability
at least $1/(4C)$.

For every accepted output, including those outside the good events,
$y^TFy\ge s/(8C)$ implies $\|y\|^2\ge s/(8C\mu)$.
Thus $r^2\le8C\mu/s$, proving the flatness bound. Moreover
$\|y\|_{\mathsf{A}}\le(\sqrt2+\zeta_0)\sqrt s\le2\sqrt s$, so
$\rho_{\mathrm{ridge}}\|Ey\|^2\le4s$, proving the leakage bound.
Finally $F\preceq D+\eta I$ and the first acceptance test imply the
curvature inequality.
\end{proof}

Square roots need not be exact primitives. Choose $\rho_{\mathrm{ridge}}$ as a power of
four, so $\sqrt{\rho_{\mathrm{ridge}}}$ is rational. Approximate each $\sqrt{F_{jj}}$ to
absolute error at most $\zeta_0\eta/8$, obtaining $\widetilde b$ with
$\|\widetilde b-b\|\le\zeta_0\sqrt{\eta s}/8$.
It suffices to certify
$\|\widetilde b-\mathsf{A}y\|^2\le\zeta_0^2\eta s/4$;
then \eqref{r:residual} follows by the triangle inequality.
Bisection computes these roots in
$O(s\log(2+1/(\zeta_0\eta)))$ arithmetic operations because $F_{jj}\le1$.
All rejection tests are finite quadratic comparisons. The acceptance probability has a lower bound independent of the condition number.

\paragraph{The strict curvature margin pays for finite proxies.}
The residual proof gives $\|y\|^2\le4s/\eta$. Dyadic normalization therefore
also gives $r^2\ge1/(4\|y\|^2)\ge\eta/(16s)$, and every accepted output obeys
\[
 h^T(F-S)h\ge\frac{\eta}{128C}.
\]
In the spectral branch $\Lambda>1$ defined below, let
$\mathcal C=(279/79)\mathcal U-\mathcal D$ be the true curvature proxy
from~\eqref{b:eq:UD}. In that branch request the finite arithmetic accuracy
\begin{equation}\label{r:proxy-accuracy}
 \|\mathcal C-(S-D)\|\le\frac{\eta}{256C}.
\end{equation}
The proxies are constructed from common nonnegative coefficients and
common Perron data, so their exact diagonal comparison survives.
The trace bound remains at most three under sufficiently accurate
normalization. The explicit retained-row formulas achieve
\eqref{r:proxy-accuracy} with inverse-polynomial coefficient and vector
accuracy, strengthening the filter accuracy requests if necessary.
Since $F\preceq D+\eta I$ and $\|h\|\le1$,
\[
 h^T\mathcal C h
 \le\eta\|h\|^2-\frac{\eta}{128C}+\frac{\eta}{256C}
 <\eta\|h\|^2.
\]
Thus each accepted output in the spectral branch satisfies the true
curvature contract. The branch
$\Lambda\le1$ uses the direct cubic potential estimate and does not request
approximation of this true curvature proxy by its auxiliary zero matrices.


\subsection{A fixed polynomial needs only constant relative accuracy}
\label{new:polynomial}
The preceding residual criterion is sufficient for arbitrary solvers.
For a fixed linear sampling map, covariance comparison permits constant
relative accuracy in place of an inverse-polynomial solve tolerance. Polynomial sampling is an established technique
\cite{ChebSample}; the point here is to verify the original signing
contracts with an absolute constant relative tolerance.

\begin{theorem}[Constant-accuracy polynomial sampler]\label{new:polysampler}
Assume the hypotheses and notation of Theorem~\ref{r:sampler}. On $H_0$
let $p$ be a real polynomial, chosen before the fresh Rademacher vectors,
such that
\begin{equation}\label{new:polycontract}
 \|I-\mathsf A p(\mathsf A)\|\le\nu,
 \qquad 0<\nu\le\frac1{100C}.
\end{equation}
Form the same exact arithmetic right-hand side
$b=P(F^{1/2}\xi+\sqrt{\rho_{\mathrm{ridge}}}E^T\omega)$ and put
$y=p(\mathsf A)b$. Apply the original tests~\eqref{r:accept} and the
same dyadic normalization. Every accepted output satisfies all the
conclusions of Theorem~\ref{r:sampler}, including its strict curvature
margin, and each trial succeeds with probability at least $1/(4C)$.
No residual test~\eqref{r:residual} is required for this implementation.
The arithmetic model in this statement permits exact square roots.
\end{theorem}
\begin{proof}
All operators below act on $H_0$ and are extended by zero outside it.
Put $X=P\mathsf A^{-1}P$ and
$B=P[F^{1/2},\sqrt{\rho_{\mathrm{ridge}}}E^T]$, so $BB^T=\mathsf A$
on $H_0$. The map $p(\mathsf A)$ is fixed independently of the signs;
therefore the covariance $Y=\mathbb E[yy^T]$ satisfies
\begin{equation}\label{new:polycov}
 Y=p(\mathsf A)\mathsf A p(\mathsf A),\qquad
 (1-\nu)^2X\preceq Y\preceq(1+\nu)^2X.
\end{equation}
The trace proof of Lemma~\ref{r:trace} gives
$\tr(FX)\le s$, $\tr(SX)\le s$, and
$\tr((F-S)X)\ge s/C$. As $F,S\succeq0$, it follows that
\[
 \mathbb E Z:=\mathbb E[y^T(F-S)y]
 \ge\frac{s}{C}-(4\nu+\nu^2)s\ge\frac{19s}{20C}.
\]
Furthermore $\|\mathsf A^{1/2}p(\mathsf A)B\|\le1+\nu$:
factor this map through $\mathsf A^{-1/2}B$, which has norm one.
The sign vector $(\xi,\omega)$ has squared norm $2s$. Hence, pathwise,
\[
 Z\le y^TFy\le y^T\mathsf Ay
 \le2(1+\nu)^2s\le21s/10<4s.
\]
An upper bound and the expectation lower bound suffice even when $Z$
can be negative. With $a=s/(8C)$, they imply
\[
 \Pr\{Z\ge a\}\ge
 \frac{19/(20C)-1/(8C)}{21/10}=\frac{11}{28C}.
\]
For every unit-norm test vector $b_j$, the linear form $b_j^Ty$ is a
Rademacher sum with squared coefficient norm at most
$(1+\nu)^2/\eta$. Thus
\[
 \Pr\{|b_j^Ty|^2>4q/\eta\}
 \le2\exp\!\left(-\frac{2q}{(1+\nu)^2}\right)\le2e^{-q}.
\]
The union bound is at most $1/(8C)$. The probability of passing every
test is at least $11/(28C)-1/(8C)=15/(56C)>1/(4C)$.

For every accepted output, $\|y\|^2\ge s/(8C\mu)$ and
$y^T\mathsf Ay\le4s$. Consequently the proof following~\eqref{r:accept}
gives, without change,
\[
 |b_j^Th|^2\le\frac{32C\chi q}{s},\qquad
 \|Eh\|^2\le\frac{32C\mu}{\rho_{\mathrm{ridge}}},\qquad
 h^T(F-S)h\ge\frac{\eta}{128C}.
\]
The latter inequality follows from $\|y\|^2\le4s/\eta$ and the
dyadic normalization. It absorbs the same finite-proxy error
as~\eqref{r:proxy-accuracy}. Feasibility follows because every polynomial
iterate lies in $H_0$. The original finite-step proof therefore applies.
\end{proof}

The polynomial must be fixed before this trial's signs are drawn. An
adaptively truncated conjugate-gradient (CG) iterate cannot be assigned covariance
\eqref{new:polycov} without a separate argument. Exact CG termination
returns $\mathsf A^{-1}b$ and is an admissible fallback. Perron data,
coefficient approximations, filter coefficients, and the exact projector
are all frozen during a trial. Section~\ref{new:polyimplementation}
gives the explicit recurrence and its full cost.


\subsection{A bounded polynomial filter eliminates the spectral basis}\label{r:filter}
Obtain a certified interval $[\ell_L,\Lambda]$ of width at most
$\delta=\eta^5/2^{60}$ for the current top eigenvalue, using
Proposition~\ref{p:certificate}. In the branch $\Lambda\le1$, use
$S=D=E=0$ and $F=\eta I$; the system is simply $\eta I$ on $H_0$.
The direct cubic estimate in Section~\ref{b:sec:perron} applies.
In the other branch retain the original endpoint $\Lambda>1$ when refining
the Perron data. It satisfies $1-\delta<L\le\Lambda\le(1001/1000)L$,
so $L>1/2$, even if $L\le1$. No exact comparison of $L$ with one is needed.
Write $M\succeq0$ for the current matrix, $L=\|M\|$, and
$T h=M'[h]v$ for the exact derivative map at the current Perron vector. By~\eqref{b:eq:gram} it has norm at most an absolute constant: $\|T\|^2\le L\|\mathcal U\|\le
L\rho_0\tr\mathcal D<1$ under the retained-row bounds.
Using a looser bound $\|T\|\le2$ suffices below.

For a prescribed $0<\varepsilon_f<1/2$, choose an even integer
$d_f=O(\log(1/\varepsilon_f))$ sufficiently large and set
\[
 f(t)=\sum_{j\ge\lceil(409/2000)d_f\rceil}
          \binom{d_f}{j}t^j(1-t)^{d_f-j}.
\]
The degree constant is absolute, although large. The binomial tail bound
gives, after increasing that constant to cover rounding,
\[
 0\le f(t)\le1\ (0\le t\le1),\quad
 f(t)\le\varepsilon_f\ (t\le1/5),\quad
 f(t)\ge1-\varepsilon_f\ (t\ge209/1000).
\]
For example, take $d_f\ge2000$ and
$d_f\ge (1/(2(2/500)^2))\log(1/\varepsilon_f)$:
rounding leaves each relevant deviation at least $2/500$.
No logarithm primitive is required; an integer multiple of the least
$q_f$ with $2^{q_f}\ge1/\varepsilon_f$ suffices.

Put $E_*=f(M/\Lambda)T$. Split at eigenvalues $\lambda>\Lambda/5$.
The high part has rank at most $k$, where $k$ counts eigenvalues
exceeding $\Lambda/5$. By~\eqref{b:eq:traceM}, $\Lambda\ge L$, and
$\Lambda>1$,
\[
 k\le5\tr M/\Lambda
 \le(1/20+5/128+5\eta)s
 =57s/640+5\eta s.
\]
The low part has norm at most $2\varepsilon_f$.
Since $209\Lambda/1000<21L/100$, on the true high space
$H=\{\lambda(M)\ge21L/100\}$ one has
\begin{equation}\label{r:true-leakage}
 \|P_HTh\|\le2\|E_*h\|.
\end{equation}
This preserves the curvature coefficient $279/79=1+2/(1-21/100)$.


Suppose the finite arithmetic operator $E$ satisfies
$\|E-E_*\|\le\varepsilon_E$. Its effective-rank tail is then
$d=2\varepsilon_f+\varepsilon_E$.
The dimension bound of Lemma~\ref{m:lem:dimension} leaves a clipped-covariance margin of at least $s/600$.
Choose $d^2\le\eta/(1200\rho_{\mathrm{ridge}})$, so Lemma~\ref{r:trace} retains
the sampler hypothesis with $C=1200$.
For a desired true high-space leakage $e_{\rm high}>0$, choose $\rho_{\mathrm{ridge}}$
as the least power of four at least $512C\mu/e_{\rm high}^2$, and
also require $\varepsilon_E\le e_{\rm high}/4$.
Then every accepted output satisfies, by \eqref{r:flat} and
\eqref{r:true-leakage},
\[
 \|P_HM'[h]v\|\le e_{\rm high}.
\]
The final choice of $e_{\rm high}$ is the larger tolerance $e_*$
in Proposition~\ref{r:tolerance}. This also controls the top directional
derivative, since $v\in H$.

\subsubsection{Finite operator and proxy accuracy}\label{r:finite}
The sum of the absolute values of the Bernstein polynomial's monomial
coefficients is at most
$3^{d_f}$. For matrices $\mathsf{A},B$ of norm at most two,
\[
 \|f(\mathsf{A})-f(B)\|\le d_f6^{d_f}\|\mathsf{A}-B\|.
\]
Thus approximating $M$ and the current Perron vector sufficiently to
make the two resulting operator errors at most $\varepsilon_E/4$ each
requires only inverse-polynomial precision: $d_f=O(\log(1/\varepsilon_f))$
and $6^{d_f}$ is polynomial in $1/\varepsilon_f$.
The exponent is an absolute constant, not a dimension-dependent exponent.
Scalar coefficient precision and the already certified Perron accuracy
must include this factor. This is an arithmetic-operation statement,
not a practical or bit-complexity claim.
For example, absolute error at most
$\varepsilon_E/(2^{16}(m+1)d_f6^{d_f})$ in each row coefficient, and
Perron-vector error at most $\varepsilon_E/(2^{16}(m+1))$, suffice while retaining the original certified scalar endpoint
$L\le\Lambda\le1.001L$. The elementary termwise estimate uses
$\|z_i\|\le1$, $\|g_i\|\le17/10$, and $\alpha_i\le32$;
it also ensures the computed normalized matrix has norm at most two.
All approximated coefficients, the Perron vector, the scalar certificate,
and polynomial coefficients are frozen for the entire trial. They define
one fixed finite arithmetic linear operator $E$, evaluated exactly on
each matrix-vector call; $E^T$ is its actual transpose. The solve and
residual test use precisely this same operator. Independently rounded or
state-dependent matrix-vector products are not covered by the proof.

Both $T$ and $T^T$ can be applied in $O(ms)$ operations using the
retained-row formula
\[
 T=\sum_{i,\sigma}\alpha_{i\sigma}\beta(z_i^Tv)
                 z_i g_{i\sigma}^T,
 \quad\beta=2/7.
\]
An application of $f(M/\Lambda)$ uses $O(d_f)$ positive-Gram
matrix-vector products, hence $O(msd_f)$ operations. Polynomial
coefficients can be constructed in $O(d_f^2)$ scalar operations. Exact
application of $P$, already maintained by the anchored-row method,
costs $O(s^2)$. Consequently $\mathsf{A}$ can be applied in
$\widetilde O(ms+s^2)$ operations, and the random right-hand side,
residual, curvature, and flatness checks have the same soft cost.

Proposition~\ref{r:cost} accounts separately for the solve cost, which
may require many operator applications.


\subsection{Perron response and the finite-step guarantee}\label{r:response}
The positive ridge also controls the response of the Perron vector.
The following estimate permits a larger leakage tolerance while
preserving the finite-step constants.
Let $M=M_0+\eta\mathbf1\mathbf1^T$, where $M_0$ is symmetric and entrywise
nonnegative and $\eta>0$. Let $Mv=Lv$, $v>0$, $\|v\|_2=1$, and write
$S_v=\mathbf1^Tv$. Denote the reduced resolvent on $v^\perp$ by $R_0$.

\begin{lemma}\label{r:response-lemma}
The matrix $R_b=(LI-M_0)^{-1}$ exists and is entrywise nonnegative.
For every $a\perp v$,
\[
 R_0a=(I-vv^T)R_ba,
 \qquad
 \max_i\frac{|(R_0a)_i|}{v_i}
 \le\frac{\max_j a_j-\min_j a_j}{\eta S_v}
 \le\frac{2\|a\|_\infty}{\eta S_v}.
\]
In addition,
\[
 \|R_0a\|_2\le\frac{2\|a\|_\infty}{\eta S_v},
 \qquad 0\le a^TR_0a\le\frac{\|a\|_\infty^2}{\eta}.
\]
\end{lemma}
\begin{proof}
Since $M_0v=Lv-\eta S_v\mathbf1<Lv$, the nonnegative matrix $M_0$ has
spectral radius below $L$ (apply the weighted maximum norm).
The Neumann series gives $R_b\ge0$ entrywise, and
\[
 R_b\mathbf1=v/(\eta S_v),\qquad
 \mathbf1^TR_b\mathbf1=1/\eta.
\]
If $a\perp v$, symmetry implies $\mathbf1^TR_ba=0$; consequently
$(LI-M)R_ba=a$. Subtracting its component along $v$ gives the unique
solution in $v^\perp$, proving the identity.
Set $y=R_ba$. The ratio $y_i/v_i$ is $1/(\eta S_v)$ times a convex
combination of the coordinates of $a$, since the row sum of $R_b$ is
$v_i/(\eta S_v)$. Also $v^Ty=\sum_i v_i^2(y_i/v_i)$ is a convex
combination of these ratios. Their differences are therefore at most
$(\max a-\min a)/(\eta S_v)$ in absolute value. The Euclidean bound
follows from $\|v\|_2=1$.
Finally $a^TR_0a=a^TR_ba\ge0$, while entrywise nonnegativity gives
$a^TR_ba\le\|a\|_\infty^2\mathbf1^TR_b\mathbf1$.
\end{proof}

\begin{proposition}[Larger leakage with unchanged step constants]\label{r:tolerance}
Assume an admissible current matrix with $1/2\le L\le9/8$ and
$\eta\le1/4096$, and a row-feasible direction $h\in H_0$ with
$1/2\le\|h\|\le1$, $x^Th=0$, $\eta\le K\le1$, and
$h^T\mathcal C h\le\eta\|h\|^2$. The row curves satisfy $|a_i|\le K$ and
$0\le b_i\le K^2$.
Let $\tau$ be the largest power-of-two step satisfying~\eqref{b:eq:stepschedule}. It suffices to require
\[
 \|P_{[21L/100,L]}M'[h]v\|\le e_*,\qquad
 e_*:=\min\{K\eta/256,\ \eta\tau/128\}.
\]
Then the pointwise remainder in Lemma~\ref{b:lem:perron} has the same
constants $128$ and $4096$. Moreover
$|\Phi'[h]|\le\eta\tau/128$ and
$\Phi''[h,h]\le-1.9\eta\|h\|^2$.
For both signs, full moves decrease $\Phi$ by at least $\eta\tau^2/32$;
short moves increase it by at most $\eta\tau^2/128$, and all row barriers
remain feasible.
\end{proposition}
\begin{proof}
First split at $L/4$, as in Lemma~\ref{b:lem:perron}. Write
$b=P_{(L/4,L]}M'[h]v$, $\alpha=v^Tb$, and
$z_{\rm hi}=R_0(b-\alpha v)$. Lemma~\ref{r:response-lemma} gives
\[
 |z_{{\rm hi},i}|/v_i\le4e_*/(\eta\mathbf1^Tv)\le K/64,
 \qquad\|z_{\rm hi}\|\le K/64,\qquad |\alpha|\le e_*\le K/64.
\]
To avoid a factor equal to the reciprocal of the smallest Perron
coordinate in the low-response bound, use
\[
 b=\alpha v+(LI-M)z_{\rm hi}.
\]
Since $M\ge0$ entrywise and $Mv=Lv$, for every integer $j\ge0$,
\[
 \frac{|((M/L)^jb)_i|}{Lv_i}
 \le\frac{|\alpha|}{L}+\frac K{32}.
\]
The whole forcing obeys $|M'[h]v|\le KLv$. Thus each term in the
low-response Neumann series, after division by $L$, is bounded relative
to $v_i$ by
\[
 K+|\alpha|/L+K/32\le17K/16.
\]
Its spectral bound remains $K4^{-j}/\nu$, $\nu=\min_i v_i$.
The elementary geometric-tail sum is therefore at most
$(51/16)K(1+\log(1/\nu))$. Including the high response gives
$|z_i|\le4Kk_vv_i$, with $k_v=1+\log(1/\nu)$.
The low Euclidean norm is at most $4K/3$, so $\|z\|<2K$.
These are the three inputs used after the forcing split in the proof of Lemma~\ref{b:lem:perron}: $\|z\|\le2K$, $|z_i|\le4Kk_vv_i$, and
$|\alpha|\le K/64$. The response is the full exact $R_0(M'[h]v-\alpha v)$,
so the Rayleigh quadratic coefficient is still the true $L''/2$.
The ground-state congruence, residual, Temple bound, and Taylor constants
therefore remain unchanged.

For the curvature comparison split instead at $21L/100$, as required by
$279/79$. The removed Perron component has no resolvent contribution,
and the high-space response contributes at most $8e_*^2/\eta$.
Consequently, using the true proxy contract already certified above,
\[
 L''[h,h]\le \eta\|h\|^2+8e_*^2/\eta.
\]
Here $e_*\le\eta/256$, hence $8e_*^2/\eta\le\eta/8192$.
Using $\varphi(L)=(L-1)_+^3$, $\varphi'(L)\le3/64$,
$\varphi''(L)\le3/4$, $\eta\le1$, and $\|h\|^2\ge1/4$ gives
\[
 \frac{\Phi''}{\eta\|h\|^2}
 \le-2+\frac3{64}+\frac3{131072}+\frac3{65536}<-1.9.
\]
This includes the nonnegative chain-rule term $\varphi''(L)\alpha^2$.
Since the energy first derivative vanishes,
\[
 |\Phi'|\le3e_*/64\le3\eta\tau/8192<\eta\tau/128.
\]
The step caps make the remainder divided by $\eta t^2$ at most
$128\cdot2^{-20}+4096\cdot2^{-40}<1/128$. Hence the full-step and short-step
budgets follow from Proposition~\ref{b:prop:onestep}.
The exact supporting planes and the row-curve bounds preserve row
feasibility in both signs. The certified low branch uses the direct
cubic estimate of Section~\ref{b:sec:perron}.
\end{proof}


\subsection{Parameter choices and the certified-solve interface}\label{r:parameters}

Use $C=1200$. Choose once the global integer
\[
 q=Q:=\left\lceil\log_2\bigl(16C(n+2m)+2\bigr)\right\rceil
\]
by integer doublings; increasing it to dominate
$J=1+\lceil\log_2n\rceil$ is harmless. This global choice covers the coordinate and normalized-gradient tests
in every epoch.
First compute the row-feasible data and
$\mu=\eta+4096/s$, $\chi=\mu/\eta$, and
$B_{\mathrm{flat}}=32C\chi q/s$.
If $B_{\mathrm{flat}}\ge1$, take $K=1$; otherwise take the least power of two
at least $\sqrt{B_{\mathrm{flat}}}$, capped at one, using squared
comparisons. Thus
\[
 \min(1,B_{\mathrm{flat}})\le K^2
       \le\min(1,4B_{\mathrm{flat}}).
\]
The coordinate tests give
$b_i=\sum_jz_{ij}h_j^2\le K^2$ because $\sum_jz_{ij}\le1$.
The normalized-gradient tests give
$|a_i|=|(2/7)g_i^Th|\le K$, since $(2/7)/(10/17)=17/35<1$.
The unconditional unit-direction bounds cover the cap $K=1$.
In particular $\eta\le K\le1$, and none of these facts depends on the
ridge penalty.

Compute $\tau$ from~\eqref{b:eq:stepschedule}, then set
\[
 e_*=\min\{K\eta/256,\eta\tau/128\}.
\]
In the branch $\Lambda>1$, select $\rho_{\mathrm{ridge}}$ as the least
power of four at least $512C\mu/e_*^2$.
Next select positive dyadic $\varepsilon_f,\varepsilon_E$ satisfying
\[
 \varepsilon_f<1/2,\qquad
 \varepsilon_E\le\min\{1,e_*/4\},\qquad
 (2\varepsilon_f+\varepsilon_E)^2
             \le\frac{\eta}{1200\rho_{\mathrm{ridge}}}.
\]
With these parameters fixed, construct the finite polynomial operator
and numerical proxies, enforce~\eqref{r:proxy-accuracy}, and then sample.
The branch $\Lambda\le1$ instead uses $S=D=E=0$, $F=\eta I$.
Its system is $\eta I$ on the row-feasible space, so the solve is
multiplication by $\eta^{-1}$ there. These choices have no circular
dependence on the sampled direction or on the linear-system solution.

\begin{proposition}[Penalty and condition-number bounds]\label{r:conditioning}
For these parameters, throughout the large metadata epochs,
\[
 \rho_{\mathrm{ridge}}=\widetilde O(s^2),\qquad
 \kappa(\mathsf{A})\le
       \frac{\mu+9\rho_{\mathrm{ridge}}}{\eta}
       =\widetilde O(s^3).
\]
\end{proposition}
\begin{proof}
The outer halving-phase size is within a constant factor of $s$, giving
$\eta^{-1}=O(Js)$, $\mu=\Theta(1/s)$, and $\chi=O(J)$.
The a priori dyadic choice gives $K^{-2}=O(s)$ and
$K^2=O(\min\{1,Jq/s\})$. Write
$H_\eta=\lceil\log_2(2/\eta)\rceil$.
The largest-dyadic rule in~\eqref{b:eq:stepschedule} yields
\[
 \tau^{-1}
 =O\!\left(\max\{1,K^2/\eta,KH_\eta/\sqrt\eta\}\right)
 =O\!\left(\max\{1,J^2q,J\sqrt q\,H_\eta\}\right).
\]
The absolute step constant $2^{-20}$ is included in these bounds.
They also cover the cap $K=1$, because that case implies $s=O(Jq)$.
Rounding the penalty to a power of four changes it by at most four, so
\[
 \rho_{\mathrm{ridge}}
 =O\!\left(\mu\eta^{-2}\max\{K^{-2},\tau^{-2}\}\right)
 =\widetilde O(s^2).
\]
The bounded filter has $\|E_*\|\le2$, and $\varepsilon_E\le1$ gives
$\|E\|\le3$. Consequently the compressed system's eigenvalues lie
between $\eta$ and $\mu+9\rho_{\mathrm{ridge}}$.
\end{proof}

\begin{proposition}[The certified one-solve reduction]\label{r:cost}
All direction work outside the residual-certified solves costs expected
$\widetilde O(ms+s^2)$ per accepted direction. The complete modular
algorithm therefore has arithmetic cost
\begin{equation}\label{r:total}
 \widetilde O(mn^2+n^3)
 +\mathbb E\sum_{\text{sampling trials }t}
   \operatorname{SolveCost}\!\left(
       \mathsf{A}_t,\widetilde b_t,\zeta_{0,t}^2\eta_ts_t/4\right).
\end{equation}
Here $\operatorname{SolveCost}$ counts preparation, updates, and the work
needed to return a vector in the exact row-feasible space with the
indicated squared-residual certificate. This proposition leaves the solve term explicit.
\end{proposition}
\begin{proof}
The retained-row formulas apply $M,T,T^T$ in $O(ms)$ arithmetic
operations. A degree-$d_f$ polynomial and its transpose require $O(d_f)$
such products; its coefficients cost $O(d_f^2)$ scalar operations.
The exact row projector costs $O(s^2)$ per application. Consequently
every system application, random right-hand side, residual check,
curvature test, and collection of row and coordinate tests costs
$\widetilde O(ms+s^2)$. The coefficient approximations and certified
Perron computation have the same expected soft bound at the prescribed
accuracies (Proposition~\ref{p:certificate}). Theorem~\ref{r:sampler}
gives a constant expected number of trials per accepted direction.
The global sum, including cleanup, exact row maintenance, and optional
cache updates, is proved in Proposition~\ref{a:accounting}. Its proof
uses this per-trial bound and the pathwise move count; it does not
assume a bound on the solve term.
\end{proof}

The condition-number bound alone does not guarantee polylogarithmically
many iterations per solve. Equation~\eqref{r:total} is therefore a
reduction to a specified changing positive-definite system, not an
unconditional cubic total-time theorem. It does establish all direction,
finite-step, and discrepancy guarantees without forming an explicit
high-eigenspace basis.


\section{The complete modular randomized algorithm}\label{a:algorithm}

The following algorithm combines the preceding components. The sampler
uses either a residual-certified solve or the fixed polynomial of
Theorem~\ref{new:polysampler}, with its cost charged explicitly.
The positive-shift Gram cache is optional and does not affect the selected
directions or their analysis. Its maintained inverse is not used in place
of the ridge sampler's inverse.

\paragraph{Algorithm overview.}
\emph{Input:} a real matrix with column norms at most one and a selected
sampler implementation. \emph{Output:} a signing with discrepancy less
than $D_*$. The detailed parameter order and tie-breaking conventions
are specified in the subsections below.
\begin{enumerate}
\item Initialize $x=0$, the row data, and the first halving phase. Handle
the bounded-size stopping case directly.
\item Perform required cleanup and phase changes. If at most 99
coordinates remain, round to the nearest signs and return. Below the
large-epoch threshold, finish by the deterministic continuation.
\item Update the active row planes and frozen-coordinate equations, and
maintain their exact kernel projector.
\item Set the step cap and leakage allowance. Certify the current Perron
data, determine the low or high branch, and freeze the finite curvature
proxies and filter.
\item Draw a direction using the selected residual-certified or fixed-polynomial
sampler. Reject failed samples and repeat with fresh signs; the state and
operator remain fixed during these retries.
\item Compute the two capped boundary distances and draw the mean-zero
boundary coin. Update $x$, record boundary events, and return to step 2.
\end{enumerate}
Every accepted direction passes the same curvature and flatness tests.
The finite-step proof controls both possible endpoints, so correctness
holds for every realized boundary choice.

Use $n$ for the original number of columns, $N$ for the active count at the
start of the current halving phase, $S$ for the number of coordinate
indices retained in the current cleanup epoch, and $a$ for the current
number of physically unfrozen coordinates.  Thus a large-epoch direction
is selected with
\[
 N/2<a\le S\le N,\qquad S\ge10000.
\]
The full vector $x\in[-1,1]^n$ is always retained, including already frozen
coordinates. Its squared norm supplies a common energy measure across changes in
the retained coordinate set.

\subsection{Initialization, cleanup, and the stopping rule}\label{a:state}
If $n\le99$, output all positive signs and stop; this includes the empty
input. Otherwise put $x=0$, set
\[
 J=1+\lceil\log_2 n\rceil,\qquad N=n,\qquad
 \eta=\frac1{4096JN},\qquad\delta=\frac{\eta^5}{2^{60}},
\]
and set $C=1200$. Fix, for the entire run,
\[
 Q=\left\lceil\log_2\bigl(16C(n+2m)+2\bigr)\right\rceil,\qquad q=Q.
\]
Thus $q\ge J$. Initialize the retained-row state and perform the cleanup of
Lemma~\ref{b:lem:cleanup}. At the start of each epoch its retained indices
are exactly the currently unfrozen coordinates and $S=a$.

Throughout a large epoch, keep the retained sets, offsets, row classes,
scales, and coordinate indices fixed. A coordinate which reaches a sign
is held in these data and imposes the exact equation $h_j=0$. Its retained
contributions are absorbed into offsets only at cleanup. Cleanup is
triggered before selecting another direction if at least $\lceil S/100\rceil$
coordinates have frozen in the epoch, if $a\le N/2$, or if the algorithm is
entering the small-dimension continuation described below.

At any cleanup, first absorb all frozen retained entries into their row
offsets and restrict to the remaining coordinates. Then perform the
specified row-class transitions, medium-row deliberate deletions, and
scale decreases, always using the preceding lexicographic conventions.
Discard the old anchors and initialize the new epoch's anchors from its
new barriers. If the optional Gram cache is used, rebuild it for the new
profiles and scales. This cleanup does not itself change $\eta$.

After cleanup, stop the walk when $a\le99$ and round every remaining
coordinate to its nearest sign, with zero rounded to $+1$. If $a>99$ and
$a\le N/2$, start the next phase by setting
\[
 N\leftarrow a,\qquad\eta\leftarrow\frac1{4096JN},\qquad
 \delta\leftarrow\frac{\eta^5}{2^{60}}.
\]
An outer reset also starts a new epoch. If both cleanup and a phase reset
are required, cleanup is performed with the old ridge first and the ridge
is increased afterwards. There is no restart of $x$ or of the retained
row offsets.

When fewer than $10000$ unfrozen coordinates remain, perform cleanup,
discard the anchors and optional caches, and use the baseline rational
current-state direction routine for all remaining moves. Its cleanup is
performed after every move, its phase parameter still follows the above
halving rule, and its stopping rule remains $a\le99$. The routine continues from the current feasible state. It can use
the baseline deterministic boundary sign. Alternatively it can use the
boundary coin below; both signs have the required finite-step guarantee.
This fallback has bounded dimension and costs
$\widetilde O(m+1)$ in the original input parameters. The polylogarithms
still include the original $n$, because $J$ and the requested precision
are retained.

\subsection{Exact row constraints at a large-epoch state}\label{a:rows}
For each active anchor retain its activation point $A_i$ and the exact
normal $g_i(A_i)$. At the start of a direction-selection state, release all
anchors with
\[
 E_i=\sum_jz_{ij}(x_j-A_{ij})^2>1/350,
\]
then activate every inactive medium row-sign pair with $u_i\ge-1/100$,
setting $A_i=x$. Released pairs have $u_i<-1/100$, so they are not
immediately reactivated. Equality in the release test retains the anchor;
equality in the activation test activates it. Apply simultaneous events
in least-index order.

Maintain the exact common kernel of the large-row normals, active anchor
normals, and held frozen-coordinate equations. Append the two current
rows
\[
 x^Th=0,\qquad
 \left(\sum_{i\text{ active}}\operatorname{diag}(z_i)(x-A_i)\right)^Th=0
\]
and denote the resulting exact orthogonal projector by $P$.
Dependent rows are handled by the basis-exchange rules in
Theorem~\ref{m:thm:rowcost}; they are not discarded from the logical list
of active constraints. All current medium-row gradients, including those
of inactive and anchored pairs, are still used in the flatness tests.
The dimension certificate is Lemma~\ref{m:lem:dimension}.

\subsection{Parameter order and one certified direction}\label{a:direction}
Let $R$ be the number of current medium row-sign pairs. The test vectors
are the $S$ coordinate vectors and the $R$ normalized current gradients
$(10/17)g_i$. Use the fixed global $q=Q$ selected at initialization;
it satisfies $q\ge J$ and $2^q\ge16C(S+R)$. All of these test vectors
have norm at most one.
The following order fixes the step cap and leakage penalty before
determining the sampling distribution.

\begin{enumerate}
\item Set
\[
 \mu=\eta+4096/S,\qquad\chi=\mu/\eta,
 \qquad A_K=32C\chi q/S.
\]
If $A_K\ge1$, take $K=1$. Otherwise choose $K=2^{-k}$ for the
largest nonnegative integer $k$ with $K^2\ge A_K$. Only squared
comparisons are needed. The sampler's coordinate and normalized-gradient bounds imply
$|a_i|\le K$ and $b_i\le K^2$ for every accepted output. The unconditional
unit-direction bounds cover the cap $K=1$. In particular
\[
 \eta\le K\le1,\qquad K^2\le\min\{1,128C\chi q/S\}.
\]
These bounds do not depend on the ridge penalty.

\item Compute $H=\lceil\log_2(2/\eta)\rceil$ by integer doublings and
choose the largest power of two $\tau\le1$ satisfying
\[
 \tau\le2^{-20},\qquad K^2\tau\le2^{-20}\eta,
 \qquad K^2H^2\tau^2\le2^{-40}\eta.
\]
Set
\[
 e_* =\min\{K\eta/256,\eta\tau/128\}.
\]
The values $K,\tau,e_*$ are now fixed for this state, before any sampling.

\item Obtain a certified interval $[\ell,\Lambda]$ for the current true
Perron eigenvalue of width at most $\delta$, using
Proposition~\ref{p:certificate}. If $\Lambda\le1$, use the low
branch: put the curvature data $\mathcal S=\mathcal D=0$, $F=\eta I$, and
$E=0$. The sampler's system is then $\eta P$ on the row-feasible space,
so its solve is just multiplication by $\eta^{-1}$ on that space.

If $\Lambda>1$, the certificate implies $L>1-\delta>1/2$; in particular
there is no unresolved equality test at $L=1$. Choose $\rho_{\mathrm{ridge}}$ as the least
power of four not smaller than $512C\mu/e_*^2$. For this modular version keep the original
$\Lambda>1$ as the filter normalization endpoint. The secant
implementation of Section~\ref{d:ideal-factor} explicitly replaces
this scalar by a refined upper approximation to $\max\{1,L\}$. Refine the current Perron
certificate and nonnegative coefficient approximations to the accuracies
required for the operator in~\eqref{r:operator}. They must jointly satisfy the
curvature-proxy error budget and the stronger polynomial-filter error
budget. Form the arithmetic curvature data
$\mathcal S=(279/79)\overline{\mathcal U}$,
$\mathcal D=\overline{\mathcal D}$, and
\[
 F=\operatorname{diag}\bigl(\min\{\mathcal D_{jj},4096/S\}\bigr)+\eta I.
\]
Construct the bounded filter and its arithmetic operator
$E$ as in Section~\ref{r:filter}, with
\[
 d^2\le\eta/(1200\rho_{\mathrm{ridge}}),\qquad\varepsilon_E\le e_*/4.
\]
The high spectral rank count remains valid even when $L\le1$:
$\Lambda>1$ and $L\le\Lambda$ imply the required bound on the number of
eigenvalues exceeding $\Lambda/5$. No exact comparison to one is used.

\item With the arithmetic data frozen, define the single linear operator
\[
 \mathsf A=P(F+\rho_{\mathrm{ridge}} E^TE)P\big|_{\operatorname{ran}P}
\]
in the high branch, and $\mathsf A=\eta I$ on $\operatorname{ran}P$ in
the low branch. Every application of $E$ uses the same frozen rational
coefficients, vector approximation, scalar certificate, and polynomial;
$E^T$ is its actual transpose. Rejection of a random sample does not change
these data or the state.

Set $\zeta=\min\{\eta/(1000C),1/(100S)\}$. Approximate the diagonal square
roots of $F$ to the prescribed error $\zeta\eta/8$. Draw independent
Rademacher vectors $\xi,\omega$ and form
\[
 \widetilde b=P(\widetilde{F^{1/2}}\xi+\sqrt{\rho_{\mathrm{ridge}}} E^T\omega)
\]
in the high branch, omitting the second term in the low branch.
Here $\sqrt{\rho_{\mathrm{ridge}}}$ is rational because $\rho_{\mathrm{ridge}}$ is a power of four.
Obtain $y\in\operatorname{ran}P$ with the finite residual certificate
\[
 \|\widetilde b-\mathsf Ay\|^2\le\zeta^2\eta S/4.
\]
A failed residual test requires further solve work or a solver retry; it
never produces an accepted direction. The default is exact elimination
on a rational independent basis. It may be replaced by any certified
solver which terminates almost surely; all its work is charged below.

As an alternative in the real-RAM model with square roots, form the exact
arithmetic right-hand side
$b=P(F^{1/2}\xi+\sqrt{\rho_{\mathrm{ridge}}}E^T\omega)$ and use the
fixed polynomial sampler of Theorem~\ref{new:polysampler}. Its polynomial
and certified spectral enclosure are fixed before drawing these signs.
This alternative replaces the preceding solve and residual test; it uses
the same acceptance test below. The original secant controller in
Section~\ref{d:section} retains the residual-certified version.

Accept the sample precisely when
\[
 y^T(F-\mathcal S)y\ge S/(8C),\qquad
 |b_j^Ty|^2\le4q/\eta\quad\hbox{for every test vector }b_j.
\]
Otherwise discard this random sample and repeat this item with new signs.
On acceptance choose a positive power of two $r$ with
$1/4\le r^2\|y\|^2\le1$ and put $h=ry$.
Theorems~\ref{r:sampler} and~\ref{new:polysampler} prove success
probability at least $1/(4C)$ per trial for the respective implementations. Proposition~\ref{r:tolerance} supplies the first-derivative and curvature
bounds and the finite-step guarantees for both signs of this $h$.
\end{enumerate}

In the residual-certified version the square-root and exponential
evaluations above mean finite rational approximations with stated error bounds.
The polynomial version uses the explicitly permitted exact square roots.
Neither version uses exact transcendental or eigenvector primitives. If the input is rational, the exact row
projectors, accepted directions, step lengths, and states remain rational
when the solve output and all approximations are rational. No polynomial
bit-length bound is asserted.

\subsection{One accepted move and the boundary coin}\label{a:move}
Let $t_+$ and $t_-$ be the distances to the cube boundary in directions
$h$ and $-h$, each capped at $\tau$. The minimum is over unfrozen indices
with nonzero direction coordinate; held frozen indices have $h_j=0$.
Since $\|h\|\ge1/2$, both capped distances are strictly positive.
Independently of all previous sampling, choose
\[
 d=\begin{cases}
 t_+,&\text{with probability }t_-/(t_++t_-),\\
 -t_-,&\text{with probability }t_+/(t_++t_-),
 \end{cases}
 \qquad x\leftarrow x+dh.
\]
A uniform binary expansion and exact comparisons implement this coin with
at most two revealed random bits in expectation. Conditional on the entire
past and the accepted $h$, its displacement has mean zero and second
moment $t_+t_-$. The coin selects the endpoint without evaluating the spectral potential.

Mark all newly boundary coordinates frozen. If cleanup is triggered,
perform the state transition in Section~\ref{a:state}. Otherwise process
anchor releases and activations and update the fixed-row factorization
before selecting another direction. If the optional Gram cache is used,
refresh every coefficient whose rational log-weight increment since its
cache anchor has absolute value at least $1/8$. At a cleanup one may
instead discard the old cache directly and rebuild it; its nonnegative
counting residual need not be paid again. The next iteration begins at
Section~\ref{a:rows}.

Every realized move with $|d|<\tau$ freezes at least one new coordinate.
Every realized move with $|d|=\tau$ has the full-step descent guarantee,
regardless of its sign. Sampling rejection and solver retry do not count
as moves and do not change the potential or the row data.

\subsection{Global guarantees and the exact unresolved cost}\label{a:guarantees}
\begin{proposition}\label{a:budgets}
The modular algorithm terminates almost surely. Every returned signing
has discrepancy strictly below $D_*=105+3\sqrt{99}<135$. Along every
execution, at most $O(nJ^4Q^2)$ accepted moves occur, where
\[
 Q=\left\lceil\log_2\bigl(16C(n+2m)+2\bigr)\right\rceil.
\]
These discrepancy and accepted-move bounds hold pathwise, not merely in
expectation. The number of sampling trials, certified-solver work, and
optional weight refreshes are separately charged random quantities.
\end{proposition}
\begin{proof}
Either sampler, the polynomial filter, and Proposition~\ref{r:tolerance} verify
the direction and two-sided step contracts in
Section~\ref{m:sec:contracts}, with $q=Q$. The anchored row constraints
preserve feasibility for both endpoint choices. Thus
Proposition~\ref{m:prop:coin} gives the pathwise bound $L<9/8$ and at most
$O(nJ^4Q^2)$ accepted moves, including the phase ridge changes.
Cleanup retains the original full-vector quadratic-energy accounting.
The bounded-dimension continuation satisfies the same weaker move bound.

Each sampling trial has uniformly positive conditional success
probability. The certified Perron primitive terminates almost surely;
the solve interface can always be completed by exact elimination, and the
fixed polynomial alternative has finite prescribed length.
These facts give almost-sure termination. All retained-row invariants
remain valid, and terminal cleanup removes the medium rows before
rounding the at most $99$ remaining coordinates. The tracking estimate
and nearest-sign rounding in Proposition~\ref{b:prop:terminal} therefore
give discrepancy strictly below $D_*$ on every returned execution.
\end{proof}

\begin{proposition}\label{a:accounting}
Let $\mathcal I$ index all attempted random direction samples in the large
epochs, including rejected samples. For $\iota\in\mathcal I$, let
$\mathsf{Solve}_\iota$ be the total arithmetic work used to obtain that
sample's certified solution of $\mathsf A_\iota y=\widetilde b_\iota$,
including every internal solver retry or failed residual test. Then the
expected arithmetic cost of the modular algorithm is
\[
 \widetilde O(mn^2+n^3)
       +\mathbb E\sum_{\iota\in\mathcal I}\mathsf{Solve}_\iota.
\]
The same bound holds if the optional positive-shift cache is maintained.
No estimate for the second term by $\widetilde O(mn^2+n^3)$ is established
here.
\end{proposition}
\begin{proof}
The row-maintenance theorem gives $\widetilde O(n^3)$ total exact row
projector work. Current coefficient evaluation, Perron certification,
filter applications, right-hand-side construction, residual checks, and
flatness/curvature tests cost $\widetilde O(mS+S^2)$ per trial or state,
with a constant expected number of direction trials per accepted move. A phase
starting with $N$ has $O(NJ^4Q^2)$ accepted moves. Summing the geometrically
decreasing phase sizes therefore pays
$\widetilde O(mn^2+n^3)$ for these operations. Cleanup formation and scans
are included; the sum of large-epoch starting sizes is $O(n)$, so dense
reference rebuilds, when requested, have total cost
$O(m\sum S^2+\sum S^3)=O(mn^2+n^3)$.
The optional cache's expected refresh count is charged by
Theorem~\ref{m:thm:cachecost}, using exactly the mean-zero boundary coin
above. The bounded-dimension continuation has additional
$\widetilde O(m+1)$ cost. The remaining operation is the certified solve
with the changing ridge matrix, which has been displayed explicitly and
charged for every attempted sample.
\end{proof}

The certified-solve interface has a finite fallback: form the current
arithmetic operator and solve on an independent row-feasible basis by
exact elimination. This cost belongs to the displayed solve term.
The deterministic algorithm of Theorem~\ref{b:thm:main} remains a
separate implementation with its own unconditional arithmetic bounds.


\section{Certified endpoint steps and a Perron-free cubic regime}
\label{ep:section}
The implementation note~\cite{CertSteps} reports that testing cube
endpoints can replace the prescribed tiny steps on the tested instances.
The subsequent analysis~\cite{CertAnalysis} identifies inexpensive
certificates for those endpoints. We incorporate that analysis here,
including the stopping safeguards and the full constraint-maintenance
cost. The resulting conditional cubic theorem has an observable
hypothesis: cheap endpoint trials succeed and freeze coordinates
throughout the run. It requires neither the production hypothesis
\eqref{d:conjecture} nor a bound on a sampler condition number.

\subsection{Endpoint feasibility and a direct inactive-branch certificate}
Use the state, cleanup, and phase rules of
Section~\ref{b:sec:state}, and write $\mathcal E(x)=\|x\|^2$, including
frozen coordinates. The discrepancy argument uses less information than
the finite-step argument.

\begin{lemma}[Endpoint correctness]\label{ep:correctness}
Start from the initialized state. Before each cleanup, require only
\begin{equation}\label{ep:feasibility}
 \begin{gathered}
 \Delta_j=0\quad(j\notin V),\qquad x'=x+\Delta\in[-1,1]^n,\\
 b_i^\T\Delta=0\quad(i\text{ large}),\qquad
 u_{i\sigma}(x')\le0\quad(i\text{ medium},\ \sigma=\pm1).
 \end{gathered}
\end{equation}
Every such sequence that reaches $s\le99$ gives discrepancy less than
$D_*$ after cleanup and nearest-sign rounding. This conclusion requires
neither a spectral certificate nor feasibility along the open segment
from $x$ to $x'$.
\end{lemma}
\begin{proof}
The large-row equations preserve $d_i=0$ until a row becomes medium.
Lemma~\ref{b:lem:cleanup} preserves the endpoint barriers and tracked
sums through cleanup; its barrier argument needs only
\eqref{ep:feasibility}. The tracking identity of
Lemma~\ref{b:lem:tracking} holds for arbitrary moves inside the cube.
Proposition~\ref{b:prop:terminal} therefore applies at termination.
\end{proof}

In particular, one should test the cube endpoint before truncating at
the first barrier root. With the retained data fixed,
\begin{equation}\label{ep:barriercurve}
 u_{i\sigma}(x+th)=u_{i\sigma}(x)+t g_{i\sigma}^\T h
       -\frac72t^2\sum_j\frac{p_{ij}}{r_i^2}h_j^2.
\end{equation}
A concave quadratic can be nonpositive at both endpoints and positive
between them. For example, $-1/100+(2/5)t-t^2/10$ is negative at $0$
and $4$, but positive at $2$. Endpoint correctness permits that move.

\begin{proposition}[Tracked-sum certificate]\label{ep:envelope}
At a cube-contained state with the current retained data, if every
medium row satisfies $|d_i|\le D\le129/2$, then
\begin{equation}\label{ep:envelope-bound}
 L\le\eta s+\max\left\{\frac1{128},
      \frac{256}{9}\exp\left(\frac{4D-174}{21}\right)\right\}.
\end{equation}
In particular, $|d_i|\le25$ for all medium rows implies
\[
 L<\frac{21}{25}=0.84,\qquad u_{i\sigma}\le-\frac{37}{3}.
\]
Thus the tracked-sum, cube, and large-row tests alone certify a feasible
inactive-branch endpoint using rational arithmetic. No exponential
weights, eigenvectors, or spectral filters are evaluated.
\end{proposition}
\begin{proof}
Since $G_i\le q_i\le r_i^2$, we have
$u_{i\sigma}\le(D-75)/r_i+21$. Cauchy--Schwarz gives
$p_ip_i^\T\preceq q_i\diag(p_i)$, so the two medium Gram terms are
bounded above by
\[
 \frac{64}{r_i^2}
   \exp\left(6-\frac{2(75-D)}{7r_i}\right)\diag(p_i).
\]
The scalar coefficient increases with $r_i\le3/2$ when $D\le129/2$:
its logarithmic derivative is
$-2/r_i+2(75-D)/(7r_i^2)\ge0$. Substitute $r_i=3/2$, combine with
the large-row coefficient $1/128$, and use $\sum_i p_{ij}\le1$.
The ridge has norm $\eta s$, proving \eqref{ep:envelope-bound}.
At $D=25$, the bound is at most
$(256/9)e^{-74/21}+1/4096=0.83898234\ldots<21/25$.
For an entirely rational verification of this strict inequality,
replace $e^{74/21}$ by its positive Taylor sum through degree $11$.
Finally $u_{i\sigma}\le21-50/r_i\le-37/3$.
\end{proof}

The hypothesis concerns tracked sums at the candidate endpoint,
before cleanup, rather than final discrepancy. Cleanup preserves these
sums, and newly medium rows have $d_i=0$. Hence the certificate remains
available after cleanup and after a phase reset, when the new ridge
still satisfies $\eta s\le1/4096$.

A less restrictive, one-pass certificate is the exact row-sum bound
\begin{equation}\label{ep:rowsum}
 L\le\Lambda_{\one}:=\eta s+
 \max_j\left\{\frac1{128}\sum_{i\text{ large}}p_{ij}
       +\sum_{i\text{ medium}}q_i(w_{i+}+w_{i-})p_{ij}\right\}.
\end{equation}
It is $\max_j(M\one)_j$ and costs $O(ms+s)$ operations once the
weights are supplied. Rational upper weight approximations give a
valid upper certificate. When the tracked-sum test passes, these
approximations are unnecessary: use $\Lambda=21/25$ directly.

\subsection{A bounded hybrid with the deterministic fallback}
\label{ep:hybrid}
Endpoint feasibility alone does not guarantee progress. The following
wrapper retains the complete deterministic guarantee while allowing any
cheap candidate routine. Define the certified displacement floor
\begin{equation}\label{ep:floor}
 B_*=2\cdot4096\cdot25600Jq,\qquad
 \ell_*:=\frac{3\cdot2^{-20}}{8B_*}.
\end{equation}
Proposition~\ref{b:prop:count} proves that every full deterministic
fallback move has length at least $\ell_*$. Keep a rational ledger
$\mathfrak b\ge\Phi$, initially $0$, and use the following rules.
\begin{enumerate}
\item Fix finite caps on candidate construction, endpoint trials,
certificate products, and scalar refinement before beginning the run.
Fix their order, with least-index and positive-sign tie breaking.
Test the two cube endpoints first; if shorter dyadic trials are used,
cap their number as well. For example, for unit candidates the cap
$2+\lceil\log_2(2\sqrt n/\ell_*)\rceil$ per sign suffices to cover
the range down to $\ell_*$. Normalization is unnecessary when only
cube endpoints are tested.
\item At a proposed endpoint $x'$, check \eqref{ep:feasibility} and
$\mathcal E(x')\ge\mathcal E(x)$. If no coordinate freezes, also
require $\mathcal E(x')-\mathcal E(x)\ge\ell_*^2$.
Obtain a rational upper certificate $\Lambda\ge L(x')$ and accept only if
\begin{equation}\label{ep:ledger}
 \widehat\Phi=(\Lambda-1)_+^3-\eta\mathcal E(x')\le\mathfrak b.
\end{equation}
On acceptance set $\mathfrak b\leftarrow\widehat\Phi$.
Proposition~\ref{ep:envelope} can supply both the row-barrier and
spectral certificates in this step.
\item If the optional work cap is exhausted, take one deterministic
step from Section~\ref{b:sec:progress}. Add $-\eta\tau^2/32$ to
$\mathfrak b$ for a full step, or $+\eta\tau^2/128$ for a shorter
boundary step, as justified by Proposition~\ref{b:prop:onestep}.
\item Perform the original cleanup. At a phase reset, with the new
ridge $\eta'$, set
\begin{equation}\label{ep:reset}
 \mathfrak b\leftarrow\mathfrak b+
     \frac{3}{16\cdot4096J}-(\eta'-\eta)\mathcal E(x).
\end{equation}
Stop at $s\le99$ and apply nearest-sign rounding. For $n\le99$ use
the original immediate completion.
\end{enumerate}
The optional routine may use projected vectors without any negative
curvature or leakage calculation. In particular, a direction satisfying
the large-row equations and $x_V^\T h=0$ has energy gain
$t^2\|h\|^2$ for either endpoint. The exact projector machinery in
Section~\ref{m:subsec:projector} handles dependent rows.

\begin{theorem}[Bounded endpoint hybrid]\label{ep:wrapper}
The wrapper terminates and returns discrepancy less than $D_*$ for
every choice of its optional candidates. All cleaned states, including
phase resets, satisfy $L<9/8$. For $n\ge100$ it makes at most
\begin{equation}\label{ep:movecount}
 N\le n-99+\frac{n}{\ell_*^2}=O(nJ^2q^2)
\end{equation}
moves. If every move freezes a coordinate, then $N\le n-99$.
Capping the optional work at a polylogarithmic multiple of the
deterministic per-state budget preserves the deterministic arithmetic
bound in Theorem~\ref{b:thm:main}, with the same arithmetic primitives.
\end{theorem}
\begin{proof}
The ledger remains an upper bound on $\Phi$ by \eqref{ep:ledger},
the fallback descent estimates, cleanup monotonicity, and the reset
charge from Proposition~\ref{b:prop:budget}. Within a phase starting
with $S$ coordinates, the energy can increase by at most $S$ and there
are at most $S$ shorter fallback moves. Optional moves never increase
the ledger. Thus every phase prefix increases the bound on
$(L-1)_+^3\le\mathfrak b+\eta\mathcal E$ by at most
$129/(128\cdot4096J)$. Each reset adds at most
$3/(16\cdot4096J)$. The same induction as in
Proposition~\ref{b:prop:budget} gives
$(L-1)_+^3\le153/(128\cdot4096)<1/512$.
Every move increases energy nonnegatively, and every nonfreezing move
increases it by at least $\ell_*^2$. Total energy gain is at most $n$;
there are at most $n-99$ freezing moves before the stopping rule.
This proves termination and \eqref{ep:movecount}.
Lemma~\ref{ep:correctness} gives discrepancy. The same argument within
a phase bounds its moves by $S+S/\ell_*^2$; summing the capped
per-state work over geometrically decreasing phase sizes proves the
arithmetic assertion. A dense optional spectral computation must obey
the chosen cap; it does not inherit the fast-arithmetic clause for free.
\end{proof}

The wrapper does not use the modular boundary coin. Consequently it
does not inherit either the lazy-cache martingale estimates or the
secant production hypothesis, which concern their specified probability
laws. Its guarantees above use the deterministic fallback.

\subsection{The conditional cubic theorem, including projector construction}
\label{ep:cubic-section}
Let $p_{\rm try}\ge1$ bound the number of cheap construction and
verification passes allowed per accepted move, including failed
candidate projections and rejected endpoint trials. This is a work
cap, not merely a bound on accepted steps.

\newpage
\begin{theorem}[Perron-free cubic regime]\label{ep:cubic}
Use the bounded wrapper with optional directions in the exact kernel
of the large-row normals and $x_V^\T$. Suppose that throughout an
execution with $n\ge100$:
\begin{enumerate}
\item each state produces an accepted cube endpoint within
$p_{\rm try}$ passes, with no fallback;
\item every accepted endpoint passes $|d_i|\le25$ for all currently
medium rows, or an upper row-sum certificate $\overline\Lambda_{\one}\le1$;
\item candidate construction and testing consist of a bounded number
of row scans and projector applications, together with $O(ms+n^2)$
scalar and vector operations. Any approximations for the row-sum alternative
are certified with inverse-polynomial accuracy within the pass budget.
\end{enumerate}
Then the run returns discrepancy less than $D_*$ in at most $n-99$
moves and
\begin{equation}\label{ep:cubic-cost}
 \widetilde O\bigl(p_{\rm try}(mn^2+n^3)\bigr)
\end{equation}
arithmetic operations, including projector construction, dependent-row
handling, cleanup, and all rejected trials. It computes no Perron
vector, eigengap, high spectral subspace, or filtered sampling system.
With the tracked-sum certificate alone, it also evaluates no
exponentials; rational input admits field arithmetic throughout.
For $m=\Theta(n)$ and $p_{\rm try}=\operatorname{polylog}(2+m+n)$,
the bound is $\widetilde O(n^3)$.
\end{theorem}
\begin{proof}
Cube endpoints freeze coordinates, and $x_V^\T h=0$ gives
nonnegative energy gain. Each accepted endpoint is row-feasible by
the wrapper and has $L\le1$. Hence its exact potential is
$-\eta\mathcal E(x')\le-\eta\mathcal E(x)\le\Phi(x)\le\mathfrak b$;
the ledger test costs no further spectral work. Correctness and the
move count follow from Theorem~\ref{ep:wrapper}.

Here the maintenance cost has a direct proof that does not assume the
sampler's flatness or anchored-row contracts. Initially there are fewer
than $4n/9$ large rows, since each has $q_i>9/4$ and
$\sum_iq_i\le n$. A row can leave this class only once and can never
enter it later. Maintain their original normals in $\R^n$, together
with equations $h_j=0$ for frozen coordinates. Large rows undergo no
deliberate deletion, so this fixed-normal representation gives exactly
the required constraints on active increments. There are $O(n)$
logical rows and $O(n)$ row-exit or coordinate-freezing events.

Construct an independent basis and its projector once in $O(n^3)$
operations. Equations~\eqref{m:eq:insert}--\eqref{m:eq:delete} maintain
them, including dependent-row records and basis exchanges, in $O(n^2)$
work per event. Temporarily adjoining the energy row and applying the
projector costs $O(n^2)$ per pass. This kernel is nonzero for $s\ge100$:
the current large rows have rank less than $4s/9$, and the energy row
adds at most one. If a trial projection vanishes it is rejected within
the cap; a nonzero projector column can alternatively be selected by
the least positive diagonal entry, without a spectral computation.

Candidate endpoints, tracked sums, barriers, and row-sum certificates
cost $\widetilde O(ms+n^2)$ per pass. The active sizes decrease at
every accepted move, so $N\le n-99$ and
$\sum_t s_t\le n(n+1)/2$. The projector passes cost
$O(p_{\rm try}n^2N)$, and all permanent projector events cost $O(n^3)$.
An initial input scan costs $O(mn)$. Cleanup can maintain row masses
and per-row trees of coefficient maxima: each retained entry is removed
once, while current row scans and scale searches fit within the
displayed pass bound, up to logarithms. Terminal rounding costs
$O(mn)$. Summing these costs proves \eqref{ep:cubic-cost}.
\end{proof}

The theorem is conditional on the observable success of the cheap
endpoint routine for the whole execution. Neither the dimension count
nor a small final discrepancy implies that success. The certificate
changes how progress is obtained: it permits boundary moves directly,
without invoking the high-branch direction construction. Relative to
the dense bound $\widetilde O(mn^3+n^4)$, these executions save one
polynomial factor of $n$.

\subsection{Positive certificates beyond the inactive branch}
\label{ep:positive}
For a proposed endpoint, the ledger permits any positive rational
target $\mu$ satisfying
$(\mu-1)_+^3\le\mathfrak b+\eta\mathcal E(x')$.
Testing $My\le\mu y$ for a positive $y$ certifies $L\le\mu$ without
approximating a Perron vector. If the right-hand side of the target
inequality is negative, reject the optional endpoint. Otherwise a
finite prescribed list of rational targets avoids any exact cube-root
operation.

\begin{lemma}[Certificate cost depends on target slack]\label{ep:series}
Let $M$ be symmetric, PSD, and entrywise nonnegative. For $\mu>0$ set
\[
 Q=M/\mu,\qquad u_0=\one,\quad y_0=0,\qquad
 y_{k+1}=y_k+u_k,\quad u_{k+1}=Qu_k.
\]
For any $k\ge1$, the componentwise test $u_k\le\one$ certifies
$My_k\le\mu y_k$. If $L\le(1-\zeta)\mu$ with $0<\zeta<1$, it
succeeds within
\begin{equation}\label{ep:series-cost}
 k\le\max\left\{1,
 \left\lceil\frac{\log\sqrt s}{-\log(1-\zeta)}\right\rceil\right\}
 =O\bigl(1+\zeta^{-1}\log(2+s)\bigr)
\end{equation}
factored matrix-vector products, independently of $L-\lambda_2(M)$.
\end{lemma}
\begin{proof}
The vector $y_k=\sum_{j=0}^{k-1}Q^j\one$ is positive and
$My_k=\mu(y_k-\one+u_k)$. The test implies the stated inequality.
It implies $L\le\mu$, for example by taking the inner product with a
nonnegative top eigenvector. Finally,
$\|u_k\|_\infty\le\|Q^k\one\|_2\le(1-\zeta)^k\sqrt s$.
\end{proof}

For finite arithmetic, round each nonnegative Gram coefficient upward
to obtain a PSD matrix $\overline M\ge M$ entrywise, and verify
$\overline My\le\mu y$ exactly. If
$\|\overline M-M\|\le\zeta\mu/2$, the lemma applies with slack
$\zeta/2$. In the representation
$W=\sum\alpha_i z_i z_i^\T$, $z_i=p_i/r_i^2$, feasible medium rows
have $\|z_i\|^2\le1/100$ and $0\le\alpha_i\le32$.
Thus absolute coefficient accuracy $O(\zeta\mu/(m+1))$ suffices.
Lemma~\ref{b:lem:scalar} supplies such bounds with logarithmic scalar
accuracy cost. The number of matrix products is nevertheless
proportional to $1/\zeta$, not to $\log(1/\zeta)$.
Cap both refinement and product counts; failure to certify in the
allotted work triggers the fallback. Finite success is not promised
at zero target slack.

One can also exploit influential rows directly. Along a fixed direction
write the exact curve, indexing medium row-sign pairs by $i$, as
\[
 M(t)=C+\sum_i w_i e^{a_it-b_it^2}p_ip_i^\T,
 \qquad a_i=\tfrac27g_i^\T h,\quad
 b_i=\sum_j(p_{ij}/r_i^2)h_j^2\ge0,
\]
where $C$ is the fixed large-row diagonal plus ridge. Impose
$g_i^\T h=0$ on a selected set $\mathcal H$ of pairs. Their weights
then decrease for both signs and at all lengths. For $y>0$, put
\[
 C_j=\frac{(Cy)_j}{y_j},\qquad
 H_j=\sum_{i\in\mathcal H}\frac{w_i p_{ij}(p_i^\T y)}{y_j},\qquad
 R_j=\sum_{i\notin\mathcal H}\frac{w_i p_{ij}(p_i^\T y)}{y_j}.
\]
For $b_{\mathcal H}\le\min_{i\in\mathcal H}b_i$ and
$z_R(t)\ge\max(0,\max_{i\notin\mathcal H}(a_it-b_it^2))$,
entrywise positivity gives the certificate
\begin{equation}\label{ep:protected}
 L(t)\le\max_j\{C_j+e^{-b_{\mathcal H}t^2}H_j+e^{z_R(t)}R_j\}.
\end{equation}
For an empty protected set take $b_{\mathcal H}=0$.
Keeping the individual exponents gives an at least equally sharp
positive-vector bound. This uses the whole protected subspace, rather
than a single scalar curvature estimate. Selecting its rows and
maintaining their changing projector remain additional charged work;
\eqref{ep:protected} alone does not establish a cubic high-branch bound.

\subsection{Existing implementation evidence and the scope of the speed claim}
\label{ep:evidence}
The implementation note~\cite{CertSteps} reports complete natural-instance
runs for six families through $n=2000$, including rectangular inputs,
with no fallback and at most $n-99$ moves. The implementation scripts are
in \texttt{algorithm/} of this repository and the stored run records in
\texttt{results/}; their floating-point scope is unchanged. The separate stored checks
in~\cite{CertAnalysis} use four square-input families at
$n=128,256,512$: column-normalized Gaussian, normalized Rademacher,
normalized Hadamard, and Gaussian with one eighth of the rows amplified
by six before column normalization. These twelve previously completed
runs test both cube endpoints of a projected random direction and use
the row-sum certificate \eqref{ep:rowsum}. They implement the retained
sets, offsets, deletions, scales, phase resets, cleanup, and terminal
rounding, and abort if neither endpoint passes. They do not implement
the high-branch fallback. Those records were retained without rerunning that particular suite.
The new matched profile comparison is reported separately in
Section~\ref{opt:comparison}.

\begin{table}[ht]
\centering\small
\caption{Stored complete low-branch runs: four families at each size.
Maxima are over the recorded runs and states; values are rounded.}
\label{ep:table}
\begin{tabular}{@{}rrrrr@{}}
\toprule
$n=m$ & Moves in every run & Largest $|d_i|$ & Largest $\Lambda_{\one}$
 & Final discrepancy range\\
\midrule
128 & 29  & 1.927 & 0.007264 & $[2.672,5.583]$\\
256 & 157 & 2.753 & 0.007124 & $[2.450,5.939]$\\
512 & 413 & 4.779 & 0.007024 & $[2.747,6.872]$\\
\bottomrule
\end{tabular}
\end{table}

Every recorded move froze a coordinate, so each run used exactly
$n-99$ moves. The tracked sums stayed well below $25$, and therefore
the analytic certificate also explains the observed inactive branch.
The scripts check tracking identities, large-row constraints before
completion, and energy increments; the largest recorded orthogonality
error is below $1.2\cdot10^{-15}$. Some small Gaussian instances lose
all medium rows during initial deletion and provide no evidence about
the active spectral regime.

These are floating-point traces, not interval-certified executions of
the exact arithmetic algorithm. Some projector factorizations are
recomputed in the scripts. They support the endpoint-success pattern
in Theorem~\ref{ep:cubic}; the theorem supplies the maintained
implementation and its cubic operation bound. They do not measure a
cubic wall-clock exponent. In particular, maxima of timings over
different instance and seed collections, as reported in the supplied
note, are not a matched scaling study. The supported speed statement is
a cubic-work implementation for this observed low-branch pattern,
saving one factor of $n$ over the dense baseline.

The high branch needs separate accounting. If $\mathcal I_{\rm high}$
indexes dense direction constructions, including retries, their cost is
\begin{equation}\label{ep:high-cost}
 \sum_{t\in\mathcal I_{\rm high}}\widetilde O(ms_t^2+s_t^3).
\end{equation}
Linear many moves can still leave this charge quartic for square
inputs. The reported $27$--$55$ high-branch moves at $n=600$--$800$
do not control its asymptotic growth. A cheap positive certificate
replaces verification; it does not itself replace an expensive
direction construction. Likewise, the protected-row tests include
supplied local states with gap $286.2167\ldots\eta$ and a certified
cube step of length $16$, but their reachability from zero is unproved.
These observations motivate the certificates, without extending
Theorem~\ref{ep:cubic} to unrestricted inputs. Nor does embedding a
Gram--Schmidt direction transfer a theorem for the ordinary
Gram--Schmidt walk: the added constraints and endpoint selection
change its probability law.


\section{An implementation parameterized by conditioning}\label{ifc:sec:kappa}
An explicit dependence on conditioning permits a complete running-time
statement for the full walk. We apply the classical conjugate-gradient estimate
\cite[Section~6.11.3]{Saad} to the sampler's certified systems.

At sampling trial $t$, let $s_t\le n$ be the active ambient dimension,
$H_t$ the exactly maintained row-feasible subspace, and $P_t$ its
orthogonal projector. The finite arithmetic sampler system is
\[
 A_t=P_t(F_t+\rho_t E_t^TE_t)P_t\big|_{H_t}\succ0,
 \qquad b_t\in H_t.
\]
All coefficients defining this operator are frozen during the trial.
Let $c_t=\widetilde O(ms_t+s_t^2)$ be the cost of applying $A_t$ and
checking its residual. The polynomial filter, transpose action, exact
projection, and finite coefficient accuracy are included in this cost.
Proposition~\ref{a:accounting} gives
\begin{equation}\label{ifc:eq:kappa-base}
 \mathbb E\sum_t c_t=\widetilde O(mn^2+n^3),
\end{equation}
including rejected sampling trials. Write $\delta_t>0$ for the requested
absolute residual tolerance. The original sampler permits
$\delta_t=\zeta_{0,t}\sqrt{\eta_ts_t}/2$ for its computed right-hand
side, and any smaller tolerance preserves its proof.

Optionally choose a symmetric positive definite (SPD) preconditioner
$B_t$ on $H_t$, fixed during the trial. Its inverse must map $H_t$ into itself. Let $d_t$ be the cost
of one inverse application, and let $S$ be the total cost of constructing,
updating, and certifying the preconditioners during the run. The choice
$B_t=I_{H_t}$ requires no preconditioner construction. Define
\[
 \kappa_t^A=\frac{\lambda_{\max}(A_t)}{\lambda_{\min}(A_t)},
 \qquad
 \kappa_t^{\rm eff}=
 \frac{\lambda_{\max}(B_t^{-1/2}A_tB_t^{-1/2})}
      {\lambda_{\min}(B_t^{-1/2}A_tB_t^{-1/2})}.
\]
The square roots here define the condition number; the implementation
uses only $B_t^{-1}$, not $B_t^{-1/2}$. For $b_t\ne0$, put
\begin{equation}\label{ifc:eq:kappa-q}
 \epsilon_t=\min\{1/2,\delta_t/\|b_t\|\},\qquad
 q_t=1+\min\left\{\dim H_t,
  \left\lceil\sqrt{\kappa_t^{\rm eff}}
       \log\frac{2\sqrt{\kappa_t^A}}{\epsilon_t}\right\rceil\right\}.
\end{equation}
If $b_t=0$, return zero immediately and set $q_t=1$.

\begin{theorem}[Condition-number-sensitive implementation]\label{ifc:thm:kappa}
Preconditioned conjugate gradients, started at zero and stopped by the
actual residual test $\|b_t-A_ty\|\le\delta_t$, implements every sampler
solve with the required exact feasibility and residual certificate.
In exact real arithmetic the expected cost of the complete signing
algorithm is
\begin{equation}\label{ifc:eq:kappa-sum}
 \widetilde O(mn^2+n^3)+
 O\!\left(\mathbb E\left[S+\sum_t q_t(c_t+d_t)\right]\right).
\end{equation}
In particular, use $B_t=I_{H_t}$ and let $\kappa\ge1$ be any uniform
bound on $\kappa_t^A$ over trials the algorithm can visit. If a parameter
$0<\epsilon\le1/2$ tightens each original finite-system residual
tolerance by a factor $\epsilon$, then the expected total cost is
\begin{equation}\label{ifc:eq:kappa-uniform}
 \widetilde O\!\left((mn^2+n^3)
      \min\left\{n,\sqrt\kappa\log\frac{2\kappa}{\epsilon}\right\}\right).
\end{equation}
Here the soft notation suppresses the previously established dimension
and input-size logarithms; the displayed conditioning and solve-accuracy
dependences are explicit. No numerical upper bound for $\kappa$ is
required to run the residual-stopped algorithm. Its discrepancy guarantee
is unchanged, and no controller-production or interface-rank hypothesis
is assumed.
\end{theorem}
\begin{proof}
All CG iterates remain in $H_t$: the initial point, right-hand side,
matrix action, and preconditioner action preserve that space. Thus no
new basis or constraint factorization is necessary. Suppress the index
$t$, put $y_*=A^{-1}b$, and let
$a=(\sqrt{\kappa^{\rm eff}}-1)/(\sqrt{\kappa^{\rm eff}}+1)$.
Classical preconditioned CG gives
\[
 \|y_k-y_*\|_A\le2a^k\|y_*\|_A,
 \qquad
 \|b-Ay_k\|\le2\sqrt{\kappa^A}\,a^k\|b\|.
\]
For the second inequality, use
$\|Ae\|\le\sqrt{\lambda_{\max}(A)}\|e\|_A$ and
$\|y_*\|_A\le\|b\|/\sqrt{\lambda_{\min}(A)}$.
If $\kappa^{\rm eff}>1$, then
$a\le\exp(-2/\sqrt{\kappa^{\rm eff}})$; if it equals one, one step
solves the system. These facts give the logarithmic bound in
\eqref{ifc:eq:kappa-q}. Independently, CG terminates within $\dim H_t$
steps in exact arithmetic. A breakdown before then can occur only when the exact solution has
already been found; the residual test detects this case.

Each iteration and residual check costs $O(c_t+d_t)$; vector operations
are absorbed since $c_t\ge s_t$. Preparation and changes to $B_t$ are
charged to $S$. The original rejection and move-count proofs depend only
on the certified residual, so they remain valid. Summing proves
\eqref{ifc:eq:kappa-sum} without assuming any independence between the
condition numbers and the adaptive trajectory.

For the uniform specialization, the original residual tolerances and
right-hand-side norms imply
$\log(1/\epsilon_t)=O(\log(1/\epsilon)+\operatorname{polylog}(m+n+2))$.
This follows from the inverse-polynomial ridge and scalar tolerances,
$\|E_t\|\le3$, $F_t\preceq I$, and the polynomial penalty and
right-hand-side bounds already proved in Section~\ref{a:algorithm}.
Use $\dim H_t\le n$, the uniform condition bound, and
\eqref{ifc:eq:kappa-base} in \eqref{ifc:eq:kappa-sum}. This gives
\eqref{ifc:eq:kappa-uniform}. In particular, a square instance costs
$\widetilde O(n^3\sqrt\kappa\log(2\kappa/\epsilon))$, with the sharper
$n$-iteration cap retained in the displayed statement.
\end{proof}

\paragraph{Preconditioning and adaptive conditioning.}
If $d_t=O(c_t)$ and $\kappa_t^{\rm eff}\le\kappa$ uniformly, the same
argument replaces the leading square-root dependence by
$\sqrt\kappa$ and adds $\mathbb ES$. Logarithms of the original
$\kappa_t^A$ remain in the Euclidean residual conversion in
\eqref{ifc:eq:kappa-q}. Approximate inverse applications need their own
fixed-operator or inexact-iteration justification; their accuracy and
cost are not implicit in the statement. When no uniform deterministic
condition bound is used, \eqref{ifc:eq:kappa-sum} is the appropriate result:
the random trajectory's condition numbers remain inside the expectation.

Two elementary checks clarify the accounting. Taking $B=A$ makes the
iteration count one, but its construction and inverse action are still
charged. Conversely, $A=\alpha I+\beta uu^T$, with $\alpha>0$,
$\beta\ge0$, and $\|u\|=1$, has only two eigenvalues and needs at most
two unpreconditioned CG steps, however small $\alpha$ is. Thus the
condition-number estimate is an upper bound, not a claim that a small
ridge forces slow convergence. No eigendecomposition is used in either
the implementation or its stopping rule.

\subsection{A fixed-length implementation of the polynomial sampler}
\label{new:polyimplementation}
We now apply classical Chebyshev iteration~\cite[Section~12.3.2]{Saad}
at the constant accuracy justified by Theorem~\ref{new:polysampler}.
The iteration uses the certified spectral enclosure and operator
applications; it needs no inverse or preconditioner.

\begin{proposition}[Cost of constant-accuracy sampling]\label{new:polycost}
At trial $t$, supply certified bounds
$0<\alpha_t I\preceq A_t\preceq\beta_t I$ on $H_t$, independent of the
fresh signs. Set $\bar\kappa_t=\beta_t/\alpha_t$. The polynomial sampler
can be implemented with
\[
 O\!\left(1+\min\{\dim H_t,\sqrt{\bar\kappa_t}\log(200C)\}\right)
\]
operator applications and the same order of vector work. If
$\bar\kappa_t\le\bar\kappa$ uniformly, the full signing algorithm has
expected arithmetic cost~\eqref{i:poly}. This statement uses the existing
Perron primitive of Proposition~\ref{p:certificate}.
The enclosure $[\eta,\mu+9\rho_{\mathrm{ridge}}]$ is always available
from the original parameter schedule; better enclosures must be certified
and their acquisition cost included.
\end{proposition}
\begin{proof}
Suppress $t$ and first suppose $\beta>\alpha$. For the first-kind
Chebyshev polynomial $T_k$, define
\[
 R_k(\lambda)=
 \frac{T_k((\beta+\alpha-2\lambda)/(\beta-\alpha))}
      {T_k((\beta+\alpha)/(\beta-\alpha))},\qquad
 p_{k-1}(\lambda)=\frac{1-R_k(\lambda)}{\lambda}.
\]
Since $R_k(0)=1$, the second expression is a polynomial of degree at
most $k-1$. For $\lambda\in[\alpha,\beta]$,
\[
 |R_k(\lambda)|\le
 2\left(\frac{\sqrt{\bar\kappa}-1}{\sqrt{\bar\kappa}+1}\right)^k
 \le2e^{-2k/\sqrt{\bar\kappa}}.
\]
Take $k=\max\{1,\lceil\tfrac12\sqrt{\bar\kappa}\log(200C)\rceil\}$.
This ensures~\eqref{new:polycontract} with $\nu=1/(100C)$.

For completeness, put $\theta=(\beta+\alpha)/2$,
$h=(\beta-\alpha)/2$, $\tau_0=h/\theta$, $y_0=0$, and $y_1=b/\theta$.
For $j=1,\ldots,k-1$, compute
\begin{align*}
 \tau_j&=(2\theta/h-\tau_{j-1})^{-1},\\
 y_{j+1}&=y_j+\frac{2\tau_j}{h}(b-Ay_j)
                 +\tau_j\tau_{j-1}(y_j-y_{j-1}).
\end{align*}
The Chebyshev recurrence gives $y_j=p_{j-1}(A)b$. It uses one new
matrix application per iteration and a fixed number of vectors. No
eigenproblem or solve is called inside an application. If $\beta=\alpha$,
the exact answer is $b/\alpha$. If $k\ge\dim H$, use exact-arithmetic
CG to termination instead; this returns the exact inverse sample within
$\dim H$ iterations and gives the displayed cap.

The choice between these routines is made before the random signs.
Theorem~\ref{new:polysampler} preserves conditional acceptance and all
pathwise walk bounds. Sum the per-trial application costs using
\eqref{ifc:eq:kappa-base}. The square roots in the right-hand side cost
$O(s_t)$ and are included. This proves the total bound.
\end{proof}

Constant-accuracy sampling removes the high-accuracy solve requirement,
but the polynomial degree still depends on conditioning.
For a guarantee uniform over the entire interval $[\alpha,\beta]$,
the Chebyshev residual above is the minimax polynomial subject to
$R_k(0)=1$. Its degree has square-root dependence on the enclosure ratio.
This is a statement about uniform polynomial approximation, not a lower
bound for the correlated matrices encountered by the signing walk.

\section{Cubic implementations with constructed interfaces}\label{new:structured}

We next construct and maintain solver representations under three
structural hypotheses. The fixed-profile and update-stream
results specify solver classes; the support-separator result applies
directly to the full signing algorithm. Schur identities and hereditary
rank facts used here are proved in Appendices~\ref{new:schur} and~\ref{ifc:sec:tree}.

\subsection{Nonmonotone coefficient motion with arbitrary fixed overlaps}
\label{ifc:sec:nonmonotone-fixed-overlap}

The fixed-profile overlap model of Appendix~\ref{new:schur}
also admits nonmonotone weights when the coefficients obey the actual row-barrier increment law.
The proof extends the lazy positive-Gram argument of Appendix~\ref{new:lazy}.
Its fixed-profile hypothesis is essential to the stated result.

\begin{proposition}[All principal Schur solves for fixed profiles]
During an epoch let
\[
 L_t=H+\sum_{i=1}^{p}\alpha_{i,t}w_iw_i^T,\qquad
 H\succeq0,\quad w_i\in\mathbb R^d,
 \quad \lambda I\preceq L_t\preceq MI,
\]
where $H,w_i,\lambda,M$ are fixed, $p=O(m)$, and $d=O(s)$.
The coefficients may increase or decrease. Suppose each accepted move
has a conditionally mean-zero scalar $\delta_t$ and satisfies
\[
 Y_{i,t+1}-Y_{i,t}=a_{i,t}\delta_t-b_{i,t}\delta_t^2,
 \quad Y_{i,t}=\log\alpha_{i,t},\quad
 |a_{i,t}|\le K_t,\quad 0\le b_{i,t}\le K_t^2,
 \quad K_t|\delta_t|\le 2^{-20}.
\]
Assume the total accepted-move count is deterministically bounded,
$\sum_t K_t^2\delta_t^2\le V$ pathwise, and the logarithms of $M/\lambda$
and the required numerical tolerances are polylogarithmic.
One can initialize in $\widetilde O(pd^2+d^3)$ arithmetic operations and
answer $O(T)$ certified solve queries for $L_t$ or any of its principal
Schur complements at requested states with expected additional cost
\[
 \widetilde O\bigl(pd^2V+T(ps+pd+d^2)\bigr).
\]
$T$ counts accepted states and requested solves together.
The eliminated coordinate set may change at every solve. No overlap-rank,
spectral-gap, or monotonicity assumption is required. Here coefficient
evaluation and the computation of their row-barrier values have the
usual $\widetilde O(ps)$ cost per state; otherwise their actual cost must
be added.
\end{proposition}

\begin{proof}
Set $r=1/8$. For every coefficient maintain an anchor $Y_i^a$ and
refresh it when $|Y_i-Y_i^a|\ge r$. Write $D_i=Y_i-Y_i^a$ before a
move, so $|D_i|<r$. Conditional on the current history and direction,
\[
 \mathbb E\Delta Y_i=-b_i\mathbb E\delta^2,\qquad
 \mathbb E(\Delta Y_i)^2
 \le2a_i^2\mathbb E\delta^2+2(2^{-20})^2b_i\mathbb E\delta^2.
\]
Consequently the expected increase of $D_i^2$ before refreshing is at
most $(2a_i^2+3rb_i)\mathbb E\delta^2$. A refresh discards at least
$r^2$ from this nonnegative potential. Summing to the bounded stopping
time, and discarding residual potentials at epoch ends, gives
\[
 \mathbb E N_i\le \frac{2}{r^2}\mathbb E\sum_ta_{i,t}^2\delta_t^2
       +\frac{3}{r}\mathbb E\sum_tb_{i,t}\delta_t^2
 \le (2/r^2+3/r)V.
\]
Thus the expected total number of noninitial coefficient refreshes is
$O(pV)$. This proof uses no sign or support property of $w_i$.

Let $L^a=H+\sum_i\alpha_i^a w_iw_i^T$. After each refresh sweep,
\[
 e^{-r}L^a\preceq L_t\preceq e^rL^a.
\]
The same sandwich with $r$ replaced by $1/4$ holds throughout a sweep:
every old or new anchor differs from the current coefficient in log
scale by at most $r+2\cdot2^{-20}<1/4$. Thus sequential updates do not
pass through a singular intermediate matrix.

For a finite implementation put
$W_*=1+\sum_i\|w_i\|^2$ and choose $\xi\le\lambda/32$.
Evaluate each anchor coefficient to absolute error at most $\xi/W_*$
and clip it to be nonnegative. The numerical matrix $B$ then satisfies
$\|B-L^a\|\le\xi$, also during a refresh sweep, so $B\succ0$ and
$\tfrac12 B\preceq L_t\preceq2B$.
The coefficient routine truncates a sufficiently negative exponential
to zero and evaluates the remaining bounded argument by finite Taylor
arithmetic. Its cost is polylogarithmic in the reciprocal error and in
a known coefficient upper bound; no exponential or logarithm is an
arithmetic primitive. The refresh comparison itself uses the available
barrier difference, so the logarithm in $Y_i$ is notation only. If the
profile norms or coefficient upper bounds are not polynomially bounded,
their logarithms belong explicitly in the scalar evaluation cost.

Construct $B$ and $B^{-1}$ once. At a refresh,
$B_+=B+\Delta\widehat\alpha_i w_iw_i^T$, and update
\[
 B_+^{-1}=B^{-1}-
 \frac{\Delta\widehat\alpha_i B^{-1}w_iw_i^TB^{-1}}
 {1+\Delta\widehat\alpha_i w_i^TB^{-1}w_i}.
\]
Positive and negative changes each cost $O(d^2)$. The denominator is
positive because both endpoint matrices are positive definite.
The preceding intermediate sandwich justifies this assertion even
when many anchors are refreshed in one sweep.

Preconditioned Richardson iteration with step size $1/2$ and $B^{-1}$ gives a certified
current solve using $\widetilde O(1)$ products, each costing
$O(pd+d^2)$. A current numerical matrix may be evaluated more accurately
than the anchor matrix: if a relative joint residual $0<\epsilon\le1$ is
required, coefficient errors giving operator error at most
$\lambda\epsilon/16$, together with a numerical residual at most
$\epsilon\|b\|/2$, certify the true joint residual. The same check gives
$\|y\|\le 2\|b\|/\lambda$, so the operator-error contribution is at most
$\epsilon\|b\|/8$. The true residual is therefore at most
$5\epsilon\|b\|/8$.
All additional accuracy costs are logarithmic.

Finally partition the current $L$ as
$\left(\begin{smallmatrix}A&C\\C^T&D\end{smallmatrix}\right)$ and
let $S=D-C^TA^{-1}C$. Solve the joint system with right-hand side
$(0,b)$. If its residual is $(r_1,r_2)$, its retained solution component
$y$ obeys
\[
 b-Sy=r_2-C^TA^{-1}r_1,\qquad
 \|b-Sy\|\le(1+M/\lambda)\|(r_1,r_2)\|.
\]
Hence a tighter joint residual certifies the requested Schur residual.
This requires no application of $A^{-1}$ and no construction of $S$.
Selecting a different coordinate partition only changes the zero-padded
right-hand side and the retained output coordinates. The stated cost
includes initialization, all refreshes, coefficient scans, products,
and residual checks.
\end{proof}

In the actual metadata schedule, $V=\widetilde O(1)$,
$T=\widetilde O(n)$, $\sum_e s_e^2=O(n^2)$, and
$\sum_e s_e^3=O(n^3)$. The proposition therefore gives
$\widetilde O(mn^2+n^3)$ total arithmetic for this entire fixed-profile
family. It applies in particular to all coordinate Schur complements
of the actual positive Gram matrix $M+\lambda I$, and to fixed joint
profiles $w_i=(v_i,u_i)$ with arbitrary overlaps. It does not apply
directly to $G=\alpha I+EPE^T$: in the current algorithm,
$E=f(M/\Lambda)T$ has changing profiles and the row-feasible projector changes.
Representing these changes by fixed positive profiles, at comparable
dimension and with the stated coefficient law, has not been established.

\subsection{Interface construction from a signed update stream}
Given a representation at one state and a stream of low-rank changes,
we can construct the later representations by exact batching. The theorem
below is a structured linear-algebra result; its update and rank
hypotheses have not been established for the moving filtered factor in
the Koml\'os walk.

Fix a balanced binary tree on $n$ coordinates with leaf sizes between
$b/2$ and $b$, where $b\asymp r\ge1$.  Put
$L=1+\lceil\log_2(n/b)\rceil$.  A sibling rank is the rank of the
off-diagonal block joining the two children of an internal node.

\begin{theorem}[Batched construction and exact maintenance]
\label{ifc:dyn:thm}
Let $K_0\succ0$ be supplied as a dense matrix, and let
\[
 K_j=K_{j-1}+\sigma_j u_ju_j^T\succ0,
 \qquad \sigma_j\in\{-1,1\},\qquad j=1,\ldots,R,
\]
be a stream of supplied update vectors.  Suppose all sibling ranks of
$K_0,K_r,K_{2r},\ldots$ are at most $r$.  There are $T$ requested solves,
each at an arbitrary current state of this stream.  After charging the
initial interface discovery from the dense matrix, all later interface
construction, all factorizations, and all solves use
\begin{equation}\label{ifc:dyn:cost}
 O\!\left(n^2r+nr^2L^2+RnrL^2+TnrL\right)
\end{equation}
exact real arithmetic operations.  No condition-number bound or
monotonicity assumption is needed.  The initial dense matrix and the
update vectors are the numerical input; evaluating these input objects
is an additional cost if they are given by formulas or oracles.
In particular, if a matrix--vector oracle of cost $C_{\rm mv}$ is the
only access to $K_0$, applying it to the $n$ coordinate vectors adds
$O(nC_{\rm mv})$ for dense extraction.  For the actual dual Gram this
is $\widetilde O(mn^2+n^3)$ at an epoch of size $n$; summing the
decreasing epoch sizes preserves that bound.

Consequently, $r=\widetilde O(\sqrt n)$,
$R=\widetilde O(n^{3/2})$, and $T=\widetilde O(n)$ give
$\widetilde O(n^3)$ arithmetic work.  The theorem remains valid when
the stated rank bounds are supplied as a promise checked at each
recompression.  A failed check invalidates the complexity bound; all nonzero directions
are retained.
More generally the update term fits the cubic target when
$R=\widetilde O(n^2/r)$, provided the other displayed terms fit it.
\end{theorem}
\begin{proof}
First discover the initial factors.  Exact pivoted elimination factors
an $a\times c$ matrix of rank at most $r$ as $UV^T$ in $O(acr)$
operations; this includes searching the remaining entries for nonzero
pivots.  At a tree level the sum of $ac$ is $O(n^2/2^\ell)$.
Thus all initial cross factors cost $O(n^2r)$, and dense leaf extraction
costs $O(nb)$.  Factor the resulting matrix by the recursive solver of
Proposition~\ref{ifc:prop:treecost}, at cost $O(nr^2L^2)$; its inverse
application cost is $Q=O(nrL)$.

Keep this factored matrix $B$ fixed for a batch of at most $r$ updates.
After $k$ updates write
\[
 K=B+W\Sigma W^T,\quad
 W=(u_1,\ldots,u_k),\quad
 \Sigma=\operatorname{diag}(\sigma_1,\ldots,\sigma_k),\quad
 Y=B^{-1}W,\quad H=\Sigma+W^TY.
\]
Since $\Sigma^{-1}=\Sigma$, Woodbury gives
\begin{equation}\label{ifc:dyn:woodbury}
 K^{-1}z=B^{-1}z-YH^{-1}Y^Tz.
\end{equation}
The matrix $H$ may be indefinite.  It is nonsingular by the determinant
lemma, because $B,K,\Sigma$ are nonsingular.  On receiving another
update, compute $y=B^{-1}u$ at cost $Q$, the new cross vector $W^Ty$
at cost $O(nk)$, and the diagonal entry $\sigma+u^Ty$ at cost $O(n)$.
The bordered inverse formula updates $H^{-1}$ in $O(k^2)$ operations.
Its scalar Schur pivot is nonzero: both the old and new $H$ are
nonsingular.  This uses positive definiteness at every stream state,
including intermediate updates.  Each update therefore costs $O(nrL)$,
and each solve in \eqref{ifc:dyn:woodbury} costs $O(nrL)$.

At the end of $r$ updates, construct the next tree representation.
For a sibling block whose old factor is $UV^T$, the new block is
\[
 [\,U\ \ W_a\,]\,[\,V\ \ W_b\Sigma\,]^T.
\]
These factors have at most $2r$ columns.  Exact rank-profile factorizations
of the two tall factors, followed by exact elimination of their at most
$2r$ square core, reduce it to its actual rank.  At a node of order
$N$ this costs $O(Nr^2+r^3)$, without inspecting a fresh dense cross
block.  The assumed rank bound guarantees the reduced width is at most
$r$.  Summing over nodes costs $O(nr^2L)$, since the number of nodes is
$O(n/b)$ and $b\asymp r$.  Apply the batch update to each dense leaf
at total cost $O(nb r)=O(nr^2)$.  The next complete recursive
factorization costs $O(nr^2L^2)$.

There are at most $R/r$ completed batches.  Charging each reconstruction
and factorization to its preceding $r$ updates gives $O(RnrL^2)$.
An incomplete final batch is served by \eqref{ifc:dyn:woodbury} and needs
no reconstruction.  Adding initialization and requested solves proves
\eqref{ifc:dyn:cost}.
\end{proof}

\paragraph{Why the hypotheses matter.}
The rank at a batch boundary is a structural property, not a consequence
of small update norm.  A sequence of $r$ independent cross directions can
increase a sibling rank by $r$, however small their coefficients.
Recompression removes algebraic dependence and cancellations; it does
not remove nonzero directions.  No nested-basis hypothesis is used:
a hierarchical off-diagonal low-rank (HODLR) representation suffices.  The rank promise need
hold only at reconstruction times; during a batch the displayed
cross ranks may reach $2r$.

\paragraph{Connection to the actual dual Gram.}
At fixed $E$, an inserted independent row constraint gives
$P'=P-ww^T$ for a unit feasible vector $w$.  Hence
\[
 (\alpha I+EP'E^T)-(\alpha I+EPE^T)=-(Ew)(Ew)^T.
\]
The reverse row event has the positive sign.  These are valid SPD
updates even at an arbitrarily small positive ridge.  Forming each
update vector costs one application of $E$.  The existing
$\widetilde O(n)$ elementary row-event bound therefore fits the usual
$\widetilde O(mn^2+n^3)$ input-vector budget.  If the tree ranks are
controlled, Theorem~\ref{ifc:dyn:thm} includes construction of all interfaces
caused by these events.  No independent dense spectral decomposition
or constraint factorization is hidden in that claim.

More generally, if $E'=E+UV^T$ with $U,V$ having $q$ columns and $P$
is fixed, then, with $A=EPV$ and $C=V^TPV$,
\begin{equation}\label{ifc:dyn:factorchange}
 E'PE'^T-EPE^T
 = [\,U\ A\,]
   \begin{pmatrix}C&I\\ I&0\end{pmatrix}
   [\,U\ A\,]^T.
\end{equation}
The change has rank at most $2q$.  The small symmetric core admits an
exact signed rank-one decomposition in $O(q^3)$ real arithmetic
operations, with square roots allowed.  Its positive terms may be
applied first and its negative terms second: if the endpoint matrices
are SPD, every intermediate matrix is SPD, because the positive stage
dominates the old matrix and every partial negative stage dominates the
final matrix.  The cost of $q$ applications of $EP$ to construct $A$
must also be charged.  Thus a rank bound for factor changes would help control the cost, but
the update vectors must also be constructed within the budget.
Moreover, the structured ranks of the two original endpoints do not
automatically certify the ranks at reconstruction times inside this
decomposition.  The batch-boundary hypothesis must still be verified.
If each atomic endpoint change has signed width at most $r$, one may
instead keep whole changes together and rebuild after the accumulated
width first reaches $r$.  The width is then less than $2r$, so the same
cost bound holds up to constants and the rank promise is needed only
at original endpoints.  This variant still charges the positive-first
intermediate solves and the supplied update-vector construction.

The present walk has no proved bound
$\sum_t\rank(E_{t+1}-E_t)=\widetilde O(n^{3/2})$, nor a proved
$\widetilde O(\sqrt n)$ sibling-rank bound.  Its finite polynomial
filter is not exactly a low-rank spectral projector; an effective-rank
estimate cannot be substituted for either exact hypothesis above.
Accordingly this theorem does not close the unrestricted cubic bound.


\subsection{Construction from input row supports}
\label{ifc:sec:supportinterfaces}

The tree solver of Proposition~\ref{ifc:prop:treecost} takes its factors
as input. When the row profiles have small support separators, we can
construct those factors within the required budget. The result includes the polynomial filter and the exact
constraint projector. No temporal regularity of the weights, Perron vector,
anchors, or filter coefficients is needed.

Fix a binary coordinate partition tree of height $h=O(\log n)$. Its leaf
sizes are at most $b$. For every non-root node $J$, count the input rows whose
supports meet both $J$ and its full complement $J^c$. Let $c$ be the largest
such count. The full-complement condition includes interactions through
ancestors. The tree is part of the input, or is specified by an ordering as
in the sliding-window example below. The bounds below do not include a search for such a tree.
Zeroing entries, restricting to surviving coordinates, or removing rows
cannot increase these counts.

\begin{lemma}[From supported constraints to an exact interface bound]
\label{ifc:lem:supportprojector}
Suppose a family of linear constraints consists of coordinate equations,
at most one normal supported on each input row support, and $g$ additional
arbitrary normals. Its exact kernel projector $P$ satisfies
\[
 \operatorname{rank}P_{J,J^c}\le c+g
\]
at every tree cut, even when the constraint normals are dependent.
More generally, the bound is the number of current supported normals
crossing the cut, plus $g$.
\end{lemma}
\begin{proof}
First impose the normals supported entirely on one side. Their projector
is $P_0=\operatorname{diag}(P_J,P_{J^c})$. Let $H$ collect the crossing and
additional normals, and let $U=P_0H^T$. Then
\[
 P=P_0-U(U^TU)^\dagger U^T.
\]
The correction has rank at most the number of columns of $U$. The formula
also follows by selecting any independent columns of $U$ and using their
ordinary Gram inverse; thus it requires no pseudoinverse routine.
\end{proof}

In the actual retained-row walk there is at most one current row-supported
normal per input row. Large and medium classes are disjoint. A medium
row has at most one active anchor per sign, and its two signs cannot both
be active: the state identities give
\[
 u_{i,+}+u_{i,-}
 =-150/r_i+7G_i/r_i^2+35\le-58,
\]
whereas two active signs would both belong to $[-1/50,0]$.
Released historical anchors are removed from the constraint list.
Held frozen-coordinate equations are local, and the only additional
normals are $x$ and aggregate quadratic energy. Hence $g=2$.
This conclusion also holds after cleanup, since supports only shrink.

Combine the two signs of each medium input row in the derivative factor:
\[
 M=\sum_i a_i z_i z_i^T+D+\eta uu^T,
 \qquad T=\sum_i z_i t_i^T,
 \qquad
 t_i=\frac27\sum_{\sigma}\alpha_{i\sigma}(z_i^Tv)g_{i\sigma}.
\]
Here $D$ is diagonal, $u$ is the indicator of the current coordinate set,
and $z_i,t_i$ have the same supported row index set. On the current
coordinates $a_i=\sum_\sigma\alpha_{i\sigma}$ and $u=\mathbf1$.
The same identities hold for the fixed finite-coefficient target when its
scalar coefficients are evaluated in these factored formulas.
At every tree cut,
\[
 \operatorname{rank}M_{J,J^c}\le c+1,\qquad
 \max\{\operatorname{rank}T_{J,J^c},
        \operatorname{rank}T_{J^c,J}\}\le c.
\]
For a degree-$d$ polynomial $f$, put $E=f(M/\Lambda)T$ and
$G=\alpha I+EPE^T$. The polynomial closure lemma therefore gives
\begin{equation}\label{ifc:eq:supportR}
 \operatorname{rank}G_{J,J^c}\le R,
 \qquad R=2d(c+1)+3c+2.
\end{equation}
The rank bound $R$ may always be replaced by $\min\{n,R\}$.
The next proposition constructs these interfaces directly from the row
profiles, avoiding $R$ generic dense matrix-vector products.

\begin{proposition}[Fully charged profile construction]
\label{ifc:prop:profileconstruction}
Given the support tree, the current row profiles and scalar coefficients,
and the current logical constraint list, one can construct exact
hierarchical factors for $M,T,P,E,G$ in
\begin{equation}\label{ifc:eq:profileconstruction}
 O\!\left(mb^2+(d+1)n(b^2+R^2)(1+h)^2\right)
\end{equation}
real arithmetic operations, provided $b\ge\max\{1,\min\{R,n\}\}$.
Preparing supported profile vectors and scalar coefficients is charged
separately; a full scan costs $O(mn)$ per state. One-time support-list
preparation costs $O(\operatorname{nnz}(A)(1+h)+m+n)$ comparisons and
index operations. All stated bounds include dependent-row handling,
low-rank recompression, and dense leaf formation. They have no condition
number or ridge-eigenvalue factor.
\end{proposition}
\begin{proof}
We account for each operation needed to construct the hierarchical
representation.

\emph{Thin-factor compression.}
If an off-diagonal block is supplied as $UV^T$ with $k$ columns, exact
rank-profile elimination on $U$ and $V$, followed by elimination of an
at most $k\times k$ coefficient matrix, gives factors of their exact
product rank. The cost is
$O((\operatorname{rows}U+\operatorname{rows}V)k^2+k^3)$.
No dense off-diagonal block is formed and no singular-value decomposition
is needed. This is exact real arithmetic, not a finite-bit stability claim.

\emph{Initial row sums.}
At each internal node, list the rows meeting both children. Their
restrictions give the cross-block factors for $M$ and both orientations
of $T$ directly; the Perron ridge adds one column. Across a level this
costs $O(n(c+1))$, hence $O(n(c+1)(1+h))$ in total.
For a leaf of size $b_J\le b$, a row wholly inside that leaf costs at most
$O(b_J^2)$ to accumulate. Across leaves such rows cost $O(mb^2)$.
Every other row meeting this leaf crosses its external cut, so there are
at most $c$ of them. Their leaf contributions cost at most
$O(cb_J^2)$, whose sum is $O(nbc)$. Thus leaf formation costs $O(mb^2+nbc)$, including all row contributions. Diagonal terms and all off-diagonal factors are constructed
from the supplied profiles. Ordinary exact factor compression, if used,
adds at most $O(n(c+1)^2(1+h))$.

\emph{The kernel projector.}
Assign each supported constraint to the lowest tree node containing its
support. At a leaf, impose its assigned rows by incremental orthogonal
projection. A row costs $O(b^2)$; all leaf work is
$O((m+n)b^2)$, including frozen-coordinate rows. At an internal node,
first use the already constructed child projectors as a block diagonal
projector $P_0$. At most $c$ assigned row normals meet both children.
Form $U=P_0H^T$ by hierarchical multiplication, select independent columns
$U_0$, and make the exact correction
\[
 P_0-U_0(U_0^TU_0)^{-1}U_0^T.
\]
Propagate this low-rank correction to the stored leaves and sibling factors,
recompressing the thin factors. Finally impose the two global normals in
the same way. After each complete constraint group, the logical projector is the exact kernel projector
of a subset of the original supported constraint family. Consequently its
external ranks are at most $c+2$ by the preceding lemma. Thus no unproved
rank cap is used by the recompressions.

For a node containing $N$ coordinates, applying a hierarchical matrix of
rank at most $r$ to $k$ columns costs
$O(Nbk+Nrk(1+h))$. An update of rank $k=O(r)$, including leaf updates and
exact thin-factor compression, costs
$O(Nbr+Nr^2(1+h))$. Selecting the independent columns and forming/inverting
their Gram matrix costs $O(Nc^2+c^3)$, with every small core capped by the actual selected rank, at most $N$. Summing over tree levels gives
$O((m+n)b^2+n(b^2+(c+2)^2)(1+h)^2)$ for $P$.
Dependencies cause zero projected columns and are handled by the exact
rank tests.

\emph{Products and the filter.}
For two represented matrices $A,B$, at a node the cross block is
\[
 (AB)_{12}=A_{11}U^B_{12}(V^B_{12})^T
        +U^A_{12}(B_{22}^TV^A_{12})^T.
\]
Construct these thin factors with child matrix applications. Recursively
multiply the two diagonal child pairs, and add their missing low-rank
corrections $A_{12}B_{21}$ and $A_{21}B_{12}$. Recompress each resulting
thin factor exactly. If the external ranks of $A$ and $B$ are $r_A,r_B$,
all products of principal submatrices encountered this way have external
rank at most $r_A+r_B$. A diagonal block after a correction is a restriction
of the product at its parent; after the next ancestor correction it is a
restriction of the product at that ancestor. The same rank bound applies
at each intermediate stage. This accounts for contributions from every ancestor as well as the
current sibling block.

The work at a node is bounded by
$O(Nb(r_A+r_B)+N(r_A+r_B)^2(1+h))$, beyond child products.
Dense leaf products total $O(nb^2)$. Summing gives
\[
 O\!\left(n\bigl(b^2+(r_A+r_B)^2\bigr)(1+h)^2\right)
\]
for one multiplication. Addition of a diagonal or transpose is cheaper.
Horner evaluation of $f(M/\Lambda)$ uses $d$ products. Its intermediate
external ranks are at most $k(c+1)$ after degree $k$. Forming $E$, then
$EP$, then $EPE^T$ gives precisely the rank budget in
\eqref{ifc:eq:supportR}. There are $d+3$ products, and all factor-construction
costs are included in \eqref{ifc:eq:profileconstruction}.
\end{proof}

\begin{corollary}[Cubic signing implementation on a support-separator class]
\label{ifc:cor:profilecubic}
Suppose the initial Koml\'os input is supplied with a balanced support tree
whose cut width is $c$, where $(d+1)(c+1)=\widetilde O(\sqrt n)$ for the
largest filter degree used by the walk. Then the modular signing algorithm
can use the constructed tree solver in expected
\[
                 \widetilde O(mn^2+n^3)
\]
real arithmetic operations. Its discrepancy guarantee is unchanged. This conclusion requires neither the signed-production conjecture
nor a monotonicity hypothesis for the changing weights or constraints.
\end{corollary}
\begin{proof}
Use the original ambient dimension $n$ throughout the cost analysis.
Choose the fixed tree, or coarsen it, to have leaf size
$b=\min\{n,\max\{\lceil\sqrt n\rceil,\lceil R\rceil\}\}
=\widetilde O(\sqrt n)$; coarsening does not increase any retained cut count.
Alternatively the factors may be kept in the initial tree with the same
polylogarithmic accounting. After cleanup, pad the current operators with
zero rows and columns on deleted coordinates, pad $P$ by zero, and add the
usual independent $\alpha I$ blocks to $G$. Thus the fixed tree remains
valid without rebalance or discovery. The ridge vector is simply the current
coordinate indicator, so it still contributes only one cross-rank column.

The construction and direct factorization cost per direction-selection state
is $\widetilde O(mn+n^2)$, including the profile scan. The modular walk uses
$\widetilde O(n)$ such states/trials in the already established expected
accounting. Its other work is
$\widetilde O(mn^2+n^3)$. Exact dual solves meet the existing residual
contract; applying this dual inverse as the preconditioner for the original
sampler requires only the already allowed polylogarithmic iteration factor.
The factorization may be rebuilt at every state within the total budget,
so this bound requires no persistence of spectral information between states. The sampler and discrepancy proof are unchanged.
\end{proof}

Sliding-window rows in a known coordinate order give a connected example
of the support-separator hypothesis. If every row is supported
on an interval of length at most $w$ and at most $\nu$ rows begin at each
coordinate, every interval cut is crossed by at most $2\nu(w-1)$ rows.
A balanced interval tree is constructed directly. Hence
$\nu w=\widetilde O(\sqrt n)$ suffices. Rows with exactly $w\ge100$ entries
$\pm1/\sqrt w$, one row at each eligible starting point, have column norms
at most one, row squared norm one, and pass the medium-row cleanup test.
They overlap along a connected chain. Taking $w=\lfloor\sqrt n\rfloor$
gives a connected class with growing row supports.
The theorem also permits arbitrary additional rows supported inside leaves,
provided the stated crossing count is preserved.


\subsection{Comparison with general inverse maintenance}

The fixed-profile and support-separator mechanisms above solve different
parts of the problem. The first tolerates arbitrary overlaps and uses the
actual coefficient increment law, but its profile vectors must be fixed.
The second constructs the actual filtered Gram afresh, with no temporal
regularity hypothesis, but requires a supplied support tree of small width.
The signed-update theorem constructs successive interfaces from numerical
update vectors and a verified rank bound at reconstruction times. These
are proved algorithmic statements under explicit structural hypotheses,
not additional conjectures about the unrestricted walk.

\paragraph{Inverse maintenance with a fixed design.}
Lee--Sidford~\cite[Definition 4 and Theorem 13]{LSmaint} obtain
$\widetilde O(d^\omega+T(\operatorname{nnz}(A)+d^2))$ for a fixed
design and positive diagonal weights, under leverage-weighted and
coordinatewise log stability and polynomial Loewner comparability.
Their Section 5 handles adaptive output dependence. The favorable actual
barrier law already supplies small increments for a fixed design:
leverage scores sum to $d$, so
$\|\Delta Y\|_\sigma^2\le4dK^2\tau^2$.
The existing schedule makes this small when $d=O(s)$.
The direct lazy proof above additionally exploits the mean-zero coin and
uses classical arithmetic. Neither argument covers moving profile vectors
by merely relabeling them as weights.

\paragraph{Finding hierarchical interfaces costs matrix applications.}
The accounting in Levitt--Martinsson~\cite[Section 4.4]{levitt}
is $O(rC_{\rm mv}+nr^2)$ for hierarchically semiseparable (HSS) rank $r$, where $C_{\rm mv}$ is the
cost of one matrix application. Their Remark 6 leaves a higher-level
sketch-independence issue unresolved. Amsel et al.~\cite[Theorem 13]{hss2026}
give a proved construction with fresh sketches at each level, using
$O(r\log(n/r))$ applications and $O(nr^2\log(n/r))$ additional arithmetic.
Its guarantee is in expected squared Frobenius norm; exact recovery follows
when the best HSS approximation has zero error. For HODLR approximation,
Chen et al.~\cite[Theorem 1.2]{hodlr2025} also prove an
$\Omega(r\log(n/r))$ query requirement for exact recovery in the
general black-box model. These results do not prove a lower bound for the
correlated Koml\'os trajectory or for algorithms using its explicit profiles.

In our unstructured application $C_{\rm mv}=\widetilde O(mn+n^2)$.
Rebuilding independently at $T=\widetilde O(n)$ states therefore costs
\[
 \widetilde O\bigl(Tr(mn+n^2)+Tnr^2\bigr).
\]
For $r\asymp\sqrt n$ the leading term is
$\widetilde O(mn^{5/2}+n^{7/2})$.
Thus a cheap solve with a given interface does not by itself cover the
cost of repeatedly constructing that interface. The signed-update construction and the
support-profile construction explicitly avoid those repeated probes.
Furthermore a Frobenius approximation does not by itself give relative
Loewner control at a small ridge; the normalized deflation argument needs
the stated principal normalization.

Section~\ref{ifc:sec:wideinterfaces} gives a further check specific to the
actual canonical formulas: the local retained-row invariants permit
linearly many strong external couplings in every coordinate tree.
The example has a cheap global incidence decomposition. It therefore
rules out a conclusion from those local invariants alone, while giving
no lower bound on the signing algorithm and no reachability claim.

\paragraph{The precise remaining boundary.}
Within fixed metadata let $Z$ have the profile columns $z_i$, let the
rows of $R$ be $(\sigma b_i/r_i)^T$, put $X=\diag(x)$, and let
$a_i=(2/7)\alpha_i(z_i^Tv)$. With row-sign indexing, the actual formulas are
\[
 M=M_0+Z\diag(\alpha_i)Z^T,\qquad
 T=Z\diag(a_i)(R-7Z^TX),\qquad E=f(M/\Lambda)T.
\]
Within this metadata epoch, only $Z,R,M_0$ are fixed. The proved refresh law for $\alpha_i$ does not
automatically control $a_i,X,f(M/\Lambda)$, or the changing $P$ in
$G=\alpha I+EPE^T$. A direct symmetric circuit linearization contains
blocks $\left(\begin{smallmatrix}0&C\\C^T&0\end{smallmatrix}\right)$,
which are indefinite when $C\ne0$; it is outside a positive-weight
fixed-design theorem. This is a check on that particular lift, not an
impossibility claim about every lift.

For general inputs there is still no proved, cheaply constructible
small-interface representation of $G$, and no proved update-vector budget
that meets Theorem~\ref{ifc:dyn:thm}. The established high-factor rank bound
is $(57/640+5\eta)s$, which is linear in $s$; $s/600$ is a curvature
dimension surplus and is not an upper bound on that rank. Consequently
the unrestricted expected cubic bound with no conditioning factor remains
conditional on the production hypothesis~\eqref{d:conjecture}. The complete
condition-number-sensitive bound in Theorem~\ref{ifc:thm:kappa} does not
use that hypothesis. These implementation results retain the discrepancy constant $D_*$.

\section{Two-sided inverse maintenance}\label{c:section}
This section proves the algebraic and numerical primitives used by the
controller in Section~\ref{d:section}. All statements here are
unconditional. The controller rebuilds at metadata cleanup and leaves
unused caches untouched in the low branch; its one remaining complexity
conjecture is stated in \eqref{d:conjecture}.

\subsection{Fixed-metadata solver parameters}\label{c:padding}
Let $s$ be the fixed retained dimension of a metadata epoch. Use the local
parameters of Section~\ref{a:direction}; in particular
$\eta=1/(4096JN)$ is fixed by its halving phase,
$\mu=\eta+4096/s$, and $\chi=\mu/\eta$.
The numbers $K,\tau,e_*,\rho_{\mathrm{ridge}}$ are fixed in this epoch.
Every arithmetic factor $E$ is frozen throughout
sampling retries, has its exact arithmetic transpose, and obeys
$\|E\|\le3$. Applications cost $\widetilde O(ms+s^2)$.
The actual system acts on $\mathcal H=\operatorname{ran}P$ and is
\begin{equation}\label{c:target}
 \mathsf A=P(F+\rho_{\mathrm{ridge}}E^TE)P\big|_{\mathcal H},
 \qquad \eta I\preceq F\preceq\mu I.
\end{equation}
In the low branch use $F=\eta I,E=0$ and solve directly. The remaining
formulas describe the high-branch solver. The arithmetic model permits
exact square roots and comparisons; log determinants are proof
potentials and are never evaluated by the signing algorithm.

\subsection{Two complementary Gram systems}\label{c:grams}
Write $\mathcal C$ for the primal matrix target, retaining $C=1200$
for the sampler constant. All identities refer to the actual retained solver space of dimension $s$.
Set $\alpha=\eta/\rho_{\mathrm{ridge}}$ and define
\begin{equation}\label{c:eq:two}
 \mathcal C=\eta I+\rho_{\mathrm{ridge}} E^\top E,\qquad
 G=\alpha I+EPE^\top,\qquad
 \mathcal C_H=P\mathcal C P\big|_{\mathcal H}.
\end{equation}
Both $\mathcal C$ and $G$ have a uniform condition bound
$\kappa_*=1+9\rho_{\mathrm{ridge}}/\eta=\widetilde O(s^3)$.
The primal target $\mathcal C$ is an \emph{ambient}, uncompressed matrix.

\begin{lemma}[Scalar-ridge comparison]\label{c:lem:scalar}
On $\mathcal H$, $\mathcal C_H\preceq \mathsf A\preceq\chi \mathcal C_H$.
\end{lemma}
\begin{proof}
The lower bound is $F\succeq\eta I$. For the upper bound use
$F\preceq\mu I=\chi\eta I$ and $\rho_{\mathrm{ridge}}\le\chi\rho_{\mathrm{ridge}}$.
\end{proof}
Thus replacing $F$ by the scalar ridge in the preconditioner incurs only
the polylogarithmic condition factor $\chi$.

\begin{lemma}[Dual formula with a residual certificate]\label{c:lem:dual}
For $b\in\mathcal H$,
\[
 \mathcal C_H^{-1}b=\eta^{-1}\{b-PE^\top G^{-1}Eb\}.
\]
If $r_d=Eb-Gz$ and $y=\eta^{-1}(b-PE^\top z)$, then
\begin{equation}\label{c:eq:dualres}
 \mathcal C_Hy-b=(\rho_{\mathrm{ridge}}/\eta)PE^\top r_d.
\end{equation}
\end{lemma}
\begin{proof}
Expand $\mathcal C_Hy$ and use $\alpha=\eta/\rho_{\mathrm{ridge}}$. The two terms involving
$EPE^\top z$ and $\alpha z$ combine to $Gz$. Setting $r_d=0$ gives the
inverse formula.
\end{proof}
Because the dual formula involves subtraction,
equation~\eqref{c:eq:dualres} requires a sufficiently small
\emph{absolute} dual residual; constant relative accuracy alone does not suffice. The needed accuracy is inverse polynomial,
so a constant-condition dual solver still uses only logarithmically many
iterations. Section~\ref{c:sec:solve} gives a fully residual-based interface.

We maintain SPD matrices $B^{\rm p}$ and $B^{\rm d}$ approximating $\mathcal C$ and
$G$, respectively. If either is a constant-factor preconditioner, the
original system can be solved in $\widetilde O(ms+s^2)$ work. The other cache can remain unchanged at that state. This choice matters
because input
rotations preserve the dual Gram, whereas output rotations preserve the
primal Gram.

\subsection{A constant-divergence secant correction}
The S-divergence \cite{Sra} is
\begin{equation}\label{c:eq:S}
 \mathcal{S}(B,K)=\log\det\frac{B+K}{2}
       -\frac12\log\det B-\frac12\log\det K\ge0.
\end{equation}
It is congruence invariant. Equivalently, it is the sum of
$\log\cosh(\tfrac12\log\lambda)$ over the generalized eigenvalues.
The logarithmic dependence on extreme eigenvalue ratios is useful here.
The controller never needs to evaluate this potential.

For a target $K\succ0$ and a nonzero vector $u$, use the standard Hessian
BFGS update
\begin{equation}\label{c:eq:bfgs}
 B^+=B-\frac{Bu\,u^\top B}{u^\top Bu}
          +\frac{Ku\,u^\top K}{u^\top Ku}.
\end{equation}
BFGS-based matrix-inverse approximation is established
\cite{Gower}. We prove the decrement below directly for a fixed target;
it requires no convergence theorem for changing targets.

\begin{theorem}[Exact S-divergence decrement]\label{c:thm:drop}
Let $r=(u^\top Bu)/(u^\top Ku)$. Then $B^+\succ0$, $B^+u=Ku$, and
\begin{equation}\label{c:eq:drop}
 \mathcal{S}(B,K)-\mathcal{S}(B^+,K)
 \ge \log\frac{1+r}{2\sqrt r}.
\end{equation}
In particular, a verified witness
$q=(u^\top Ku)/(u^\top Bu)\notin[1/2,3/2]$
gives a decrease at least
$\gamma=\log(5/(2\sqrt6))>0.0204$.
\end{theorem}
\begin{proof}
Whiten $K$ by congruence and rotate the transformed, normalized $u$ to
$e_1$. Write the transformed cache as
\[
 B=\begin{pmatrix}r&w^\top\\ w&D+ww^\top/r\end{pmatrix},\qquad D\succ0.
\]
Equation~\eqref{c:eq:bfgs} gives $B^+=\diag(1,D)$. The Schur determinant
formula and the rank-one determinant lemma yield the exact identity
\begin{align*}
 \mathcal{S}(B,I)-\mathcal{S}(B^+,I)
 &=\log\frac{1+r}{2\sqrt r}\\
 &\quad+\log\left(1+
       \frac{w^\top(D+I)^{-1}w}{r(1+r)}\right).
\end{align*}
The second term is nonnegative. The first term increases with
$|\log r|$. The smaller boundary value among $r=2$ and $r=2/3$
is the stated $\gamma$.
\end{proof}

Each correction requires one application of $K$ and dense rank-two updates. Store a Cholesky factor of each cache. Apply the
positive rank-one update in~\eqref{c:eq:bfgs} before its negative downdate;
the intermediate matrix dominates the final SPD matrix. Standard rank-one
Cholesky updates therefore cost $O(s^2)$ arithmetic operations.

\subsubsection{Finding a witness without a spectral decomposition}\label{c:tester}
Given $B=LL^\top$, apply
$H=L^{-1}KL^{-\top}-I$ using matrix-vector products and triangular solves.
For a chosen failure parameter $0<\delta_{\rm test}<1$, the following
finite test returns a verified witness $q\notin[1/2,3/2]$, or declares the
cache good. Except with probability $\delta_{\rm test}$, a good declaration
satisfies
\begin{equation}\label{c:eq:good}
 \tfrac14 B\preceq K\preceq\tfrac74 B.
\end{equation}
The probability statement is conditional on all information before the
fresh test randomness; every returned witness is verified exactly.

Choose an integer $k$ with $(11/12)^{2k}\le1/(16s)$.
For an independent Rademacher vector $\xi$, form and normalize $H^k\xi$,
skipping a zero result. Denote the resulting unit vector by $z$. Compute the exact
one- or two-dimensional compression of $H$ to
$\operatorname{span}\{z,Hz\}$, using square roots and scalar arithmetic.
Map its extremal Rayleigh vectors $w$ to $u=L^{-\top}w$ and directly test
$(u^\top Ku)/(u^\top Bu)$. Return any verified witness, with least-index
and positive-sign tie rules. Repeat independent starts a number $j$ such
that $(11/12)^j\le\delta_{\rm test}$. If no witness is found, declare good.
The integer choices require only rational multiplication and comparisons.

To prove the test, suppose $\|H\|>3/4$ and let $v$ be a unit eigenvector
of extremal absolute eigenvalue. For $X=(v^\top\xi)^2$, one has
$\mathbb EX=1$ and $\mathbb EX^2\le3$. Cauchy--Schwarz applied on
$\{X\ge1/2\}$ gives the Paley--Zygmund bound
$\Pr\{X\ge1/2\}\ge1/12$. On this event the normalized start has squared
overlap at least $1/(2s)$ with $v$. After $k$ powers the relative squared
mass on eigenvalues of absolute value below $11/16$ is at most
$2s(11/12)^{2k}\le1/8$. Thus
\[
 \|Hz\|\ge(11/16)\sqrt{7/8}>5/8.
\]
The compression has norm at least $\|Hz\|$, because its action on $z$
is exactly $Hz$. It therefore has a Rayleigh vector with eigenvalue of
absolute value greater than $5/8$, which gives a verified quotient outside
$[1/2,3/2]$. Each independent start detects a bad cache with probability
at least $1/12$. Failure to find such a witness after the prescribed starts
has probability at most $\delta_{\rm test}$. The complement of
\eqref{c:eq:good} has $\|H\|>3/4$, proving the claim.

The test uses $O(\log(2s)\log(1/\delta_{\rm test}))$ products with $H$.
Each product costs one target application and $O(s^2)$ for triangular
solves. This count has no spectral condition-number dependence. We use
exact real arithmetic with square roots for these scalar normalizations
and two-dimensional calculations. If scalar approximations are used
instead, the strict $5/8$ witness margin must be retained and the final
quotient still verified; no approximate product is allowed to redefine the
fixed arithmetic target between calls.

A good declaration is probabilistic. Every final linear-system residual
is nevertheless checked against the fixed arithmetic system. A capped
solve which fails that check triggers a dense solve and a reset of
\emph{both} caches to the exact current targets. Section~\ref{c:failure}
charges the small expected fallback cost without assuming that a good
declaration is infallible.

\subsection{Changing row constraints are affordable}
\begin{lemma}[Rank and monotonicity charge]\label{c:lem:drift}
For every cache $B\succ0$ and targets $K,K'\succ0$, let
$\lambda_i$ be the eigenvalues of $K^{-1/2}K'K^{-1/2}$. Then
\begin{equation}\label{c:eq:variation}
 |\mathcal{S}(B,K')-\mathcal{S}(B,K)|
 \le\frac12\sum_i|\log\lambda_i|.
\end{equation}
If both targets lie in $[aI,bI]$ and $\rank(K'-K)\le r$, this is at most
$\tfrac r2\log(b/a)$. If $K_t$ is Loewner monotone, then even for
arbitrarily changing caches $B_t$,
\begin{equation}\label{c:eq:monotone}
 \sum_t\{\mathcal{S}(B_t,K_{t+1})-\mathcal{S}(B_t,K_t)\}
 \le\tfrac12|\log\det K_T-\log\det K_0|.
\end{equation}
\end{lemma}
\begin{proof}
The derivative with respect to the target in direction $Z$ is
\[
 \tr\left[\left((B+K)^{-1}-\tfrac12K^{-1}\right)Z\right].
\]
For $Z\succeq0$, its absolute value is at most
$\tfrac12\tr(K^{-1}Z)$, since $0\preceq(B+K)^{-1}\preceq K^{-1}$.
By congruence, vary the generalized eigenvalues one at a time. Integrating
these rank-one monotone paths proves~\eqref{c:eq:variation}. A rank-$r$
perturbation has at most $r$ generalized eigenvalues different from one,
all between $a/b$ and $b/a$. Integrate along each monotone target segment
and telescope the log determinants to obtain~\eqref{c:eq:monotone}.
\end{proof}

At fixed $E$, a rank-$r$ change in $P$ changes $G$ by rank at most $r$;
it does not change $\mathcal C$. Thus the inherited $\widetilde O(s)$ elementary row events
produce at most $\widetilde O(s)$ total divergence. This includes insertions,
deletions, replacement of moving rows, and changes in dependency. It does
not require a relative stability estimate for the fine ridge.

\subsubsection{Maintaining a primal preconditioner on the current kernel}
Store the zero-extended constrained inverse
$Z=P(PB^{\rm p}P|_{\mathcal H})^{-1}P$. A rank-two change
$B^{\rm p}\mapsto B^{\rm p}+URU^\top$ gives
\[
 Z^+=Z-ZU(R^{-1}+U^\top ZU)^{-1}U^\top Z.
\]
For insertion of an independent constraint $a^\top h=0$,
$Z^+=Z-Za\,a^\top Z/(a^\top Za)$.
For a one-dimensional enlargement of $\mathcal H$, take a unit vector
$u\perp\mathcal H$ in the enlarged space and set
\[
 w=u-ZB^{\rm p}u,\quad
 d=u^\top B^{\rm p}u-u^\top B^{\rm p}ZB^{\rm p}u>0,
 \qquad Z^+=Z+ww^\top/d.
\]
The rank-two formula is Woodbury on an orthonormal basis of $\mathcal H$;
the indicated two-by-two inverse exists because the updated compressed
matrix is positive definite. The insertion formula is checked by applying
the displayed inverse to a right-hand side and verifying both the new
constraint and the normal equations. For enlargement, use the block basis
consisting of an old orthonormal basis and $u$. Its last Schur complement
is exactly $d>0$, and block inversion gives the stated $ww^\top/d$ term.
These identities therefore update the exact constrained inverse.

All costs are $O(s^2)$. The independent row basis and dependent-row
promotion rules of Section~\ref{m:sec:rows} identify these elementary events without
assuming well-conditioned pivots. Maintain $Z$ on row events even while
the primal cache is idle. This is already within the row-event budget.

\subsection{Residual-certified primal and dual solves}\label{c:sec:solve}
The cache controller and its initialization and failure schedule are
specified in Section~\ref{d:controller}. Once a cache is declared good,
the following explicit iterations enforce the true arithmetic residual.

\paragraph{Primal solve.}
When the primal declaration is correct, compression of
\eqref{c:eq:good} and Lemma~\ref{c:lem:scalar} give
\[
 \tfrac14 PB^{\rm p}P|_{\mathcal H}\preceq\mathsf A
 \preceq\tfrac74\chi PB^{\rm p}P|_{\mathcal H}.
\]
Starting from zero, the explicit iteration
\[
 y^+=y+\frac4{7\chi}\,Z(\widetilde b-\mathsf Ay)
\]
contracts error in $\mathsf A$-norm by at most
$1-1/(7\chi)$. Indeed the eigenvalues of the preconditioned
operator lie in $[1/4,7\chi/4]$. The iteration therefore has an explicit contraction bound.

\paragraph{Dual solve.}
Put $L_C=\eta+9\rho_{\mathrm{ridge}}$. For an outer residual
$r=\widetilde b-\mathsf Ay$, construct $z\in\mathcal H$ satisfying
\begin{equation}\label{c:eq:inner}
 \|\mathcal C_Hz-r\|\le\delta_0\sqrt{\eta/L_C}\,\|r\|,
 \qquad\delta_0=1/(4\chi).
\end{equation}
By Lemma~\ref{c:lem:dual}, it suffices to solve $Gw=Er$ to absolute dual
residual at most
\[
 \frac{\eta}{3\rho_{\mathrm{ridge}}}\,
       \delta_0\sqrt{\eta/L_C}\,\|r\|,
 \qquad z=\eta^{-1}(r-PE^\top w).
\]
The factor $3$ uses $\|E\|\le3$. If $r=0$, the outer solve has finished.
A correct dual declaration gives eigenvalues in $[1/4,7/4]$ for its
preconditioned target. Richardson iteration with step $4/7$ and inverse
action from the stored dual Cholesky factor therefore contracts by at
most $6/7$. The displayed absolute residual is reached in a logarithmic
number of iterations, and is explicitly checked.

We can bound the outer error directly.
If $e_z=z-\mathcal C_H^{-1}r$, then
\[
 \|e_z\|_{\mathcal C_H}\le\|\mathcal C_Hz-r\|/\sqrt\eta
 \le\delta_0\|r\|/\sqrt{L_C}
 \le\delta_0\|r\|_{\mathcal C_H^{-1}}.
\]
Use the update $y^+=y+z/\chi$. Its exact part contracts in
$\mathsf A$-norm by $1-1/\chi$, since
$\mathcal C_H\preceq\mathsf A\preceq\chi \mathcal C_H$.
For $e=y-\mathsf A^{-1}\widetilde b$,
\[
 \|e_z/\chi\|_{\mathsf A}
 \le\|e_z\|_{\mathcal C_H}/\sqrt{\chi}
 \le\delta_0\|e\|_{\mathsf A},
\]
where the last inequality uses
$\mathsf A \mathcal C_H^{-1}\mathsf A\preceq\chi\mathsf A$.
The total contraction is therefore at most $1-3/(4\chi)$.
All inner and outer iterates lie exactly in the retained row-feasible space.

\paragraph{Final certificate and dense fallback.}\label{c:failure}
In both branches stop only after checking the actual current-state test
\begin{equation}\label{c:finalresidual}
 \|\widetilde b-\mathsf Ay\|^2\le\zeta^2\eta s/4,
 \qquad\zeta=\min\{\eta/(1000C),1/(100s)\},\quad C=1200.
\end{equation}
Here $s$ is the actual retained sampler dimension, and $\widetilde b$ is its actual approximate-square-root right-hand side.
This is exactly the certificate used in Section~\ref{a:direction}.
The covariance rejection and normalization tests are unchanged.

The spectrum of $\mathsf A$ is in
$[\eta,\mu+9\rho_{\mathrm{ridge}}]$, with polynomial ratio in
$s$ and polylogarithmic factors in the original input dimensions.
Consequently the energy contraction estimates above imply the required
Euclidean residual bound in a prescribed polylogarithmic number of
iterations; the same observation covers the absolute inner tolerance.
These caps can be chosen from the displayed scalar bounds and the initial
residual norm by integer doubling tests. The final residual is checked even after a good declaration.
If a cap is reached without its required certificate, form and solve the
current arithmetic system densely. Then set
\[
 B^{\rm p}\leftarrow \mathcal C,\qquad B^{\rm d}\leftarrow G
\]
\emph{both exactly}, rebuild their Cholesky factors and $Z$, and continue
using the certified dense solution. The fallback costs
$\widetilde O(ms^2+s^3)$ and makes both finite-target cache divergences zero. The analytic
potential used later can increase by the explicitly charged reset amount.
A correct declaration cannot trigger it. At most one such fallback is
needed in a fixed state, since all further trials reuse exact current
caches.



\section{A conditional cubic implementation by two-sided secants}
\label{d:section}

The controller below uses the established direction sampler and finite-step
proof. It specifies the cache lifetime and accounting potential while preserving
the acceptance certificates and discrepancy bound. Its correctness and
termination are unconditional. The expected cubic arithmetic bound has
one explicitly stated, unproved hypothesis.

\subsection{Fixed implementation and inherited contracts}\label{d:contracts}
Fix a particular implementation $\mathcal I$ of the certified arithmetic
primitives in Sections~\ref{p:section}, \ref{r:finite}, and
\ref{a:algorithm}. Fix their coefficient evaluation, precision schedules,
polynomial degree, normalization, and tie rules before the run. Use fresh
independent random bits when a randomized primitive requests them and
accept its first output passing its stated certificate. At a state, form
the numerical coefficients before testing either cache; freeze these
coefficients throughout all sampling retries. The transpose of a finite
operator is its actual arithmetic transpose. Use the least allowed degree
and first precision level meeting each stated error test; take the
midpoint of a certified scalar interval. Break index ties by the least
index. Orient an eigenvector by making its first nonzero coordinate
positive; in the two-dimensional Ritz test choose the eigenvalue with
largest absolute value, preferring the positive one at a tie. Choose
the largest dyadic normalizer consistent with the required norm upper
bound. These conventions fix the remaining elementary choices.

These conventions fix the probability law, including the choice of
finite approximations at each state. All expectations
below include the randomness of the certified Perron primitive, cache
tests, direction samples, and boundary coins. Constants may depend on this
one fixed implementation, but never on the input matrix or its dimensions.

The signing algorithm starts at $x=0$ on a matrix $A\in\mathbb R^{m\times n}$
whose columns have Euclidean norm at most one. We use the contracts proved above. Every certified direction preserves
the row barriers at either capped endpoint, has norm in $[1/2,1]$, and
satisfies $x^Th=0$. For every past history, the sampler succeeds with
probability at least $1/(4C)$, where $C=1200$. Exact row maintenance has
$\widetilde O(n)$ elementary events, the number of accepted moves is
$\widetilde O(n)$ pathwise, and all non-solve work has expected cost
$\widetilde O(mn^2+n^3)$.
The final residual is always
\begin{equation}\label{d:residual}
 \|\widetilde b-\mathsf Ay\|^2\le\zeta^2\eta s/4,
 \qquad \zeta=\min\{\eta/(1000C),1/(100s)\},
\end{equation}
where $s$ is the actual retained dimension. The covariance rejection and
dyadic normalization tests remain those of Section~\ref{a:direction}.

\subsection{Rebuild at metadata cleanup}\label{d:epochs}
One cache epoch is one fixed-metadata epoch. Let $s$ be its retained
coordinate count and $N$ its halving-phase initial count. While the large
sampler is used, $N/2<s\le N$. Newly frozen coordinates remain in the
retained data until cleanup and are constrained by exact zero equations.
Thus the ambient dimension $s$, the ridge
$\eta=1/(4096JN)$, and all scalar parameters in
Section~\ref{a:direction} are fixed throughout this epoch. Use precisely
that local schedule for $\mu=\eta+4096/s$, $\chi=\mu/\eta$, $K$, $\tau$,
$e_*$, and $\rho=\rho_{\mathrm{ridge}}$.
All matrices act on the retained coordinates of this epoch.

At metadata cleanup, discard both solver caches. Initialize fresh caches
only when the new epoch first visits the high branch. In particular,
changes of retained profiles, scales, offsets, filters, or phase ridges
across cleanup do not enter the production hypothesis.

The decrease in epoch sizes bounds the total rebuilding cost. Except at the end of a halving
phase or the bounded-dimension continuation, cleanup freezes at least
$\lceil s/100\rceil$ coordinates and reduces the next retained dimension
to at most $99s/100$. Hence there are only an absolute constant number of
metadata epochs in each halving phase. Writing their sizes as $s_e$,
\begin{equation}\label{d:epoch-sums}
 \sum_e s_e=O(n),\qquad \sum_e s_e^2=O(n^2),\qquad
 \sum_e s_e^3=O(n^3).
\end{equation}
A fresh pair of dense caches costs
$\widetilde O(ms_e^2+s_e^3)$, so all rebuilds fit the desired total budget.

\subsection{A canonical factor for the analysis}\label{d:ideal-factor}
Within a fixed-metadata epoch let $M(x)$ be the true retained matrix,
$L(x)=\lambda_1(M(x))$, and $v(x)$ its positive unit Perron vector. Put
\begin{equation}\label{d:canonical}
 \Lambda_\#(x)=\max\{1,L(x)\},\qquad
 T(x)h=M'(x)[h]v(x),\qquad
 E_\#(x)=f\bigl(M(x)/\Lambda_\#(x)\bigr)T(x).
\end{equation}
The bounded Bernstein polynomial $f$ and its degree are fixed for the
whole metadata epoch. Define this factor analytically even at a low
state; the algorithm need not compute it there.

This normalization preserves the filter proof. In a certified high state,
$L>1-\delta$, whence
$1\le\Lambda_\#\le L+\delta\le1.001L$. The high-rank count uses
$\Lambda_\#\ge1$, and the original leakage estimate uses
$\Lambda_\#\le1.001L$. Both therefore remain valid. At every admissible
state, including a low state, the retained-row bounds give
$\operatorname{tr}\mathcal D\le2L\le9/4$ and
$\|T\|^2\le L\rho_0\operatorname{tr}\mathcal D<4$.
Since $0\le f\le1$ on $[0,1]$, we have $\|E_\#\|\le2$ everywhere.

At a high state construct a finite arithmetic factor $E$ with
\begin{equation}\label{d:factor-accuracy}
 \|E-E_\#\|\le\varepsilon_E,\qquad
 \varepsilon_E\le e_*/4,\qquad
 (2\varepsilon_f+\varepsilon_E)^2\le\eta/(1200\rho).
\end{equation}
The usual common-coefficient curvature certificate is also required.
To specify the filter tolerance completely, one may take
$\varepsilon_f=\varepsilon_E=\varepsilon$, with $\varepsilon$ the largest
negative power of two satisfying the displayed inequalities and
$\varepsilon\le1/4$. Let $q_f$ be the least integer with
$2^{q_f}\ge1/\varepsilon$, and choose the least even degree at least
$\max\{2000,125000q_f\}$. The binomial-tail bound in
Section~\ref{r:filter} proves the required filter accuracy. These choices
are made once per metadata epoch.

To implement this normalization, retain the original certified
high/low decision and independently refine a Perron interval until its
upper endpoint $\Lambda_{\rm ref}$ has the required scalar accuracy.
Use $\overline\Lambda=\max\{1,\Lambda_{\rm ref}\}$. If the refined interval
has width $a$, then
$|\overline\Lambda-\Lambda_\#|\le a$, and
\[
 \|\overline M/\overline\Lambda-M/\Lambda_\#\|
 \le\|\overline M-M\|+2a.
\]
The bound $d_f6^{d_f}$ from Section~\ref{r:finite} lets us include this
scalar error, together with the matrix and Perron-vector errors, in the
budget~\eqref{d:factor-accuracy}. All reciprocal accuracies remain polynomial in
the input dimensions. This is a refinement of the existing finite
primitive, with the same soft arithmetic cost.

Write $\alpha=\eta/\rho$. The arithmetic and analytic targets are
\begin{align}\label{d:targets}
 \mathcal C&=\eta I+\rho E^TE,&
 G&=\alpha I+EPE^T,\\
 \mathcal C_\#&=\eta I+\rho E_\#^TE_\#,&
 G_\#&=\alpha I+E_\#PE_\#^T.\notag
\end{align}
Only the arithmetic targets are used for solves, witnesses, and BFGS
updates. The analytic targets are used solely in the proof.
The error in \eqref{d:factor-accuracy} implies, for either pair of targets,
\begin{equation}\label{d:relative}
 \frac{15}{16}K_\#\preceq K\preceq\frac{17}{16}K_\#.
\end{equation}
Indeed, for $a=\varepsilon_E/\sqrt\alpha\le1/\sqrt{1200}$, the triangle
inequality on the stacked maps $[\sqrt\alpha I;E]$ gives the sharper
factors $(1-a)^2$ and $(1+a)^2$. The same argument applies to $PE^T$ for
the dual Gram. These factors lie inside the displayed interval.

\subsection{Arithmetic secants decrease the analytic potential}
\label{d:robust-drop}
Let $\phi(r)=\log((1+r)/(2\sqrt r))$ for $r>0$.
\begin{lemma}[Robust secant decrement]\label{d:drop}
Suppose \eqref{d:relative} holds and a verified arithmetic witness obeys
$u^TKu/(u^TBu)\notin[1/2,3/2]$. The BFGS update using $Ku$ satisfies
\[
 \mathcal S(B,K_\#)-\mathcal S(B^+,K_\#)
 \ge\gamma_0:=\phi(17/24)-\phi(15/16)-\log(466/465)>1/100.
\]
Also $\mathcal S(K,K_\#)<s/1000$.
\end{lemma}
\begin{proof}
Whiten $K_\#$ and rotate the normalized witness to $e_1$. Write
\[
 B=\begin{pmatrix}b&w^T\\w&D+ww^T/b\end{pmatrix},\qquad
 Ke_1=\binom{q}{a},\qquad D\succ0.
\]
Then $15/16\le q\le17/16$, $\|a\|\le1/16$, and
$B^+=\left(\begin{smallmatrix}q&a^T\\a&D+aa^T/q\end{smallmatrix}\right)$.
The same Schur determinant calculation as for the exact decrement gives
\begin{align*}
 \mathcal S(B,I)-\mathcal S(B^+,I)
 &=\phi(b)-\phi(q)
 +\log\left(1+\frac{w^T(D+I)^{-1}w}{b(1+b)}\right)\\
 &\quad-\log\left(1+\frac{a^T(D+I)^{-1}a}{q(1+q)}\right).
\end{align*}
The witness implies $b<17/24$ or $b>15/8$, hence
$\phi(b)\ge\phi(17/24)$. Also
$\phi(q)\le\phi(15/16)$, $q(1+q)\ge465/256$, and $(D+I)^{-1}\preceq I$.
These bounds prove the claimed decrement. Its strict lower bound follows
by squaring the rational inequality
\[
 e^{\gamma_0}=\frac{82\cdot465}{31\cdot466}\sqrt{5/34}>100/99,
 \qquad \log(100/99)>1/100.
\]
Finally, all generalized eigenvalues of $(K,K_\#)$ lie between $15/16$
and $17/16$, so $\mathcal S(K,K_\#)\le s\phi(15/16)<s/1000$.
\end{proof}

Thus finite coefficient or Perron-vector errors within the specified
budget require no separate variation hypothesis. Resetting both caches to
the current arithmetic targets injects at most $s/500$ into their total
analytic potential, regardless of the old caches.

\subsection{Controller and certification}\label{d:controller}
The following rules, together with Section~\ref{a:algorithm}, specify the
implementation. Every solve uses the unchanged current arithmetic
operator $\mathsf A=P(F+\rho E^TE)P|_{\operatorname{ran}P}$.
\begin{enumerate}
\item At a low state use $E=0$, $F=\eta I$ in the sampler and solve by
division by $\eta$. Do not test or repair either solver cache. Existing
caches are kept through the low excursion. Maintain the exact row
structure and, if initialized, the constrained primal inverse $Z$.
\item At the first high state of a metadata epoch, form both arithmetic
targets in \eqref{d:targets}, initialize $B^{\rm p}=\mathcal C$,
$B^{\rm d}=G$, their Cholesky factors, and the constrained inverse $Z$.
At later high states update the exact row data and $Z$.
\item Test both caches using Section~\ref{c:tester}. If one declares good,
use it, preferring the primal cache if both do. Otherwise use the verified
witness for one BFGS correction, alternating primal and dual corrections
within the state and beginning with primal. Retest both and stop as soon
as one declares good. Never repair an idle cache solely because it is stale.
\item Use the explicit primal or dual Richardson routine of
Section~\ref{c:sec:solve}, with the actual $\chi$ in place of its
conservative envelope. Enforce its prescribed iteration caps and the
actual residual \eqref{d:residual}. If a cap is reached without the
required certificate, solve the current arithmetic system densely and
reset both caches exactly to the current arithmetic targets. Rebuild
their Cholesky factors and $Z$. Then continue with that certified solution.
\item Sampling rejection keeps the state, coefficients, and prepared
caches fixed. Once a sample passes the unchanged quadratic tests, make
the mean-zero boundary coin move of Section~\ref{a:move}. At metadata
cleanup discard both caches; at bounded dimension use the established
baseline continuation.
\end{enumerate}
A correct good declaration cannot cause a dense fallback. After one
fallback in a state both caches are exact, so further sampling retries
in that state need no additional fallback.

The failure schedule is independent of the conjecture.
Put $\eta_{\min}=1/(4096Jn)$ and $H_{\max}=\lceil\log_2(2/\eta_{\min})\rceil$,
and let $\tau_{\min}$ be the largest dyadic satisfying the step tests with
$K=1$, $\eta=\eta_{\min}$, and $H=H_{\max}$. Every actual full step has
$|d|\ge\tau_{\min}$, hence Euclidean length at least $\tau_{\min}/2$.
Consequently
\[
 \overline T=\lceil4n/\tau_{\min}^2\rceil+n+2
\]
is a deterministic polynomial upper bound on the accepted states.
In fact $\eta_{\min}H_{\max}^2\le1$, so
$\tau_{\min}^{-1}=O(Jn)$ and $\overline T=O(n^3J^2)$.
Set
\[
 \overline e=\min\{\eta_{\min}^2/256,
                         \eta_{\min}\tau_{\min}/128\},\qquad
 \overline\mu=4097,
\]
let $\overline\rho$ be the least power of four at least
$512C\overline\mu/\overline e^2$, and put
$\overline\kappa=1+9\overline\rho/\eta_{\min}$.
Indeed $K\ge\eta\ge\eta_{\min}$, $\tau\ge\tau_{\min}$, and
$\mu\le\overline\mu$, so these are deterministic polynomial bounds for
all arithmetic target condition numbers. Choose the least positive
integer $L_\kappa$ with $2^{L_\kappa}\ge\overline\kappa$.
The exact arithmetic-target decrement and full-rank target-drift bound
give a deterministic repair cap
$\overline R=\lceil50n(\overline T+1)L_\kappa\rceil$.
This cap follows from a separate potential calculation for the finite targets (taking the sampler's
$E=0$ at low states): every finite-target initialization or fallback
reset has zero divergence, and each verified correction drops the
finite-target potential by more than $1/50$. Thus this cap does not rely
on the analytic production conjecture or on a bound for analytic reset
injections.
There are consequently at most
$\overline N=2(\overline T+\overline R+2)$ individual cache tests.
Use fresh test randomness with
\begin{equation}\label{d:failure}
 \delta_{\rm test}\le[\overline N(2+m+n)^{10}]^{-1}.
\end{equation}
The probability of any false good declaration is at most
$(2+m+n)^{-10}$. There are at most $\overline T$ fallbacks, each of cost
$\widetilde O(mn^2+n^3)$ and analytic-potential injection $O(n)$.
Their expected cost and injection are therefore within
$\widetilde O(mn^2+n^3)$ and $\widetilde O(n)$, respectively. All
certificates remain valid on the failure event as well.

\subsection{The precise global production statement}\label{d:production}
If a metadata epoch never visits a high state, assign it zero production.
Otherwise begin accounting at its first high state, after initialization,
and include every subsequent direction-selection state in that epoch,
including low states. Stop at its last such state before cache disposal.
No endpoint after which metadata are changed is compared with the new
metadata in this ledger.

Write $B_t^{\rm p},B_t^{\rm d}$ for the caches after all corrections and
any fallback at state $t$. At a low state their matrix values are
unchanged. Put
\[
 \widehat G_{t+1}=\alpha I+E_{\#,t+1}P_tE_{\#,t+1}^T
\]
and separate the target change into a factor change followed by a row
projector change. The corresponding changes in divergence are
\begin{align}\label{d:production-def}
 \Delta_t^\#
 &=\mathcal S(B_t^{\rm p},\mathcal C_{\#,t+1})
       -\mathcal S(B_t^{\rm p},\mathcal C_{\#,t})\notag\\
 &\quad+\mathcal S(B_t^{\rm d},\widehat G_{t+1})
       -\mathcal S(B_t^{\rm d},G_{\#,t}),\\
 \Delta_t^{\rm row}
 &=\mathcal S(B_t^{\rm d},G_{\#,t+1})
       -\mathcal S(B_t^{\rm d},\widehat G_{t+1}).\notag
\end{align}
Let $\Pi=\sum_e\sum_t\Delta_t^\#$, with the inner sum only over adjacent
states of that epoch's ledger. This signed sum retains cancellations between positive and negative terms.

Reversing the order of the factor and projector changes gives a factor
ledger differing by at most
$\sum_t r_t\log\kappa_t$, where $r_t=\operatorname{rank}(P_{t+1}-P_t)$
and $\kappa_t=1+4\rho/\eta$ uses the fixed parameters of that transition's
epoch. Both analytic targets have this common spectral-range ratio.
This follows by bounding
each of the two possible row productions by
$\tfrac12r_t\log\kappa_t$. Inserting intermediate low states merely
telescopes the total target production while the caches are unchanged;
the factor-ledger difference is again covered by the row budget.
Including all states therefore lets us condition before each boundary
coin, without assuming smoothness over an entire low excursion.

\begin{theorem}[Unconditional accounting]\label{d:accounting}
Let $R$ count all verified BFGS corrections and let $N_e^{\rm fb}$ count
dense fallback resets in epoch $e$. For every execution,
\begin{equation}\label{d:account}
 \gamma_0 R\le
 \Pi+\frac12\sum_t r_t\log\kappa_t
       +\frac1{500}\sum_e s_e(1+N_e^{\rm fb}).
\end{equation}
In particular the row term and the expected initialization/reset term
are $\widetilde O(n)$. The random variable $\Pi$ is integrable without
assuming the production hypothesis.
\end{theorem}
\begin{proof}
In each ledger use the nonnegative potential
$\mathcal S(B^{\rm p},\mathcal C_\#)+\mathcal S(B^{\rm d},G_\#)$.
Its initial value is at most $s_e/500$. Every verified correction lowers
it by at least $\gamma_0$ by Lemma~\ref{d:drop}; a fallback increases it
by at most $s_e/500$. A transition changes it by exactly
$\Delta_t^\#+\Delta_t^{\rm row}$. The row term is bounded by
$r_t\log\kappa_t/2$ by Lemma~\ref{c:lem:drift}. Telescope and discard
the nonnegative terminal potentials. Equations~\eqref{d:epoch-sums} and
\eqref{d:failure} prove the claimed budgets.
Finally the cache-uniform full-rank drift estimate bounds each
$|\Delta_t^\#|$ by $O(s_e\log\kappa_t)$. There are at most
$\overline T$ states and all condition ratios are polynomially bounded,
so $\Pi$ has a deterministic polynomial bound in absolute value.
\end{proof}

\begin{quote}
\textbf{Unproved production hypothesis for the fixed implementation $\mathcal I$.}
There are absolute constants $A_0,p>0$, independent of $m,n$ and $A$,
such that for every $A\in\mathbb R^{m\times n}$ with column norms at
most one, the preceding complete algorithm started at $x=0$ satisfies
\begin{equation}\label{d:conjecture}
 \mathbb E\Pi\le A_0 n\log^p(2+m+n).
\end{equation}
\end{quote}
This concerns the reachable states and adaptive caches of this particular
algorithm. It makes no assertion for arbitrary bounded factors,
arbitrary feasible fractional starting states, or every admissible
choice of numerical proxies. Nor does it require a separate bound
conditional on every epoch-start history. It remains unproved.

\begin{theorem}[Conditional cubic implementation]\label{d:conditional}
For the fixed implementation $\mathcal I$, the algorithm terminates
almost surely, and every returned signing has discrepancy strictly less
than $105+3\sqrt{99}$. If \eqref{d:conjecture} holds, its expected total
cost is $\widetilde O(mn^2+n^3)$ real-arithmetic operations with exact
square roots and comparisons. No bit-complexity claim is made.
\end{theorem}
\begin{proof}
Equation~\eqref{d:account} and the production hypothesis give
$\mathbb ER=\widetilde O(n)$. Every correction, cache test, or certified
iterative solve costs at most $\widetilde O(mn+n^2)$. The number of tests
is at most twice the number of states plus corrections. The expected
number of sampling trials is $\widetilde O(n)$ by the uniform conditional
success probability. Initialization costs sum to
$\widetilde O(mn^2+n^3)$ by \eqref{d:epoch-sums}; exact row maintenance,
all other work, and the rare dense fallbacks obey that same budget.
Multiplying the global repair count by the global per-repair cost proves
the asserted total. The global production hypothesis suffices for this count.

Independently of the hypothesis, the polynomial cap on verified repairs,
the certified dense fallback, and the positive conditional sampler
success probability give almost-sure termination. Every accepted solve
passes \eqref{d:residual}; hence the inherited direction, both-sign
finite-step, tracking, and terminal-rounding proofs give the stated
discrepancy on every returned output.
\end{proof}

The theorem gives an arithmetic complexity bound. A numerical experiment can
check the realized production, repair count, and all residual and
acceptance certificates for the algorithm it implements. It cannot prove
\eqref{d:conjecture}; simulations that alter the dimension thresholds,
filter, ridge, or accepted-direction law test a different model and must
be identified as such.


\paragraph{Relation to inverse-maintenance amortization.}
Lee--Sidford~\cite{LSmaint} achieve amortized
$\widetilde O(\operatorname{nnz}(A)+d^2)$ maintenance for $A^TD_tA$
under relative diagonal stability. Jiang--Lee--Song--Wong~\cite{JLSW}
obtain $O(n^2)$ additional work per cutting-plane oracle call using
layered maintenance. These results motivate preserving a cache, but
their stability hypotheses are not supplied by ordinary factor energy.
The general SDP rank inequality of Huang et al.~\cite{HJSTZ} gives
$\sum_t r_tg_{r_t}=O(T\|g\|_2\log s)$ for nonincreasing weights $g$
under their relative-matrix stability assumptions. With $g_i=1$ this
controls only $\widetilde O(T\sqrt s)$ total rank, whereas our dense
rank-update implementation needs $\widetilde O(s)$ corrections over
$\widetilde O(s)$ states. No cited maintenance theorem closes
\eqref{d:conjecture}.


\section{A deterministic variant below 67}\label{opt:section}
This section retains the earlier profile and its two-profile comparison.
The stronger final construction and the screened arithmetic bounds are
presented in Section~\ref{fin:section}; the earlier benchmark records below
are not measurements of that final construction.

The optimized-profile supplement~\cite{Below67} improves the deterministic
\emph{discrepancy guarantee}, while retaining the arithmetic exponents of
Theorem~\ref{b:thm:main}. We retain the quadratic walk as a separate method:
it underlies the modular randomized construction, has simpler energy
updates, and permits a wider tracked-sum certificate for cheap endpoint
moves. This section specifies the alternative, explains how endpoint
acceleration transfers, and separates its proved bound from numerical
performance.

\subsection{The coupled estimate and the coordinate profile}
The spectral improvement is to bound the Gram factor and the low-spectrum
resolvent together. If $CC^\T\preceq M$ and $R_a$ is the reduced
resolvent restricted to eigenvalues below $aL$, then
\begin{equation}\label{opt:coupled}
 C^\T R_aC\preceq\frac{a}{1-a}I.
\end{equation}
Indeed, $R_a^{1/2}CC^\T R_a^{1/2}\preceq R_a^{1/2}MR_a^{1/2}$,
whose eigenvalues on its support are $\lambda/(L-\lambda)<a/(1-a)$.
This is the coupled estimate in~\cite[Lemma~2.4]{LiSpectral}.
Writing $M'[h]v=C\ell$ and $\|\ell\|^2=h^\T\mathcal Uh$ gives
\[
 L''[h,h]\le h^\T\left(\frac{1+a}{1-a}\mathcal U-\mathcal D\right)h
     +\frac2\eta\|P_{[aL,L]}M'[h]v\|^2.
\]
The improvement is available for the quadratic profile as well: at its
cutoff $a=21/100$, the factor $279/79$ can be replaced by $121/79$ in
this local curvature estimate. It does not by itself improve the
quadratic profile's terminal discrepancy constant.

For the alternative profile, put
\begin{equation}\label{opt:profile}
 f(x)=\frac{275}{34}\log\frac{25-17|x|}{8}
       +\frac{11}{2}(|x|-1),\qquad -1\le x\le1.
\end{equation}
This even concave function vanishes at the endpoints, is $C^{2,1}$,
and satisfies $0\le f(x)<93/25$ and
\[
 f'(x)=-\frac{187x}{2(25-17|x|)},\qquad
 -f''(x)=\frac{4675}{2(25-17|x|)^2},\qquad
 \frac{(1+(2/11)|f'(x)|)^2}{-f''(x)}=\frac{50}{187}.
\]
Thus a diffuseness threshold $p_{ij}\le(4/121)q_i$ gives the pointwise
comparison needed for diagonal curvature domination. The profile is
optimal for this specified pointwise condition, not among all signing
algorithms.

Use the same retained sets and affine offsets, but replace the medium-row
energy by $E_i=\sum_jp_{ij}f(x_j)$ and take
\begin{equation}\label{opt:parameters}
 \begin{gathered}
 R=149/100,\quad \Delta=4/121,\quad \theta=1001/1000,\quad
 H_0=689/20,\quad b=218/25,\quad \beta=4/5,\\
 u_{i\sigma}=(\sigma d_i-H_0)/r_i+E_i/r_i^2+b,\qquad z_i=p_i/r_i^2,\\
 M=\sum_{i,\sigma}36e^{\beta u_{i\sigma}}z_iz_i^\T
   +\frac4{625}\sum_{i\text{ large}}\diag(p_i)+\eta\one\one^\T.
 \end{gathered}
\end{equation}
Large rows have $q_i>R^2$ and remain exactly constrained. Delete medium
entries with $p_{ij}>\Delta q_i$ and retain scales
$q_i\le r_i^2<\theta^2q_i$, $r_i\le R$. The original cleanup order
preserves feasibility and decreases the spectral potential for these
parameters too; the identity $\beta(b-93/25)=4$ supplies scale
monotonicity.

The coupled cutoff $a=59/250$ gives
$\mathcal U_{jj}\le(40/187)\mathcal D_{jj}$ and a curvature factor
$309/191$. The protected-row, spectral-cancellation, and curvature
budgets leave more than $s/200$ feasible negative directions for
$s\ge513$. At a phase start of size $S$, use
$\eta=(2^{40}JS)^{-1}$, accuracy $\delta=\eta^5/2^{100}$,
and spectral ceiling $257/256$. The full dimension calculation,
rational scalar interfaces, and finite-step proof for the $C^{2,1}$
curves are given in~\cite[Sections~4--7 and 9]{Below67}.
The third derivative can jump when a coordinate crosses zero; the
quadratic-weight finite-step lemma must therefore be replaced by that
$C^{2,1}$ argument.

\subsection{Terminal rounding and the improved guarantee}
At $s\le512$, conditional expectation chooses the remaining signs so that
\begin{equation}\label{opt:rounding}
 \|A_V(\varepsilon_V-x_V)\|_2^2+
       4\|\varepsilon_V-x_V\|_2^2\le5s.
\end{equation}
Independent signs with means $x_j$ have expected left side
$\sum_{j\in V}(1-x_j^2)(\|A_{:j}\|^2+4)\le5s$.
Fixing a sign with the smaller conditional expectation preserves this
bound, using $O(ms+s)$ work with a maintained partial row-sum vector.

Write $A_\Delta=(11+\sqrt{117})/2$.
The deleted absolute mass of a medium row is at most
$A_\Delta(R-\sqrt{q_i})$: the decrease of $\sqrt q$ at a deliberate
deletion pays for its absolute coefficient, while freezing only
reduces the remaining mass. Hence a medium row's final discrepancy is
at most
\[
 H_0+2RA_\Delta+
       (\sqrt{640}-b-2A_\Delta)\sqrt{q_i}<H_0+2RA_\Delta.
\]
Here its barriers give $|d_i|\le H_0-b\sqrt{q_i}$, and
\eqref{opt:rounding} bounds the rounding movement by $\sqrt{640}$.
Completed rows inherit $|d_i|<H_0$ and the deletion allowance; large
rows have no deliberate deletion and discrepancy at most
$\sqrt{2560}$. The same rounding directly handles inputs with
$n\le512$.

\begin{theorem}[Optimized deterministic profile, \cite{Below67}]
\label{opt:guarantee}
There is a deterministic signing algorithm with
\[
 \|A\varepsilon\|_\infty<C_{67}:=
 \frac{1271}{25}+\frac{149}{100}\sqrt{117}
 =66.956814201324\ldots<67.
\]
It uses $O(nJ^2q^2)$ updates and
$\widetilde O(mn^3+n^4)$ field operations and comparisons, with rational
states for rational input. With the imported fast primitives and square
roots its cost is $\widetilde O(mn^\omega+n^{\omega+1})$.
\end{theorem}
\begin{proof}[Implementation and scope]
Use the parameters and terminal rule above with the supplement's
exactly feasible negative subspace, deterministic flattening, and
finite-step schedule. Its spectral-budget proof gives $O(nJ^2q^2)$
updates, each costing $\widetilde O(ms^2+s^3+m)$, or
$\widetilde O(ms^{\omega-1}+s^\omega)$ with the fast interface.
The geometric phase sum gives the stated costs. The preceding tracking
and rounding argument gives $C_{67}$. These are arithmetic-operation
bounds; neither a polynomial bit bound nor a practical time bound is
asserted.
\end{proof}

This theorem does not change Theorem~\ref{i:modular}. The anchored
quadratic energy identity used there is not an identity for
\eqref{opt:profile}; replacing its constants alone would not prove a
below-67 filtered-covariance algorithm.

\paragraph{The terminal improvement also transfers to the quadratic walk.}
One need not change the coordinate profile to use
\eqref{opt:rounding} at $s\le512$. In the quadratic state let
$A_{\rm old}=10+\sqrt{99}$. Keeping the remaining mass in the deletion
estimate bounds the deliberately removed mass by
$A_{\rm old}(3/2-\sqrt{q_i})$. A surviving medium row therefore has
final discrepancy at most
\[
 75+3A_{\rm old}+
       (\sqrt{640}-35/2-2A_{\rm old})\sqrt{q_i}<D_*.
\]
The coefficient is negative. Completed rows retain their old bound, and
large rows are bounded by $\sqrt{2560}<D_*$. Thus the quadratic walk
can stop at 512 and use conditional expectation while preserving $D_*$.
This variant transfers a useful part of the new algorithm to the
existing method. Comparing it with the optimized profile matches both
the stopping threshold and the terminal objective.

\subsection{Endpoint acceleration for the alternative profile}
Starting from zero, any sequence of cube-contained moves preserving the
large-row sums and the new medium barriers gives the bound $C_{67}$
when it reaches $s\le512$ and uses \eqref{opt:rounding}. Cleanup and
the tracking argument require only endpoint feasibility, as in
Lemma~\ref{ep:correctness}. Thus the practical endpoint strategy remains
available despite the nonquadratic energy.

A cheap certificate is narrower than the quadratic one. If every medium
row at an endpoint satisfies $|d_i|\le9$, then
\begin{equation}\label{opt:envelope}
 L\le\eta s+\max\left\{\frac4{625},\frac{72}{R^2}
   \exp\!\left(\frac45\left[\frac{9-H_0}{R}+\frac{93}{25}+b\right]\right)\right\}
 <\frac45,
 \qquad u_{i\sigma}<-4.64.
\end{equation}
To see this, use $E_i/r_i^2\le93/25$ and
$p_ip_i^\T\preceq q_i\diag(p_i)$. The resulting medium coefficient
$(72/r_i^2)\exp(\beta[(9-H_0)/r_i+93/25+b])$ increases up to $R$;
combine its maximum with the large-row deposit using the column-norm
bound. Replacing the reciprocal exponential by the positive Taylor
sum through degree 19 proves the strict numerical inequality with
rational arithmetic. The accompanying schedule verifier checks it.
The comparable sufficient threshold for the quadratic walk is $25$
(Proposition~\ref{ep:envelope}); failure of either sufficient test need
not mean the endpoint itself is infeasible.

For a complete exact hybrid, put
$D_{67}=2^{41}\cdot3200Jq$ and
$\ell_{67}=2^{-40}/(64D_{67})$. Use the bounded trial rules of
Section~\ref{ep:hybrid}, with this displacement floor, the new endpoint
barriers, and ledger $\mathfrak b\ge\Phi$. On fallback use the
supplement's step, charging $-\eta\tau^2/4096$ for a full move and
$+\eta\tau^2/2^{16}$ for a shorter move. At a reset charge
\[
 \mathfrak b\leftarrow\mathfrak b+\frac1{2^{40}J}
                 -(\eta'-\eta)\|x\|^2.
\]
The same ledger induction gives
$(L-1)_+^3<3\cdot2^{-40}<(1/256)^3$ and
$O(nJ^2q^2)$ moves. With boundedly many cheap accepted cube trials at
every state, the maintained-projector accounting gives
$\widetilde O(mn^2+n^3)$ work, now with terminal bound $C_{67}$.
This is an execution-dependent endpoint result; it is not a transfer of
the modular randomized or secant-controller analysis.

\subsection{Which implementation to retain}\label{opt:comparison}
The theoretical schedules favor neither a claim of practical acceleration
nor replacement of the existing implementation. At $m=n=10^4$ in the
first phase, taking the conservative bound $K=1$, the largest permitted
step parameters are $2^{-50}$ for the quadratic walk and $2^{-98}$ for
the optimized profile. The latter also rescales the direction to norm
at most $1/16$, giving displacement at most $2^{-102}$.
Its prescribed accuracy has $\log_2(1/\delta)\simeq386$, compared with
$206$ for the quadratic walk. These are schedule calculations, not
observed running times: endpoint methods avoid these steps when their
certificates succeed.

The implementation keeps the quadratic profile as the default and adds
an explicit \texttt{below67} option. Both retain the same affine state,
cleanup order, constraint maintenance, and endpoint tests. The new option
uses the actual nonquadratic energy and derivatives, the coupled
curvature proxy, and conditional-expectation terminal rounding. The
quadratic option with Gram--Schmidt directions remains available.
A separate \texttt{quadratic-coupled} option tests the local curvature
improvement without changing the old profile or terminal guarantee.
An independent conditional-expectation finish at 512 is available for
the quadratic profile, with either direction routine.

We compared seven configurations on 48 matched cases: six matrix families
(Gaussian, Rademacher, sparse incidence, Hadamard, prefix, and heavy-row),
two seeds, $n=640$ with $m\in\{80,640,2560\}$, and $n=m=1024$.
All 336 runs completed and passed the numerical sign, discrepancy, and
retained-state checks. Table~\ref{tab:matched-profiles} reports the median of paired wall-time
ratios to the original quadratic configuration, the observed discrepancy
range, and the number of walk moves. CE512 and CE99 mean conditional
expectation at the indicated threshold; GSW means constrained
Gram--Schmidt directions.

\begin{table}[htbp]
\centering\small
\caption{Matched numerical comparison: 48 cases per configuration.}
\label{tab:matched-profiles}
\begin{tabular}{@{}lrrr@{}}
\toprule
Variant & Time/original & Discrepancy range & Walk moves\\
\midrule
Quadratic, original finish & 1.00 & $[0.336,6.022]$ & $537$--$925$\\
Below 67, CE512 & 0.33 & $[0.296,3.938]$ & $128$--$512$\\
Below 67, CE99 & 1.24 & $[0.336,5.101]$ & $536$--$925$\\
Quadratic, GSW finish & 1.60 & $[0.000,3.438]$ & $512$--$925$\\
Below 67 + GSW, CE512 & 1.05 & $[0.059,16.062]$ & $124$--$512$\\
Quadratic, CE512 & 0.27 & $[0.296,3.938]$ & $128$--$512$\\
Quadratic + GSW, CE512 & 0.84 & $[0.059,16.062]$ & $124$--$512$\\
\bottomrule
\end{tabular}
\end{table}

With the same CE512 finish, the new profile had median time ratios
$1.19$ for random kernel directions and $1.15$ for Gram--Schmidt directions,
relative to the corresponding quadratic option. Their discrepancies agreed
to $10^{-9}$ in 46 and 47 of the 48 cases, respectively; the new profile
was lower in the remaining cases. Thus the apparent speedup over the old
99-coordinate configuration largely comes from earlier completion, which
both profiles can use. Compared with the original quadratic GSW finish,
the new profile with GSW directions and CE512 had lower, equal, and higher
discrepancy in 14, 7, and 27 cases. These results favor retaining both
profiles and the original terminal routine.

Every run stayed in the low branch, so these timings do not measure the
coupled high-spectrum improvement. The largest tracking residual was below
$1.74\cdot10^{-14}$ and the largest large-row residual was below
$6.82\cdot10^{-13}$. Both profiles include the same repair for loss of
rank in the constraint projector: cleanup now transfers frozen entries
into the affine offsets before rebuilding a singular projector.


Early rounding can also worsen the observed maximum discrepancy.
If the remaining columns are orthonormal and $x_V=0$, the terminal
quadratic objective is $5s$ for every choice of signs. It supplies no
preference among their row maxima. For a normalized $256\times256$
Hadamard matrix, positive-sign tie breaking returns all positive signs,
with discrepancy $16$, while still satisfying the terminal energy bound.
Thus the early completion rule is a speed option, not an assurance of
better observed discrepancy. The original Gram--Schmidt finish is retained.


The comparison scripts time construction, walk, and completion together,
use matched matrices and seeds, and report failures explicitly. They
also include a below-67 run continued to 99 coordinates to distinguish
profile effects from the earlier terminal rule. All prototypes use
floating-point arithmetic and abort when no acceptable step is resolved;
they do not implement the exact fallback. Their final discrepancy tests
and residual checks are numerical evidence, not exact certificates or
validation of an asymptotic running-time exponent.

\section{The final constant, arithmetic bounds, and practical execution}
\label{fin:section}
The final construction~\cite{Final37} lowers the deterministic bound to
$C_{37}=1877/50=37.54$. Its proof uses a normalized spectral response
and counts protected rows, large rows, and positive curvature jointly.
This improves the signing constant without identifying the smallest
constant with the fastest arithmetic model or the fastest measured code.
We keep those three comparisons explicit.

\subsection{A joint response construction with discrepancy below 37.54}
\label{fin:construction}
Use the retained-row state of Section~\ref{opt:section}, with
\[
 R=\frac{7139}{5000},\quad \Delta=\left(\frac{52}{173}\right)^2,
 \quad\beta=\frac{17}{25},\quad A_0=\frac{25}{4},\quad B_0=\frac{7463}{10000},
\]
\[
 H_0=\frac{94893}{5000},\quad b=\frac{341613}{42500},\quad
 a_0=\frac{5389}{2500},\quad\theta=\frac{1001}{1000},
\]
\[
 f(x)=\frac{519}{104}\log(3-2|x|)+\frac{173}{52}(|x|-1).
\]
Here $A_0$ is the exponential amplitude and $B_0$ the large-row
diagonal coefficient. The profile is even, concave, and $C^{2,1}$,
with $f(\pm1)=0$, $0\le f<a_0$, and $\beta(b-a_0)=4$.
For $d=52/173$, its pointwise identity is
\[
 \frac{(1+d|f'|)^2}{-f''}=\frac{78}{173},\qquad
 \mathcal U_{jj}\le\rho\mathcal D_{jj},\quad \rho=\frac{1326}{4325}.
\]
The potential remains $(L-1)_+^3-\eta\|x\|^2$, with
$\eta=(2^{40}JS)^{-1}$ and $\delta=\eta^5/2^{100}$.

Write $M=B+W+\eta\one\one^\T$, where $B$ is the large-row diagonal,
and at a high-branch anchor define
\[
 S_L=(LI-B)^{-1/2},\qquad K_L=S_L(W+\eta\one\one^\T)S_L.
\]
The normalized matrix has Perron eigenvalue one. For $y\perp v$,
\begin{equation}\label{fin:response}
 y^\T(LI-M)^\dagger y
 =(S_Ly)^\T(I-K_L)^\dagger(S_Ly).
\end{equation}
This identity concerns quadratic forms on the indicated subspace;
pseudoinverses do not in general commute with congruence.
Let $N_K=V_KD_KV_K^\T$ be the rational reconstructed normalized
matrix, with $V_K$ exactly orthogonal, and put $k_* =\max_j(D_K)_{jj}$.
Retain the computed high columns $F$ with eigenvalues at least
$(18/25)k_*$, and let $R_K$ be the low resolvent with denominators
$k_*-(D_K)_{jj}$. With the warm-pair forcing matrix $\bar B_1$ and
diagonal inverse-square-root approximation $S$, use
\begin{equation}\label{fin:proxy}
 \bar Q=\bar{\mathcal U}-\bar{\mathcal D}
           +2\bar B_1^\T S R_K S\bar B_1,
 \qquad (V_K)_F^\T S\bar B_1h=0.
\end{equation}
The older scalar proxy $\gamma\mathcal U-\mathcal D$ cannot be
substituted into the new dimension calculation.

The reciprocal-Gram inertia bound combines cancellation and curvature.
Its supporting line is
\[
 \frac{1-\lambda}{1+\lambda}
 \ge\frac{23}{25}-\frac{32}{25}\lambda,
\]
with equality at $\lambda=1/4$. Put
$c_p=\theta^4/[A_0(1-\beta/100)]$,
$L_*=257/256$, and $\ell_0=999999/10^6$. The exact certificate gives
\[
 B_0R^2\left(c_p+\frac{(32/23)\Delta}{\ell_0-B_0}\right)<1.
\]
Consequently the combined loss for large-row equations, protected
gradients, high normalized spectral equations, and positive curvature
is at most
\[
 \left[R^{-2}+\frac{25}{23}\rho+\frac{32}{23}\Delta
                 +c_p(L_*-B_0)\right]s.
\]
Protected gradients have already been charged in this expression.
Including the two exact tangencies, the high row-energy space, and
the precision allowance leaves
\begin{equation}\label{fin:dimension}
 1-R^{-2}-\frac{25}{23}\rho-\frac{32}{23}\Delta-c_p(L_*-B_0)
 -\frac1{10000}-\frac2{337}-\frac1{4094}
 >\frac1{400}.
\end{equation}
The tangencies are $x_V^\T h=0$ and the maintained row-energy gradient
equation. The latter and the exclusion of computed row-Gram eigenvalues
at least 4095 give the row-energy bound used for screening.
Flatten coordinates and $g_r/3$, construct the feasible negative basis
at accuracy $\eta/2^{24}$, and set $h=h_0/16$. The dyadic flatness test
$9q/k$ gives $K^2\le\min\{1,14400q/s\}$. The step tests use
$c_*=2^{-56}$; the normalized finite-comparison proof in~\cite{Final37}
controls the remainder by
$2^{16}K^3|t|^3+2^{20}K^4t^4/\eta$.

\begin{theorem}[Final deterministic guarantee]\label{fin:guarantee}
For every real input with column Euclidean norms at most one, the
specified deterministic construction returns one sign per original
column with
\[
 \|A\varepsilon\|_\infty<C_{37}:=\frac{1877}{50}.
\]
It uses $O(n\log^4(2+n))$ updates and
$O(mn)+\widetilde O(n^4)$ ordinary field operations and comparisons.
Rational input gives rational states. No polynomial bit-length bound
or cubic total-time bound is asserted.
\end{theorem}
\begin{proof}[Proof and implementation dependencies]
The complete normalized-response and finite-step proof is in
\cite{Final37}; the screened finite core is supplied in~\cite{Screened67}.
The latter is a frozen source dependency, distinct from the earlier
unscreened supplement and its numerical implementation.
First discard rows with $\|a_i\|_1\le37$, which are safe for every
signing. The remaining count satisfies $m_{37}<n^2/37^2$, since each
surviving row has squared norm greater than $37^2/n$ and the total
squared mass is at most $n$. Preserve every original column. Empty
inputs and an empty surviving row set return immediately.

The new dimension bound and finite comparison preserve the phase and
descent arguments. Screening retains $w_s=O(n+s\log^2(2+n))$ warm
rows; evaluating their response and the dense direction matrices costs
$\widetilde O(w_ss^2+s^3)$ per state. A phase starting at $S$ has
$O(SJ^2q^2)$ updates, so the geometric sum of this work is
$\widetilde O(n^4)$. Input processing costs $O(mn)$, and the charged
support and row-statistic maintenance costs
$\widetilde O(m_{37}n^2)\subseteq\widetilde O(n^4)$.

At $s\le336$, conditional expectation with weight $19/6$ ensures
\begin{equation}\label{fin:rounding}
 \|A_V(\varepsilon_V-x_V)\|^2+\frac{19}{6}\|\varepsilon_V-x_V\|^2
 \le\frac{25}{6}s\le1400.
\end{equation}
Here $A$ denotes the surviving rows. Since $52^2+165^2=173^2$,
the deletion coefficient is $A_\Delta=13/2$.
Thus the remaining medium-row bound is
\[
 H_0+13R+\bigl(\sqrt{8400/19}-b-13\bigr)\sqrt{q_i}<H_0+13R=C_{37}.
\]
The strict inequality follows from $(b+13)^2>8400/19$.
Large rows have discrepancy at most $\sqrt{1400}<C_{37}$, and completed
rows retain the strict tracking bound. For $n\le336$, apply the same
rounding directly at zero. Discarded rows have discrepancy at most 37.
\end{proof}

\subsection{Which arithmetic bound is proved for which constant}
\label{fin:arithmetic}
The strongest constant and the lowest known exponent within these
constructions are separate claims. The screened below-67 baseline proves
\begin{equation}\label{fin:fast67}
 O(mn)+\widetilde O(m_{66}n^2+n^{\omega+1}),\qquad
 m_{66}=|\{i:\|a_i\|_1>66\}|<n^2/66^2,
\end{equation}
with the imported fast diagonalization and exact rank primitives,
allowing square roots~\cite[Section~11]{Screened67}.
For $m\le n$, the second term is $\widetilde O(n^{\omega+1})$.
For taller inputs, the explicit $m_{66}n^2$ maintenance term remains.
One may select this implementation or the ordinary
$O(mn)+\widetilde O(n^4)$ implementation after the input scan.

The final construction proves Theorem~\ref{fin:guarantee} using rational
Jacobi rotations and ordinary products. Its normalized joint count
uses an exactly orthogonal reconstructed spectral decomposition.
The fast interface of Appendix~\ref{app:arithmetic} has different
computed-vector contracts; reusing it here requires a compatible
reconstruction and error analysis. We therefore retain
\eqref{fin:fast67} at its proved discrepancy constant rather than
attach its exponent to $C_{37}$ without that argument.

The optional randomized screening in~\cite[Section~6]{Final37}
reduces \emph{row processing} to
$O(mn)+\widetilde O(n^3)$ expected work. It uses unbiased boundary
moves, a variance clock, final verification, and restart. The changing
dense negative subspace still costs quartic total work in the specified
ordinary implementation. This is not a cubic signing theorem.
The quadratic sampler and secant results elsewhere in this paper retain
their own bounds and hypotheses; the new nonquadratic profile does not
establish their below-37.54 counterparts.

\subsection{Fast execution with a checked output}
\label{fin:practical}
Worst-case step schedules need not dictate the first attempt to find a
signing. A fixed number of candidate signings costs $O(mn)$ to evaluate.
Choose a candidate by its measured row maximum, then verify it against
the original matrix. Any signing with a certified row maximum below
$C_{37}$ may be returned, regardless of how it was proposed.
If the bounded attempt fails, the exact algorithm of
Theorem~\ref{fin:guarantee} remains a complete fallback with the same
asymptotic bound. This preserves the strongest discrepancy guarantee
while allowing easy inputs to finish before the spectral walk begins.
Every accepted fixed-budget random-proposal execution costs $O(mn+n)$
work. It proves no uniform acceptance rate for the cheap attempts.

For a represented binary64 matrix, the implementation encloses each
row sum by adding its signed entries with outward rounding. If
$[l_i,u_i]$ contains the partial sum, the next interval is obtained
by rounding $l_i+a_{ij}\varepsilon_j$ downward and
$u_i+a_{ij}\varepsilon_j$ upward. Multiplication by a sign is exact.
Induction encloses the exact sum of the represented entries. A strict
comparison of all interval endpoints with a downward-rounded
$1877/50$ therefore certifies the requested output bound.
This check costs $O(mn)$ and does not rely on a BLAS error estimate or
on floating-point feasibility of the candidate's trajectory.

The practical code offers eight random sign proposals, a checked
quadratic GSW proposal, and endpoint walks for all three profiles.
GSW proposals are charged at their full construction cost; they are not
part of the linear-work random-proposal bound.
The final-profile high branch uses \eqref{fin:proxy}; its terminal rule
uses \eqref{fin:rounding}. It retains both row-energy constraints.
Its optional input filter is implemented, but lazy warm/cold screening,
rational finite arithmetic, and the exact tiny-step fallback are not.
Unresolved numerical walks report failure. Thus the present code
certifies each returned proposal but is not yet a complete executable
implementation of the worst-case theorem.

For comparison with the numerical steps, at $m=n=10^4$ in the first
phase and $K=1$, the final prescribed parameter is $\tau=2^{-114}$
and $\|\tau h\|\le2^{-118}$. A simple sufficient tracked-sum certificate
is also narrower: $|d_i|\le3/5$ on every medium row implies $L<1$,
compared with the sufficient thresholds 9 and 25 for the older profiles.
Indeed, the medium coefficient is at most
\[
 \frac{2A_0}{R^2}
 \exp\left(\beta\left[\frac{3/5-H_0}{R}+a_0+b\right]\right),
\]
which, combined with $B_0$ and the ridge, is below $0.993$.
The monotonicity condition $\beta(H_0-3/5)>2R$ and a positive Taylor
lower bound certify this rationally. These sufficient tests and tiny
schedules do not measure the actual endpoint or proposal performance.

We compared seven implementations on 48 matched cases: the six matrix
families of Section~\ref{opt:comparison}, two seeds, $n=640$ with
$m\in\{80,640,2560\}$, and $n=m=1024$. All 336 executions passed the
strict outward-rounded $37.54$ output check. Table~\ref{tab:final-comparison}
reports the new comparison; it includes construction and verification
costs and retains the original GSW finish as the reference. Method order
rotates, and all methods receive the same matrices and seeds.

\begin{table}[htbp]
\centering\small
\caption{Final-profile comparison. Times include construction and output verification; ratios are paired medians relative to the original quadratic GSW configuration.}
\label{tab:final-comparison}
\begin{tabular}{@{}lrrr@{}}
\toprule
Implementation & Time/GSW & Discrepancy range & Moves\\
\midrule
Quadratic, original GSW finish & 1.0000 & $[0.000,3.438]$ & $512$--$925$\\
Quadratic, CE336 & 0.4012 & $[0.415,4.384]$ & $304$--$688$\\
Below 67, CE336 & 0.4709 & $[0.415,4.384]$ & $304$--$688$\\
Final 37.54, no input filter & 0.6029 & $[0.316,4.250]$ & $304$--$690$\\
Final 37.54, input filter & 0.0022 & $[0.593,32.000]$ & $0$--$688$\\
Final 37.54 + GSW, CE336 & 0.8750 & $[0.059,10.375]$ & $289$--$688$\\
Eight proposals + interval check & 0.0016 & $[0.257,11.006]$ & $0$--$0$\\
\bottomrule
\end{tabular}
\end{table}

The eight-proposal path took a median of $0.0077$ seconds and used no
walk fallback. Its median paired time ratio to original GSW was
$0.00159$, about a $630$-fold speedup on this suite. Its discrepancy was
higher than GSW in 39 of 48 cases, equal in four, and lower in five;
the observed ranges were $[0.257,11.006]$ and $[0,3.438]$, respectively.
Thus the fastest observed option and the option with smaller typical
discrepancy are both useful. The input filter also avoided all walk moves
in 36 cases, but its all-positive completion reached discrepancy 32.

With the terminal rule fixed at CE336, the unfiltered final walk took
median paired times $1.52$ and $1.11$ times those of the quadratic and
below-67 walks. This comparison changes the profile and its row-energy
constraints together. The stronger theorem alone gives no practical
speed ranking. All matched runs stayed in the low branch, so their
results do not assess the normalized high-response improvement.

Eighteen additional fast-path checks at $n=2048$ and $4096$, with square
and wide matrices, all passed without fallback. They took $0.008$--$0.231$
seconds, including interval verification, and had discrepancy at most
$15.431$. These larger cases have no matched old-walk timing and are
excluded from the paired speed ratios. All timings are single-run
measurements; they do not establish an asymptotic exponent or a uniform
proposal-acceptance guarantee.


\subsection{Sources and verification}
The source of the final construction is
\nolinkurl{archive/komlos_constructive_final.tex} in this repository,
and its checks are under \nolinkurl{certificates/}.
The six certificate checks were rerun: two exact scalar certificates,
the protected-charge and randomized-screening checks, and the two
numerical response audits. The latter cover 300 inertia examples,
480 normalized-response identities, and 464 finite curves.
They support the written proofs and do not replace them.
The new implementation tests exercise the actual normalized proxy,
terminal decisions, outward row-sum enclosures, filtering, fixed
high-branch states, and explicit failure/fallback paths. Earlier
benchmark records remain unchanged; the final-profile comparison has
separate raw records, source hashes, and timing metadata.

\appendix
\section{Finite arithmetic primitives}\label{app:arithmetic}


\subsection{Finite arithmetic implementations and their scope}\label{b:sec:implementation}
\subsubsection{Direct rational arithmetic}

\begin{lemma}[Positive coefficient approximation]\label{b:lem:scalar}
Given a real arithmetic number $a\le0$ and $0<\epsilon<1$, one can
compute $\bar\alpha\in[0,32]$ with
$|\bar\alpha-32e^a|\le\epsilon$ in $O(\log(2/\epsilon))$
arithmetic operations and comparisons. Rational input gives rational output.
\end{lemma}
\begin{proof}
Find the smallest positive integer $H_0$ with $2^{H_0}\ge128/\epsilon$
by doubling. If $a\le-H_0$, return zero: the discarded value is at most
$32e^{-H_0}\le32\,2^{-H_0}\le\epsilon/4$.
Otherwise put $u=-a\in[0,H_0]$, compute the positive Taylor sum
$E_0=\sum_{k=0}^{6H_0}u^k/k!$, and return $32/E_0$.
Taylor's remainder divided by $e^u$ is at most
$u^{6H_0+1}/(6H_0+1)!\le2^{-(6H_0+1)}$, using
$k!\ge(k/e)^k$ and $e<3$. Since $E_0\ge1$, the returned value lies in
$[0,32]$, and inversion changes it by at most
$64\,2^{-(6H_0+1)}<\epsilon$. Successive terms cost constant arithmetic
work each. Taking $\epsilon=\delta/(1000(m+1))$ for each of at most $2m$
row-sign pairs gives the summed coefficient tolerance used above.
\end{proof}


For~\eqref{b:eq:diaginterface}, form an exactly orthogonal accumulator of
rational plane rotations. If $b$ is a largest off-diagonal entry in a
$2\times2$ block with diagonal entries $a,c$, use
$\cos\theta=(1-z^2)/(1+z^2)$, $\sin\theta=2z/(1+z^2)$.
The new pivot numerator is
$b(1-6z^2+z^4)+2z(1-z^2)(c-a)$.
Bisection between zero and
$-\tfrac12\operatorname{sgn}(b(c-a))$, using $1/2$ if $a=c$,
reduces the pivot magnitude by at least a factor two.
If $E$ is squared off-diagonal Frobenius norm, each rotation reduces it by
at least $3b^2/2\ge3E/(2s(s-1))$.

Stop once the squared off-diagonal Frobenius norm is at most $\xi^2$.
Before stopping, a largest pivot has $|b|\ge\xi/s$. On $|z|\le1/2$
the displayed quartic has derivative bounded by an absolute constant
when the input norm is at most two. Its values at the two bisection
endpoints have opposite signs: at the nonzero endpoint the contribution
from $b$ is $-7b/16$, and the other term has the opposite sign to $b$.
Consequently $O(\log(s/\xi))$ exact bisections find a rotation whose
new pivot is at most $|b|/2$; the denominator is $(1+z^2)^2\ge1$.

For norm-bounded matrices, $O(s^2\log(s/\xi))$ rotations suffice for
$E\le\xi^2$. A tournament tree finds the largest pivot and updates its
$O(s)$ changed entries in $O(s\log s)$ work. Each bisection costs
$O(\log(s/\xi))$. The diagonal part and exact orthogonal accumulator give
the interface in
$O(s^3\log s\log(s/\xi)+s^2\log^2(s/\xi))$ operations.
This also applies to the indefinite proxy matrix $Q$.

Extract independent computed constraint rows to obtain a full-row-rank
matrix $E_0$. The exact projector is
$P=I-E_0^\T(E_0E_0^\T)^{-1}E_0$; use $P=I$ for no rows.
Classical exact elimination and projection cost $O(s^3)$.
No conditioning bound is required in this arithmetic model.
All Gram/proxy matrices and cancellation maps cost $O(ms^2)$ to form.
Flat selection, including forming the row tests against its basis, costs
$O(ms^2+(m+s)s\log(2(m+s)s))$.
Thus the per-state bound is $\widetilde O(ms^2+s^3+m)$.
Scalar approximation precisions have logarithms $O(\log(2+m+n))$.

A phase beginning at size $S$ has $O(SJ^2q^2)$ updates.
Summing the geometrically decreasing values of $S$ proves
$\widetilde O(mn^3+n^4)$ arithmetic operations. Cleanup and scale
computations are absorbed; extra scans following deletions cost at most
$O(mn^2)$ overall. All computed quantities are rational for rational input.

\subsubsection{Fast real arithmetic with square roots}
Sobczyk~\cite[Corollaries 2.1 and C.2]{Sob} supplies~\eqref{b:eq:diaginterface}
deterministically in
$O(s^\omega\log s+s^2\operatorname{polylog}(s/\xi))$ operations,
including exact square roots. The witness orthogonal $U$ is used only in the
analysis. Our computed constraints use $V$ exactly; their leakage was proved
in Section~\ref{b:sec:subspace}.
Exact rank-profile algorithms such as~\cite{DPS}, followed by fast inversion
and products, construct $P$ in $\widetilde O(s^\omega)$ operations.
Projection of the approximate negative basis gives~\eqref{b:eq:basisgram}
without a conditioning assumption or exact eigenvectors.

The remaining large matrix constructions are rectangular products. If
$Z_0$ contains the row vectors $z_i$ and $G$ contains the row gradients,
examples are
$Z_0^\T\diag(\bar\alpha)Z_0$,
$G^\T\diag(\bar\alpha\beta^2(z_i^\T w)^2)G$, and
$Z_0^\T\diag(\bar\alpha\beta z_i^\T w)G$.
Blocking the $m$ rows into blocks of at most $s$ gives
$O(ms^{\omega-1}+s^\omega)$ work. Forming all flat-selection tests against
$T$ has the same bound; the subsequent scalar greedy procedure costs only
$O((m+s)s\log(2(m+s)s))$.
Both eigensolver calls, exact rank computation, projection, and all matrix
products therefore cost $\widetilde O(ms^{\omega-1}+s^\omega)$ per state.
The phase count proves
\[
 \widetilde O(mn^\omega+n^{\omega+1}).
\]
This establishes the fast clause of Theorem~\ref{b:thm:main}. The square-root
operation belongs to the imported real-RAM model; it is not silently included
in the rational-arithmetic clause.




\subsection{An imported Perron primitive with finite certificates}
\label{p:section}
This section imports the randomized Perron algorithm of
Ahmadinejad, Jambulapati, Saberi, and Sidford (AJSS)~\cite{AJSS}; it does not
supply a new eigensolver.  We give the implicit-matrix implementation and
an independently verifiable accuracy contract needed here.  The cost is
an expected arithmetic-operation bound, with the polylogarithmic factors
specified below.  It is separate from the cost of constructing a signing
direction.

\paragraph{Input representation and size.}
For $s\ge1$, the arithmetic matrix is supplied as
\begin{equation}\label{p:gram}
 \overline M=B+\sum_{i=1}^{r}\alpha_i z_i z_i^{\mathsf T}
                 +\eta\mathbf1\mathbf1^{\mathsf T},
 \qquad B\text{ diagonal and nonnegative},\quad
 \alpha_i\ge0,\quad z_i\ge0,\quad0<\eta\le1.
\end{equation}
Zero factors can be omitted.  Write
$N_0=s+\sum_i|\operatorname{supp}z_i|$; the factor supports are part of the
input, and a matrix-vector product costs $O(N_0)$.  When $r\le2m$,
$N_0=O((m+1)s)$, including $m=0$.  No dense Gram product is formed.
The signing problem with $m=0$ itself is handled by returning any signing
in $O(n)$ time; the primitive remains well-defined with an empty sum in
\eqref{p:gram}.  For $s=0$ there is no spectral computation, and for $s=1$
the scalar eigenpair is computed directly.

\paragraph{The imported solver and its applicable hypotheses.}
Formally, $\overline M=CC^{\mathsf T}$ with nonnegative columns
$\sqrt{\alpha_i}z_i$, $\sqrt{B_{jj}}e_j$, and $\sqrt\eta\mathbf1$.
AJSS Theorem~9 gives approximate Perron computation.  Section~6.3,
Lemmas~13--14, implements its nonnegative-Gram version using the factor
supports rather than the entries of $CC^{\mathsf T}$.  The symmetric
specialization of Theorem~8 and Appendix~10 uses symmetric diagonally
dominant (SDD) solvers.  It has polylogarithmic overhead; importing only
the paper's general asymmetric ``nearly linear'' statement would not
justify that stronger assertion.

Here is the reason this representation fits those results.  At each
symmetric scaling stage the SDD matrix to solve has form
$D_0-\sum_i a_i u_i u_i^{\mathsf T}$ with $u_i\ge0$, $a_i\ge0$.
It is a nonnegative diagonal remainder plus the weighted clique
Laplacians
\begin{equation}\label{p:cliques}
 a_i\bigl(\|u_i\|_1\operatorname{diag}(u_i)
                    -u_i u_i^{\mathsf T}\bigr).
\end{equation}
The diagonal remainder is nonnegative precisely because the matrix is
SDD.  Lemma~13 sparsifies each clique in its support size times a
polylogarithmic factor; Lemma~14 gives the resulting SDD solve.
Positive diagonal scaling preserves factor supports.  Their proof is a
sum over columns, so it also applies to rectangular $C$.  Thus the
symmetric Perron reduction uses
$N_0\operatorname{polylog}(s,1/\eta,1/\delta)$ operations at requested
inverse-polynomial accuracy $\delta$, provided the Perron-vector
condition number is polynomially bounded.  This is an imported
randomized algorithm, including its SDD subroutines.

The necessary condition bound is automatic here.  If
$\overline L=\lambda_1(\overline M)\le2$ and $v$ is the positive unit
Perron vector, then
\begin{equation}\label{p:ridgebounds}
 v_j\ge\frac{\eta\sum_kv_k}{\overline L}\ge\frac\eta2,
 \qquad \frac{\max_jv_j}{\min_jv_j}\le\frac2\eta,
 \qquad \overline L-\lambda_2(\overline M)\ge\eta.
\end{equation}
For the gap, $\overline M-\eta vv^{\mathsf T}$ is entrywise nonnegative,
has positive eigenvector $v$ with eigenvalue $\overline L-\eta$, and
leaves every other eigenvalue unchanged.  Perron--Frobenius gives the
last inequality.  This checks irreducibility, the vector condition bound,
and the gap without lower bounds on individual $\alpha_i$, $B_{jj}$,
or nonzero entries of $z_i$.  Also $\overline L\ge\eta s$ and
$\|\overline M\|_1=\|\overline M\|_\infty\le2\sqrt{s}$,
so the size and spectral-scale logarithms in the imported reduction are
covered by the same bound.

\begin{proposition}[Certified current Perron data]\label{p:certificate}
Let $M$ and $\overline M$ both have the form \eqref{p:gram}, with the same
ridge $\eta$, operator norms at most two, and
$\|M-\overline M\|\le\epsilon_M$.  For prescribed
$0<\epsilon_L,\epsilon_v\le1$, put
\begin{equation}\label{p:budgets}
 e=\min\{\epsilon_L/2,\eta\epsilon_v/16,\eta/4\},\qquad
 \epsilon_M\le\min\{\epsilon_L/8,\eta\epsilon_v/64\},\qquad
 \delta\le\frac{\eta^2e}{1024s}.
\end{equation}
There is a Las Vegas implementation of the imported primitive returning
an interval $I$ of width at most $\epsilon_L$ containing
$L=\lambda_1(M)$, and a vector $\widehat v$ satisfying
$\|\widehat v-v(M)\|\le\epsilon_v$.  Its expected cost is
\[
 N_0\operatorname{polylog}(2+s+1/\eta+1/\epsilon_L+1/\epsilon_v).
\]
The vector can be exactly normalized if square roots are available.
Otherwise, normalization to the asserted accuracy uses field operations
and comparisons only.
\end{proposition}
\begin{proof}
Run AJSS at accuracy $\delta$, rescale the candidate positive vector $r$
to $\|r\|_\infty=1$, and compute
\begin{equation}\label{p:collatz}
 \ell=\min_j\frac{(\overline Mr)_j}{r_j},\qquad
 u=\max_j\frac{(\overline Mr)_j}{r_j}.
\end{equation}
Reject if any coordinate is nonpositive or $u-\ell>e$.  The check costs
$O(N_0)$.  For every accepted output, Collatz--Wielandt and the matrix
error bound certify
$I=[\ell-\epsilon_M,u+\epsilon_M]$; its width is at most
$e+2\epsilon_M\le3\epsilon_L/4$.

Set $w=r/\|r\|_2$ and $q=w^{\mathsf T}\overline Mw$.  The numbers in
\eqref{p:collatz} bracket $q$ as well as $\overline L$, and their
componentwise bounds give
$\|(\overline M-qI)w\|\le e$.  By \eqref{p:ridgebounds},
$q-\lambda_2(\overline M)\ge\eta-e$.  Expanding in an orthonormal
eigenbasis, and using positivity to fix the Perron sign, yields
\[
 \|w-v(\overline M)\|
 \le\frac{\sqrt2e}{\eta-e}\le\frac{3e}{\eta}.
\]
The same elementary residual argument applied to the perturbation
$M-\overline M$, with Weyl's bound and $\epsilon_M\le\eta/64$, gives
$\|v(\overline M)-v(M)\|\le4\epsilon_M/\eta$.
Consequently $\|w-v(M)\|\le\epsilon_v/4$ under
\eqref{p:budgets}.  If exact square roots are unavailable, bisection on
$1\le\|r\|_2^2\le s$ gives a rational scaling of $r$ within
$\epsilon_v/4$ of $w$, with logarithmically many scalar operations.
This proves the claimed vector accuracy with slack.

It remains to bound retries.  On the imported algorithm's success event,
its scalar $\sigma$ satisfies $|\sigma-\overline L|\le2\delta$, and its
infinity-normalized vector satisfies
$\|(\overline M-\sigma I)r\|_\infty\le\delta$.
The gap and $\|r\|_2\ge1$ imply
$\|w-v(\overline M)\|\le3\sqrt{s}\delta/\eta\le\eta/4$.
Hence every coordinate of $r$ is at least $\eta/4$, and
$u-\ell\le8\delta/\eta<e$: the certificate passes.
Set the imported success probability to at least $3/4$, truncate each
attempt at its prescribed work bound, and restart with independent
randomness after failure.  Every returned answer is certified, and the
expected number of attempts is at most $4/3$.  This proves the asserted
expected cost.
\end{proof}

\paragraph{Arithmetic and coefficient precision.}
The square roots in the formal factor $C$ need not be evaluated:
\eqref{p:cliques} uses $\alpha_i$ directly.  The checks
\eqref{p:collatz} and the non-square-root normalization use only
$+,-,\times,/$ and comparisons, and preserve rationality when their
inputs are rational.  We charge the imported numerical routines in their
arithmetic model; this observation about our interface alone is not a
new rational-arithmetic implementation of every imported SDD routine.
For coefficient approximations, the explicit sufficient error budget is
\[
 \|B-\overline B\|+
 \sum_i|\alpha_i-\overline\alpha_i|\,\|z_i\|^2\le\epsilon_M.
\]
In particular $\|z_i\|\le1$ permits a sum-of-coefficient-errors budget.
The ridge is retained exactly and coefficients remain nonnegative.
Thus any separate exponential-evaluation routine meeting this finite
budget is compatible with Proposition~\ref{p:certificate}; no exact
exponential operation is assumed here.  If all requested error
reciprocals and $1/\eta$ are polynomial in $m+n$, the displayed cost is
$\widetilde O((m+1)s)$.

\begin{corollary}[Finite endpoint evaluation]\label{p:endpoint}
Suppose the supplied matrices at two points satisfy the hypotheses of
Proposition~\ref{p:certificate}, their true top eigenvalues lie in $[0,2]$,
and the quadratic energy at both points is available exactly.  For
$\Phi(x)=(\lambda_1(M(x))-1)_+^3-\eta\|x\|^2$, two calls with
$\epsilon_L\le\Delta/100$ certify a decrease of $\Delta$ whenever the
actual decrease is at least $2\Delta$.
\end{corollary}
\begin{proof}
Intersect each certified eigenvalue interval with $[0,2]$.
The increasing function $f(t)=(t-1)_+^3$ is $3$-Lipschitz there, so each
potential interval has width at most $3\epsilon_L$.  The certified
lower bound on the decrease differs from the true decrease by at most
$6\epsilon_L<\Delta$.
\end{proof}
This evaluation primitive neither constructs a direction nor proves a
positive acceptance probability for a proposed direction distribution.
In particular, its cost does not include a solve in a gradient-coupled
covariance operator.


\subsection{A certified power-iteration alternative}
\label{new:power}
The nearly linear Perron primitive can be replaced by ordinary power
iteration when dependence on the top spectral gap is acceptable.
This is an optional implementation, not an improvement of the existing
worst-case Perron budget. The convergence estimate is standard
\cite[Section~4.1]{SaadEig}; the stopping certificate is the same one used
in Proposition~\ref{p:certificate}.

Use the arithmetic matrix $\overline M$ and error budgets~\eqref{p:budgets}.
The scalar case is immediate. For $s\ge2$, put
\[
 \overline L=\lambda_1(\overline M),\qquad
 \kappa_{\rm P}=\frac{\overline L}
     {\overline L-\lambda_2(\overline M)},\qquad
 w_0=\one/\sqrt s,
 \quad w_{k+1}=\frac{\overline M w_k}{\|\overline M w_k\|}.
\]
Every iterate is positive. The matrix is PSD, and its positive unit
Perron vector $v$ satisfies $v^Tw_0\ge1/\sqrt s$. Spectral expansion gives
\[
 \tan\angle(w_k,v)\le\sqrt s\,
       (\lambda_2/\overline L)^k
 \le\sqrt s\,e^{-k/\kappa_{\rm P}}.
\]
At each iteration compute the Collatz bounds $\ell_k,u_k$ of
\eqref{p:collatz}, and stop when $u_k-\ell_k\le e$, with $e$ as
in~\eqref{p:budgets}. The proof of Proposition~\ref{p:certificate}
then certifies the scalar interval and Perron-vector accuracy for the
true $M$, including the error $\|M-\overline M\|$.

To bound termination, recall $v_i\ge\eta/2$ and $\overline L\le2$.
If $\|w_k-v\|\le d\le\eta/4$, then $(w_k)_i\ge\eta/4$ and
\[
 \left|\frac{(\overline Mw_k)_i}{(w_k)_i}-\overline L\right|
 \le\frac{8d}{\eta}.
\]
Thus $d\le\eta e/16$ suffices. The iteration terminates after
\[
 O\!\left(1+\kappa_{\rm P}
       \log\frac{16\sqrt{2s}}{\eta e}\right)
\]
matrix applications, each costing $O(N_0)$. It needs no prior knowledge
of the gap to run the certificate-stopped loop. Under a uniform bound
$\kappa_{\rm P}^*$ for every Perron call, the full walk's Perron work
becomes $\widetilde O((m+1)n^2\kappa_{\rm P}^*)$; add the selected
sampler cost. The present ridge only supplies
$\kappa_{\rm P}^*=\widetilde O(n)$, whereas the imported Perron method
already fits the cubic budget.

The inverse relative gap is different from the SPD condition number.
For example,
\[
 M_\eta=I_2+\eta
 \begin{pmatrix}2&1\\1&1\end{pmatrix}
 =\diag(1+\eta,1)+\eta\one\one^T
\]
has eigenvalues $1+(3\pm\sqrt5)\eta/2$. Its ordinary condition number
tends to one, while its inverse relative top gap is of order $1/\eta$.
Starting from $\one/\sqrt2$ has nonzero error in the second eigendirection,
so ordinary power iteration genuinely exhibits this gap dependence.
This is a check on the Perron primitive, not a claimed reachable state
of the signing walk.

\section{A lazy positive-Gram cache and its shifted inverse}\label{new:lazy}
\label{m:sec:cache}
This cache is an additional matrix-maintenance primitive. It provides a
preconditioner for positive shifts of $M$, and is not an assumed solver
for the filtered-ridge operator of the operator construction~\ref{r:operator}.
Within an epoch write the current matrix as
\begin{equation}
 M(x)=B+\eta\mathbf1\mathbf1^{\mathsf T}
                +\sum_i\alpha_i(x)z_iz_i^{\mathsf T},
 \qquad \alpha_i(x)=32e^{(2/7)u_i(x)}.                     \label{m:eq:cache-gram}
\end{equation}
Here $B\succeq0$ is diagonal, all $z_i\ge0$, $\|z_i\|_1\le1$, and
there are at most $2m$ summands. The profiles, $B$, and $\eta$ are fixed
inside the epoch. Along a selected direction the exact scalar identity is
\begin{equation}
 \Delta Y_i=a_id-b_id^2,\qquad Y_i=\log\alpha_i,
 \quad |a_i|\le K,\quad 0\le b_i\le K^2.                 \label{m:eq:log-increment}
\end{equation}
This follows directly from the quadratic barrier and uses no logarithm
computation in the algorithm.

Store an anchor $Y_i^*$ for each term, initialize it at the epoch's first
state, and put $r_0=1/8$. After a move refresh that coefficient if
$|Y_i-Y_i^*|\ge r_0$, resetting $Y_i^*=Y_i$.
The test is the finite comparison
$|(2/7)(u_i(x)-u_i(x_i^*))|\ge1/8$.
These scalar cache anchors are independent of the supporting-plane
anchors used for row feasibility.

\begin{lemma}[Expected number of coefficient refreshes]\label{m:lem:refresh}
For one fixed row-sign identity, let $R_i$ count its noninitial refreshes
over all large cleanup epochs in which it is present. Then
\begin{equation}
 \mathbb E R_i\le\frac2{r_0^2}\mathbb E\sum_ta_{i,t}^2d_t^2
                  +\frac3{r_0}\mathbb E\sum_tb_{i,t}d_t^2.
                                                               \label{m:eq:refresh-count}
\end{equation}
Consequently $\mathbb E\sum_iR_i=O(mJ^2q)$.
\end{lemma}
\begin{proof}
Before a move set $e=Y_i-Y_i^*$, so $|e|<r_0$.
Condition on the history and selected direction. By
\eqref{m:eq:coin} and \eqref{m:eq:log-increment},
$\mathbb E\Delta Y_i=-b_i\mathbb E d^2$, and
\[
 \mathbb E(\Delta Y_i)^2
 \le2a_i^2\mathbb E d^2+2b_i^2\mathbb E d^4
 \le2a_i^2\mathbb E d^2+2c_0^2b_i\mathbb E d^2.
\]
The last inequality uses $b_i d^2\le K^2\tau^2\le c_0^2$.
Thus the conditional expected increase of $e^2$ before resetting is at most
\[
 2a_i^2\mathbb E d^2+(2r_0+2c_0^2)b_i\mathbb E d^2
 \le2a_i^2\mathbb E d^2+3r_0b_i\mathbb E d^2.
\]
A refresh removes at least $r_0^2$ from this potential. Metadata resets or
completion of a row discard only a nonnegative residual and start any new
cache anchor at zero. Summing over the deterministically bounded number
of moves from Proposition~\ref{m:prop:coin}, and using the nonnegative
terminal potential, proves \eqref{m:eq:refresh-count}. Apply
$a_i^2,b_i\le K^2$, \eqref{m:eq:variation}, and the at most $2m$ row-sign
identities for the final bound.
\end{proof}

Let $M_*$ use the exact coefficients at the current cache anchors in
\eqref{m:eq:cache-gram}. After all refreshes have been processed,
positivity of every summand implies
\begin{equation}
                e^{-1/8}M_*\preceq M(x)\preceq e^{1/8}M_*.
                                                               \label{m:eq:cache-loewner}
\end{equation}
The fixed PSD part respects the same inequalities. This is a constant
relative spectral approximation; it makes no claim about fine eigenspace
accuracy.

\subsection{Finite coefficient accuracy and rank-one inverse updates}
\label{m:subsec:cache-arithmetic}
Fix a positive shift $\lambda$ during each epoch, for example $\lambda=\eta$.
Choose $0<\xi\le\min(1,\lambda/16)$. At initialization or refresh,
approximate an anchored $\alpha_i$ by a nonnegative arithmetic number
with absolute error at most $\xi/(2m+1)$. Then the numerical reference
$\widehat M_*$ has $\|\widehat M_*-M_*\|\le\xi$.
The cutoff and Taylor procedure of Section~\ref{b:sec:implementation}
evaluates these nonpositive exponential arguments using
$O(\log((2+m)/\xi))$ arithmetic operations per coefficient.
Small coefficients require absolute precision only.

Maintain the exact arithmetic inverse of
$A_* =\widehat M_*+\lambda I$. If a refreshed coefficient changes by $c$,
then
\begin{equation}
 A_{*,+}^{-1}=A_*^{-1}
 -\frac{c\,A_*^{-1}z_iz_i^{\mathsf T}A_*^{-1}}
               {1+c z_i^{\mathsf T}A_*^{-1}z_i}.           \label{m:eq:inverse-update}
\end{equation}
Both reference matrices are positive definite, because all stored
coefficients are nonnegative and $\lambda>0$. The denominator is therefore
strictly positive, including when $c<0$. Every refresh costs $O(S^2)$
operations. The maintained inverse is exact for the numerical reference;
no accumulation of inverse-rounding error is assumed in this arithmetic
model.

\begin{theorem}[Total positive-Gram cache cost]\label{m:thm:cachecost}
For $\lambda=\eta$ and inverse-polynomial numerical accuracy, the additional
expected cost of the reference matrix and its shifted inverse, including
coefficient evaluation, refresh tests, scans, and metadata rebuilds, is
\[
                         \widetilde O(mn^2+n^3).
\]
More generally the same bound has an additional logarithmic dependence on
$1/\lambda$ and the requested coefficient accuracy. The cache satisfies
\begin{equation}
        \tfrac12 A_*\preceq M(x)+\lambda I\preceq2A_*.
                                                               \label{m:eq:positive-preconditioner}
\end{equation}
\end{theorem}
\begin{proof}
Lemma~\ref{m:lem:refresh} bounds expected noninitial events by $O(mJ^2q)$;
\eqref{m:eq:inverse-update} charges at most $O(n^2)$ per event.
At each large epoch boundary form the reference matrix and its inverse in
$O(mS^2+S^3)$ work, paying for every changed profile, offset, scale, row
class, and ridge. Equation~\eqref{m:eq:epoch-sums} sums these rebuilds to
$O(mn^2+n^3)$. The fixed-dimension continuation contributes only
$\widetilde O(m+1)$, even with a rebuild after each move.
Scanning true barriers and cache residuals costs $O(mS)$ per move, hence
$\widetilde O(mn^2)$ by Proposition~\ref{m:prop:coin}. The stated scalar
precision factors are logarithmic. Finally, numerical error at most
$\lambda/16$ and \eqref{m:eq:cache-loewner} imply
\eqref{m:eq:positive-preconditioner}.
\end{proof}

For completeness this cache also gives a certified positive-shift solve.
Evaluate a current numerical matrix $A_{\rm cur}$ with the same fixed
shift and error at most $\lambda/16$. The constants above still give
$\tfrac12 A_*\preceq A_{\rm cur}\preceq2A_*$. Therefore
\[
 y_{j+1}=y_j+\tfrac12 A_*^{-1}(b-A_{\rm cur}y_j)
\]
contracts error in the $A_*$ norm by at most $3/4$ per iteration.
For relative energy error $0<\epsilon<1$ in the true system
$(M(x)+\lambda I)y=b$, use current coefficient error at most
$\lambda\epsilon/32$ and solve the numerical system to a sufficiently
smaller constant multiple of $\epsilon$. The perturbation estimate follows
from $M(x)+\lambda I\succeq\lambda I$. Its operation count is
\begin{equation}
 O\!\left(m\log\!\left(2+\frac{m+1}{\lambda\epsilon}\right)
 +(mS+S^2)\log(2/\epsilon)\right).                         \label{m:eq:positive-solve-cost}
\end{equation}
The full algorithm's remaining solve is the different structured operator
specified in the operator construction~\ref{r:operator}. Its cost is explicitly present in
\eqref{r:total}; neither \eqref{m:eq:positive-preconditioner} nor
\eqref{m:eq:positive-solve-cost} removes that term.

\begin{lemma}[A precise criterion for reusing the remaining preconditioner]
\label{m:reuse}
Fix a row-feasible space $H$ and an orthonormal basis matrix $U$ for it.
For positive diagonal $F_t$ and a penalty operator $E_t$, write
\[
 V_t=\begin{pmatrix}F_t^{1/2}U\\ \sqrt{\rho_{{\rm ridge},t}}\,E_tU\end{pmatrix},
 \qquad A_t=V_t^{\mathsf T}V_t\succ0.
\]
If, for some reference state $0$ and $0\le r<1$,
\[
                \|(V_t-V_0)A_0^{-1/2}\|\le r,
\]
then
\[
              (1-r)^2A_0\preceq A_t\preceq(1+r)^2A_0.
\]
Thus a fixed $r<1$ makes a cached inverse of $A_0$ a constant-quality
preconditioner for $A_t$.
\end{lemma}
\begin{proof}
The columns of $V_0A_0^{-1/2}$ are orthonormal. For every $y$, the triangle
and reverse triangle inequalities give
\[
 (1-r)\|y\|\le\|V_tA_0^{-1/2}y\|\le(1+r)\|y\|.
\]
Squaring and conjugating by $A_0^{1/2}$ proves the two Loewner inequalities.
\end{proof}
The lemma fixes the row-feasible space and measures the augmented factor
in the reference inverse metric. A change of that space requires a
separate update argument. Neither checking this operator-norm condition
nor bounding its cumulative violations is supplied by the positive-Gram
cache theorem. The missing total update bound is exactly what prevents
this reuse criterion from establishing a complete cubic algorithm.


\section{A certified allowance for additional regularization}
\label{reg:section}

The direction construction tolerates extra positive semidefinite
regularization.  Its cost to the curvature argument is an effective
dimension, rather than the size of its largest eigenvalue.  This is a
local robustness statement: selecting and maintaining a useful
regularizer remains an algorithmic task, whose work must be charged.

\begin{lemma}[Curvature loss under an added regularizer]
\label{reg:loss}
Let $A\succ0$, $R\succeq0$, and $Q=Q^{\mathsf T}\preceq A$ act on a
finite-dimensional Euclidean space.  Set $X_R=(A+R)^{-1}$ and
\[
 d_R=\operatorname{tr}(RX_R).
\]
Then
\begin{equation}\label{reg:loss-bound}
 \operatorname{tr}(QX_R)
 \ge \operatorname{tr}(QA^{-1})-d_R,
 \qquad 0\le d_R\le\operatorname{rank}R.
\end{equation}
If $R=VV^{\mathsf T}$ and $L=V^{\mathsf T}A^{-1}V$, then
\begin{equation}\label{reg:effective-dimension}
 d_R=\operatorname{tr}\bigl(L(I+L)^{-1}\bigr).
\end{equation}
No positivity assumption on $Q$ is needed.
\end{lemma}
\begin{proof}
The difference $D=A^{-1}-(A+R)^{-1}$ is positive semidefinite.  Therefore
\[
 \operatorname{tr}(QD)\le\operatorname{tr}(AD)
 =\operatorname{tr}\bigl(I-A(A+R)^{-1}\bigr)=d_R,
\]
which proves the first assertion.  Put $W=A^{-1/2}RA^{-1/2}$.  Cyclicity
of trace gives $d_R=\operatorname{tr}(W(I+W)^{-1})$; each positive
eigenvalue $w$ contributes $w/(1+w)\le1$.  Finally, with
$U=A^{-1/2}V$, the identity $(I+UU^{\mathsf T})^{-1}U=U(I+U^{\mathsf T}U)^{-1}$
proves~\eqref{reg:effective-dimension}.
\end{proof}

\subsection{A certified sampler with additional regularization}

We apply the lemma on the current row-feasible space $H_0=\operatorname{ran}P$.
Write $s$ for the retained coordinate dimension and $C=1200$.  The
current finite arithmetic proxies satisfy
\[
 \eta I\preceq F\preceq\mu I,\qquad
 Q=P(F-S)P\big|_{H_0},\qquad
 \|F-S\|\le3,
 \qquad
 A=P(F+\rho_{\mathrm{ridge}}E^{\mathsf T}E)P\big|_{H_0}.
\]
The exact diagonal comparison gives $F\preceq D+\eta I$.
Let $V\in\mathbb R^{s\times r}$, $0\le r\le s$, be fixed before the
sampling trial, and put
\[
 R=PVV^{\mathsf T}P\big|_{H_0},\qquad A_R=A+R.
\]
All matrices and the representation of $V$ remain fixed during sampling
retries at this state.  Conditional on any randomness used to select
$V$, the signs below are fresh and independent.

\begin{theorem}[The original direction contracts survive extra regularization]
\label{reg:sampler}
In the high branch, tighten the filter-tail request to
\begin{equation}\label{reg:budgets}
 (2\varepsilon_f+\varepsilon_E)^2
       \le\frac{\eta}{4800\rho_{\mathrm{ridge}}},
 \qquad d_R\le\frac{s}{4800}.
\end{equation}
Keep all other geometric parameters and finite accuracy requests from
Sections~\ref{r:regularized}--\ref{r:parameters}.  In particular, let
$b_1,\ldots,b_N$ be the normalized tests with $\|b_j\|\le1$, choose
$2^q\ge16CN$, and set
\[
 \zeta_0=\min\{\eta/(1000C),1/(100s)\}.
\]
For independent random signs
$\xi,\omega\in\{-1,1\}^s$ and $\nu\in\{-1,1\}^r$, form
\begin{equation}\label{reg:rhs}
 b=P\bigl(F^{1/2}\xi+
       \sqrt{\rho_{\mathrm{ridge}}}E^{\mathsf T}\omega+V\nu\bigr).
\end{equation}
Compute $y\in H_0$ satisfying the absolute residual certificate
\begin{equation}\label{reg:residual}
 \|b-A_Ry\|^2\le\zeta_0^2\eta s.
\end{equation}
Accept exactly when
\begin{equation}\label{reg:accept}
 y^{\mathsf T}(F-S)y\ge\frac{s}{8C},\qquad
 |b_j^{\mathsf T}y|^2\le\frac{4q}{\eta}\quad(1\le j\le N).
\end{equation}
Upon acceptance, choose a positive power of two $a$ with
$1/4\le a^2\|y\|^2\le1$ and return $h=ay$.

The acceptance probability is at least $1/(4C)$.  Every accepted output
satisfies
\begin{gather}
 h\in H_0,\qquad \tfrac12\le\|h\|\le1,
 \qquad h^{\mathsf T}(F-S)h\ge\frac{\eta}{128C},
 \label{reg:curvature}\\
 |b_j^{\mathsf T}h|^2\le\frac{32C\chi q}{s},\qquad
 \rho_{\mathrm{ridge}}\|Eh\|^2+\|V^{\mathsf T}h\|^2
       \le32C\mu,
 \qquad \chi=\mu/\eta.
 \label{reg:contracts}
\end{gather}
Consequently the original true-curvature and true-high-space leakage
contracts hold with the original $K,\tau,e_*$ and ridge penalty.
\end{theorem}
\begin{proof}
The raw dimension surplus in Lemma~\ref{m:lem:dimension} is at least
$s/600$.  Lemma~\ref{r:trace} and the stronger filter request give
$\operatorname{tr}(QA^{-1})\ge s/600-s/4800$.
Lemma~\ref{reg:loss} therefore yields
\begin{equation}\label{reg:trace-margin}
 \operatorname{tr}(QA_R^{-1})\ge
 \frac{s}{600}-\frac{s}{4800}-\frac{s}{4800}
 =\frac{s}{800}=\frac{3s}{2C}.
\end{equation}

On $H_0$, let
$B=P[F^{1/2},\sqrt{\rho_{\mathrm{ridge}}}E^{\mathsf T},V]$.
Then $BB^{\mathsf T}=A_R$ and $B^{\mathsf T}A_R^{-1}B$ is an orthogonal
projector.  For the exact solution $y_*=A_R^{-1}b$, extended by zero on
$H_0^\perp$, its covariance is exactly $X_R=PA_R^{-1}P$ and, pathwise,
\begin{equation}\label{reg:sample-energy}
 y_*^{\mathsf T}A_Ry_*\le2s+r\le3s.
\end{equation}
Because $Q\preceq A_R$, the random variable
$Z_*=y_*^{\mathsf T}(F-S)y_*$ is at most $3s$ and has expectation at least
$3s/(2C)$.  Hence
\[
 \Pr\{Z_*\ge s/(4C)\}
 \ge\frac{3s/(2C)-s/(4C)}{3s-s/(4C)}
 \ge\frac5{12C}.
\]
The compressed inverse obeys
$0\preceq X_R\preceq PA^{-1}P\preceq\eta^{-1}I$.
Each $b_j^{\mathsf T}y_*$ is therefore a Rademacher sum with squared
coefficient norm at most $1/\eta$.  The probability that any test has
$|b_j^{\mathsf T}y_*|^2>2q/\eta$ is at most
$2Ne^{-q}\le2N2^{-q}\le1/(8C)$.

Write $e=y-y_*$.  The absolute residual implies
\[
 \|e\|_{A_R}\le\zeta_0\sqrt s,
 \qquad \|e\|\le\zeta_0\sqrt{s/\eta}.
\]
Using $\|F-S\|\le3$ and~\eqref{reg:sample-energy},
\[
 \bigl|y^{\mathsf T}(F-S)y-Z_*\bigr|
 \le3(2\sqrt3\,\zeta_0+\zeta_0^2)s/\eta
 \le14\zeta_0s/\eta<s/(8C).
\]
Also $|b_j^{\mathsf T}e|\le1/(100\sqrt\eta)$, so the two good
exact-sample events imply all tests in~\eqref{reg:accept}.
Their intersection has probability at least
$5/(12C)-1/(8C)=7/(24C)>1/(4C)$.

For every accepted sample, independently of those probability estimates,
$y^{\mathsf T}Fy\ge s/(8C)$ gives
$\|y\|^2\ge s/(8C\mu)$, hence $a^2\le8C\mu/s$.
Moreover $\zeta_0\le1/100$ and
\[
 \|y\|_{A_R}\le(\sqrt3+\zeta_0)\sqrt s<2\sqrt s.
\]
Thus $\|y\|^2\le4s/\eta$ and
$a^2\ge1/(4\|y\|^2)\ge\eta/(16s)$.
These inequalities prove~\eqref{reg:curvature} and the flatness bound.
Finally
\[
 \rho_{\mathrm{ridge}}\|Ey\|^2+\|V^{\mathsf T}y\|^2
 \le y^{\mathsf T}A_Ry\le4s,
\]
which proves the combined penalty bound in~\eqref{reg:contracts}.

For completeness, let $\mathcal C$ denote the true curvature proxy and
$E_*=f(M/\Lambda)T$ the ideal filter factor.  The unchanged accuracy
$\|\mathcal C-(S-D)\|\le\eta/(256C)$ and $F\preceq D+\eta I$ give
\[
 h^{\mathsf T}\mathcal C h
 \le\eta\|h\|^2-\eta/(128C)+\eta/(256C)
 <\eta\|h\|^2.
\]
The unchanged penalty and factor-accuracy requests,
$\rho_{\mathrm{ridge}}\ge512C\mu/e_*^2$ and
$\varepsilon_E\le e_*/4$, imply
$\|Eh\|\le e_*/4$ and $\|E_*h\|\le e_*/2$.
Equation~\eqref{r:true-leakage} then gives
$\|P_HM'[h]v\|\le e_*$ on the true high space.
The coordinate and gradient bounds yield the same $K$ as before.
The original finite-step proof consequently applies with the same
$\tau$ and endpoint choices.
\end{proof}

The added noise $V\nu$ is essential.  Keeping the old right-hand side
while replacing its solve by $A_R^{-1}$ would give covariance
$A_R^{-1}AA_R^{-1}$, to which Lemma~\ref{reg:loss} does not apply.
Already for $A=Q=1$ and $R=t>0$, this would lose
$1-(1+t)^{-2}>t/(1+t)=d_R$ units of curvature.

\paragraph{Finite arithmetic and costs.}
The theorem uses the supplied finite arithmetic factor $V$ itself; any
cost of constructing, compressing, or updating that factor is additional.
For a dense $s\times r$ factor, its extra right-hand-side cost and its
extra cost per system application are $O(sr)$, besides projection onto
$H_0$.  A rank certificate $r\le s/4800$ is sufficient for the effective
dimension request, without evaluating a trace or a determinant.  A
higher-rank factor can also qualify, but certifying its effective
dimension is a separate task.

Since the ridge penalty is a power of four, its square root is rational.
Approximate only the diagonal roots of $F$ to absolute accuracy
$\zeta_0\eta/8$, as in the original sampler.  This produces
$\widetilde b$ with
$\|\widetilde b-b\|\le\zeta_0\sqrt{\eta s}/8$.  It suffices to check
\[
 \|\widetilde b-A_Ry\|^2\le\zeta_0^2\eta s/4;
\]
the triangle inequality implies~\eqref{reg:residual}.  All matrix-vector
products and this residual must use the same fixed finite operator,
including the same $V$.  Arbitrarily accurate residuals exist because
$A_R\succeq\eta I$; their computational cost is not presumed small.
In particular, an unbounded $\|R\|$ does not preserve the old polynomial
condition-number bound: the elementary bound becomes
\[
 \kappa(A_R)\le
 \frac{\mu+9\rho_{\mathrm{ridge}}+\|R\|}{\eta}.
\]

\paragraph{Scope for the signing walk.}
An optional variant may use any predictable supplied regularizer meeting
these conditions in a high state, while keeping $R=0$ in low states.
Fresh sampler signs and the original mean-zero endpoint coin preserve
all previously proved local contracts.  In particular, $x^{\mathsf T}h=0$
still gives the pathwise identity
$\|x+dh\|^2-\|x\|^2=d^2\|h\|^2$, so the accepted-move bound and final
discrepancy are unchanged.  A constant expected number of trials per
state also survives.  This statement does not choose a regularizer,
bound its maintenance work, or provide a new cubic algorithm.  Its
sampling law differs from that of the fixed controller in
Section~\ref{d:controller}; that controller and its stated conjecture
remain unchanged.

\subsection{Regularization, effective dimension, and determinant scores}
\label{reg:dpp}

The admissible amount of regularization is measured by effective dimension,
not by its operator norm.  This observation has an exact determinantal
interpretation, but the interpretation does not provide a bound on the cost
of updating the regularizer.

Let $A\succ0$ and $R=VV^\top\succeq0$, and write
\[
 d_R=\operatorname{tr}\bigl(R(A+R)^{-1}\bigr),\qquad
 L=V^\top A^{-1}V.
\]
The effective-dimension bound~\eqref{reg:effective-dimension} charges
regularization by $d_R$.  Woodbury's identity gives
$d_R=\operatorname{tr}(L(I+L)^{-1})$ with no commutativity assumption.

Write $V=[v_1,\ldots,v_k]$.  The ordinary $L$-ensemble is
\begin{equation}\label{reg:lensemble}
 \mathbb P\{\mathcal J=J\}=
 \frac{\det L_{J,J}}{\det(I+L)},\qquad J\subseteq\{1,\ldots,k\}.
\end{equation}
The empty determinant equals one.  The principal-minor expansion of
$\det(I+L)$ proves normalization.  Its marginal kernel and expected size
are
\begin{equation}\label{reg:marginal}
 K=L(I+L)^{-1}=V^\top(A+R)^{-1}V,\qquad
 \mathbb P\{j\in\mathcal J\}=K_{jj},\qquad
 \mathbb E|\mathcal J|=d_R.
\end{equation}
These are the usual ridge-leverage marginals of a determinantal process;
the additional information is in the joint law, not a new collection of
individual scores~\cite{DMdpp}.  For example,
$\mathbb P\{i,j\in\mathcal J\}=K_{ii}K_{jj}-K_{ij}^2$.
The role of the same effective dimension in Bayesian experimental design
is developed in~\cite{DLMdesign}.  The regularized experimental-design law
in that paper includes a Bernoulli reference measure; it should not be
identified with~\eqref{reg:lensemble} without the corresponding change of
kernel.

There is also an exact hard-constraint interpretation.
For each positive-probability $J$, define
\[
 C_J=A^{-1}-A^{-1}V_J(L_{J,J})^{-1}V_J^\top A^{-1},
 \qquad C_\varnothing=A^{-1}.
\]
Then $C_J\succeq0$, $V_J^\top C_J=0$, and
\begin{equation}\label{reg:averageconstraint}
 \mathbb E C_{\mathcal J}=(A+R)^{-1}.
\end{equation}
To prove this identity, differentiate
\[
 \log\!\sum_J\det(V_J^\top A^{-1}V_J)
 =\log\det(A+VV^\top)-\log\det A
\]
with respect to $A$ in an arbitrary symmetric direction.  On the left,
the derivative of each positive term divided by that term is the negative
trace pairing with $A^{-1}V_J(L_{J,J})^{-1}V_J^\top A^{-1}$.
Rank-deficient subsets have identically zero determinant and contribute
nothing.  Equality of the trace pairings proves~\eqref{reg:averageconstraint}.
Equivalently, in coordinates whitened by $A$, an ordinary determinantal
sample selects a random subspace whose mean orthogonal projector is
$A^{-1/2}V(I+L)^{-1}V^\top A^{-1/2}$.
This is the standard implicit-regularization identity, specialized to
consistent regression data; see~\cite[Theorem~3]{DMdpp}.
The identity gives a covariance interpretation, not an implementation of
the direction sampler with all its acceptance guarantees.

The determinant also measures the charge differentially.  For
$\phi(t)=\log\det(A+tR)-\log\det A$,
\[
 t\phi'(t)=d_{tR},\qquad
 \phi''(t)=-\operatorname{tr}\bigl(((A+tR)^{-1}R)^2\bigr)\leq0.
\]
For positive weights $w_j$ and $M=A+\sum_jw_jv_jv_j^\top$,
\[
 \frac{\partial\log\det M}{\partial\log w_j}
 =w_jv_j^\top M^{-1}v_j,\qquad
 \frac{\partial^2\log\det M}{\partial w_i\partial w_j}
 =-(v_i^\top M^{-1}v_j)^2.
\]
The first quantity is the corresponding weighted $L$-ensemble inclusion
probability.  Thus determinant gradients recover the same exact charge.

\paragraph{What the interpretation does not imply.}
An expected cardinality is not a pathwise rank bound.  Conditioning on
$|\mathcal J|=k$ changes the marginal formula: for
$L=\operatorname{diag}(1,3)$, the ordinary marginals are $(1/2,3/4)$,
whereas conditioning on one element gives $(1/4,3/4)$.
Nor does the identity assert that the unweighted partial sum
$V_{\mathcal J}V_{\mathcal J}^\top$ spectrally approximates $VV^\top$.
The individual scores and subset law depend on the chosen column
factorization, although $d_R$ depends only on $A$ and $R$.

Even the best coupling of consecutive determinantal samples need not
have small cost.  Let $\Pi_1,\Pi_2,\Pi_3$ project onto disjoint coordinate
sets of size $r$, and put $L_i=9\Pi_i$.
The matrices $\alpha(I+L_i)$ all have condition number $10$ and the same
determinant, while every law has expected size $9r/10$.
A sample from $L_i$ is supported on its own coordinate set.  Hence, for
any coupling of samples $\mathcal J_i,\mathcal J_j$ with $i\ne j$,
\begin{equation}\label{reg:dpprecourse}
 \mathbb E|\mathcal J_i\mathbin\triangle\mathcal J_j|
 =\mathbb E|\mathcal J_i|+\mathbb E|\mathcal J_j|=\frac{9r}{5}.
\end{equation}
The equality is pointwise before taking expectations because the sets
are disjoint.  This concerns samples on a fixed coordinate ground set.
It is not a lower bound for arbitrary matrix solvers or a counterexample
arising along the signing walk.  It shows why static determinant and
effective-dimension identities alone do not bound the work of following
a changing subspace.


\subsection{A complete accounting result for a fixed reference matrix}

Let $A\succ0$ be an $s\times s$ matrix and $R\succeq0$.  Define
\[
 d_A(R)=\operatorname{tr}\bigl(R(A+R)^{-1}\bigr),\qquad
 g_A(R)=\log\det(A+R)-\log\det A.
\]
If $\lambda_1,\ldots,\lambda_s\ge0$ are the eigenvalues of
$A^{-1/2}RA^{-1/2}$, then
\[
 d_A(R)=\sum_j\frac{\lambda_j}{1+\lambda_j},\qquad
 g_A(R)=\sum_j\log(1+\lambda_j).
\]
In particular, $d_A(R)\le g_A(R)$.  If $A+R\preceq\kappa A$,
with $\kappa\ge1$, then
\begin{equation}\label{regmaint:detdim}
 g_A(R)\le (1+\log\kappa)d_A(R).
\end{equation}
Indeed, $\log(1+x)\le x$ gives
$\log(1+x)\le [x/(1+x)]\{1+\log(1+x)\}$, and
$1+x\le\kappa$.  The case $x=0$ follows by continuity.

\begin{proposition}[Monotone fixed-reference updates]
\label{regmaint:fixed}
Suppose $R_0\succeq0$ and
$R_t=R_{t-1}+v_tv_t^\top$.  Put
\[
 \ell_t=v_t^\top(A+R_{t-1})^{-1}v_t.
\]
If every accepted update has $\ell_t\ge\epsilon>0$, and at the
last update $d_A(R_T)\le b$ and $A+R_T\preceq\kappa A$, then
\begin{equation}\label{regmaint:count}
 T\le\frac{b(1+\log\kappa)}{\log(1+\epsilon)}.
\end{equation}
Given an initial inverse and candidate vectors, each score evaluation
and each accepted inverse update requires $O(s^2)$ real-arithmetic
operations.  Thus $N$ candidate evaluations and the accepted updates
cost $O(Ns^2+Ts^2)$ after one $O(s^3)$ initialization.  This count does
not include constructing candidate vectors or finding a violating
direction.
\end{proposition}
\begin{proof}
The matrix determinant lemma gives
\[
 g_A(R_t)-g_A(R_{t-1})=\log(1+\ell_t).
\]
Telescope, discard $g_A(R_0)\ge0$, and use
\eqref{regmaint:detdim}.  With $H=A+R_{t-1}$, the maintained inverse
is updated by the exact Sherman--Morrison formula
\[
 (H+v_tv_t^\top)^{-1}
 =H^{-1}-\frac{H^{-1}v_tv_t^\top H^{-1}}{1+\ell_t}.
\]
No determinant needs to be computed by the algorithm.
\end{proof}

The upper comparison with $\kappa A$ is necessary for this proof and
cannot be omitted from its statement.  For $A=I$ and a rank-one
orthogonal projector $\Pi$, take $R_t=(2^t-1)\Pi$.  Every update has
$\ell_t=1$, while
\[
 d_A(R_t)=1-2^{-t}<1,\qquad g_A(R_t)=t\log2.
\]
Thus an effective-dimension allowance of one permits arbitrarily many
constant determinant-gain updates in a single direction.  The
determinantal gain measures both the number of affected directions and
their accumulated strength; effective dimension saturates in each
direction.

For a changing reference matrix $A_t\succ0$, the same calculation has
an additional exact term.  If $R_t=R_{t-1}+v_tv_t^\top$ and
$\ell_t=v_t^\top(A_t+R_{t-1})^{-1}v_t$, put
\[
 \delta_t=g_{A_t}(R_{t-1})-g_{A_{t-1}}(R_{t-1}).
\]
Then
\begin{equation}\label{regmaint:dynamic}
 \sum_{t=1}^T\log(1+\ell_t)
 =g_{A_T}(R_T)-g_{A_0}(R_0)-\sum_{t=1}^T\delta_t.
\end{equation}
For example, $A_t\succeq A_{t-1}$ implies $\delta_t\le0$:
differentiate $g_A(R)$ in a positive semidefinite direction and use
$(A+R)^{-1}\preceq A^{-1}$.  Consequently the moving-reference term
cannot simply be discarded with the sign needed for an upper bound.
Equation~\eqref{regmaint:dynamic} is an accounting identity, not a
bound for the signing walk.  The latter would need a separately proved
bound for its actual reference-matrix motion, as well as a rule that
keeps the effective-dimension allowance and supplies candidates cheaply.

\subsection{The generic cache obstruction survives selective regularization}

Fix $s\ge16$, $r=\lfloor s/16\rfloor$, and $\alpha>0$.  Let
$\Pi_1,\Pi_2,\Pi_3$ project onto three mutually orthogonal
$r$-dimensional subspaces, and set
\[
 K_i=\alpha(I+9\Pi_i).
\]
For $b\ge0$, allow an arbitrary added regularizer at every state and
define the enlarged good-cache class
\[
 \mathcal N_i(b)=\left\{B\succ0:
 \begin{array}{l}
 \text{there exists }R\succeq0\text{ with }d_{K_i}(R)\le b,\\[-2pt]
 (4/7)(K_i+R)\preceq B\preceq4(K_i+R)
 \end{array}\right\}.
\]
The regularizer may have full rank and arbitrary magnitude.  It may be
chosen after seeing the target and the cache.

\begin{proposition}[Rank separation with an effective-dimension budget]
\label{regmaint:separation}
If $i\ne j$, $B_i\in\mathcal N_i(b)$, and
$B_j\in\mathcal N_j(b)$, then
\begin{equation}\label{regmaint:rankgap}
 \operatorname{rank}(B_i-B_j)
 \ge2\max\{0,r-\lfloor10b/3\rfloor\}.
\end{equation}
For two classes, the threshold for a common cache is sharp:
\[
 \mathcal N_i(b)\cap\mathcal N_j(b)\ne\varnothing
 \quad\Longleftrightarrow\quad b\ge3r/10.
\]
\end{proposition}
\begin{proof}
Choose witnessing regularizers $R_i,R_j$, and let
$U_j=K_j^{-1/2}R_jK_j^{-1/2}$.  At most
$k=\lfloor10b/3\rfloor$ eigenvalues of $U_j$ are at least $3/7$,
because each contributes at least $3/10$ to $d_{K_j}(R_j)$.
Intersect $\operatorname{ran}\Pi_i$ with the span of the remaining
eigenvectors of $U_j$.  This gives a subspace $W$ of dimension at least
$r-k$.  On $\operatorname{ran}\Pi_i$ we have $K_j=\alpha I$, so for
every nonzero $x\in W$,
\[
 x^\top R_jx<\frac{3\alpha}{7}\|x\|^2.
\]
Therefore
\[
 x^\top B_ix\ge\frac{40\alpha}{7}\|x\|^2,
 \qquad
 x^\top B_jx<4\left(1+\frac37\right)\alpha\|x\|^2
 =\frac{40\alpha}{7}\|x\|^2.
\]
Thus $B_i-B_j$ has at least $r-k$ positive eigenvalues.  Exchanging
$i$ and $j$ gives at least $r-k$ negative eigenvalues, proving
\eqref{regmaint:rankgap}.

For necessity of the intersection threshold, suppose the same $B$
belongs to both classes.  On $\operatorname{ran}\Pi_i$, its lower
bound from state $i$ and upper bound from state $j$ imply
$x^\top R_jx\ge(3\alpha/7)\|x\|^2$.  The min--max principle gives
at least $r$ eigenvalues of $U_j$ at least $3/7$, whence
$d_{K_j}(R_j)\ge3r/10$.  Conversely, let
\[
 R_i=\frac{3\alpha}{7}\Pi_j,\quad
 R_j=\frac{3\alpha}{7}\Pi_i,\quad
 B=\frac{40\alpha}{7}(\Pi_i+\Pi_j)
    +\alpha(I-\Pi_i-\Pi_j).
\]
Both effective dimensions equal $3r/10$, and direct restriction to
the three orthogonal summands verifies both good-cache inequalities.
\end{proof}

Let $D=r-\lfloor10b/3\rfloor>0$.  Rank distance between two
different classes is at least $2D$.  By the rank triangle inequality,
a fixed cache has rank distance strictly less than $D$ from at most
one of the three classes.  With two caches an adaptive sequence can
therefore choose a class at rank distance at least $D$ from both
current caches.  A rank-two repair changes this distance by at most
two, so reaching either good-cache class needs at least
$\lceil D/2\rceil$ repairs.  Taking $b=s/2400$ gives
$D=\lfloor s/16\rfloor-\lfloor s/720\rfloor=\Omega(s)$; over
$s$ such selections the cost is $\Omega(s^2)$ rank-two repairs.
Here $b=s/2400$ is the conservative allowance obtained by preserving
half of the already guaranteed curvature trace $s/1200$.
The same lower bound therefore also applies under the smaller
allowance $s/4800$ used in Theorem~\ref{reg:sampler}.  A particular
state can have a larger actual curvature trace and could permit a larger
regularizer; the result does not restrict that additional margin.

This statement allows either cache its own optimally chosen
regularizer, so it also applies to any restricted choice within these
enlarged classes.  It concerns controllers whose cache changes are
charged as rank-two repairs.  It is not a lower bound for arbitrary
linear solvers: the model matrices are explicitly simple, and even
their unregularized condition numbers equal ten.

As in the unregularized model, choosing $\alpha=s^{-3}$ and
$E_i=3\sqrt\alpha\,\Pi_i$ gives
\[
 K_i=\alpha I+E_i^\top E_i,\qquad
 \sum_{t=0}^{s-1}\|E_{i_{t+1}}-E_{i_t}\|_F^2
 =18rs\alpha\le\frac{9}{8s}
\]
for any $s$ switches between distinct classes.  Hence the vanishing
ordinary factor-energy budget and the full selective-regularization
allowance together still do not imply a near-linear repair count for
this generic controller.  The missing input must use further structure
of the actual walk, or a different solver whose complexity is proved
without this repair count.

\section{Schur identities and fixed-profile maintenance}\label{new:schur}
\subsection{A rank-one update pays only for its residual overlap}
Write
\[
 L=\begin{pmatrix} A&B\\ B^T&C\end{pmatrix}\succ0,
 \qquad S=C-B^TA^{-1}B\succ0.
\]
For a vector $z=(v,u)$ and a scalar $\delta\ge0$, put
\[
 \ell=v^TA^{-1}v,\qquad r=u-B^TA^{-1}v.
\]
\begin{lemma}[Exact overlap discount]\label{ifc:lem:update}
The Schur complement $S_+$ of the leading block of
$L_+=L+\delta zz^T$ is
\begin{equation}\label{ifc:eq:update}
 S_+=S+\frac{\delta}{1+\delta\ell}rr^T.
\end{equation}
Moreover,
\begin{equation}\label{ifc:eq:detshare}
 \frac{\det L_+}{\det L}
 =(1+\delta\ell)
 \left(1+\frac{\delta}{1+\delta\ell}r^TS^{-1}r\right).
\end{equation}
The formulas remain valid for negative $\delta$ when $L_+\succ0$.
\end{lemma}
\begin{proof}
Use
$(A+\delta vv^T)^{-1}
=A^{-1}-\delta A^{-1}vv^TA^{-1}/(1+\delta\ell)$
in the definition of the new Schur complement and collect terms.
Positive definiteness ensures $1+\delta\ell>0$. The determinant
factorization $\det L=\det A\det S$ and the rank-one determinant
lemma then give \eqref{ifc:eq:detshare}.
\end{proof}
This is the useful form of overlap cancellation: the new direction is
the residual $r$, and its strength is discounted by $1+\delta\ell$.
If $r=0$, it makes no change at all to the Schur complement, regardless
of its size in the full space. Applying \eqref{ifc:eq:update} separately
to every scalar coefficient at every step is not yet a cubic algorithm:
there may be linearly many changing coefficients per step. The next
section avoids that implementation cost.

\subsection{A cubic theorem with arbitrary fixed overlaps}
Fix feature matrices $V\in\R^{q\times p}$ and $U\in\R^{r\times p}$,
and fixed $H\succeq0$, $\alpha>0$. Let
$D_t=\diag(a_1(t),\ldots,a_p(t))\succeq0$ and define
\begin{align}
 L_t&=\begin{pmatrix}
 H+VD_tV^T&VD_tU^T\\
 UD_tV^T&\alpha I+UD_tU^T
 \end{pmatrix},\label{ifc:eq:lift}\\
 K_t&=\alpha I+UD_tU^T
 -UD_tV^T(H+VD_tV^T)^{-1}VD_tU^T.\label{ifc:eq:short}
\end{align}
Assume $L_0\succ0$ and $a_i(t+1)\ge a_i(t)$ for every $i,t$.
No disjointness, orthogonality, commutation, or angle condition is imposed
on the feature columns or on their overlaps.

\begin{proposition}[Monotonicity survives overlaps]\label{ifc:prop:mono}
Both $L_t$ and $K_t$ are nondecreasing in Loewner order. For every $y$,
\begin{equation}\label{ifc:eq:min}
 y^TK_ty=\alpha\|y\|^2+
 \min_z\left\{z^THz+
       \sum_i a_i(t)(u_i^Ty-v_i^Tz)^2\right\}.
\end{equation}
When $H=0$ and $VD_tV^T\succ0$, let
$F_t=UD_t^{1/2}$ and let $P_t$ project onto
$\ker(VD_t^{1/2})$. Then $K_t=\alpha I+F_tP_tF_t^T$.
\end{proposition}
\begin{proof}
Put $Z=(V^T,U^T)^T$. Then
$L_{t+1}-L_t=Z(D_{t+1}-D_t)Z^T\succeq0$.
Completing the square in $z$ gives \eqref{ifc:eq:min}. Its objective
increases for every fixed $z,y$, so its minimum increases as well.
Finally substitute the usual full-row-rank formula for the orthogonal
projector. All inverses exist because principal blocks of $L_t$ are SPD.
\end{proof}

A radial model is obtained by taking $U=I$, $V=\mathbf1^T$,
$H=0$, with $a_i$ equal to block energies. Arbitrary $U,V$ permit
overlapping outputs and several overlapping constraints. For the
hard-constraint interpretation, the leading block must already have
full rank at the initial state. The proposition makes no claim about
the cost of rank births before that state.

Define Sra's divergence
\[
 \Sc(B,K)=\log\det((B+K)/2)
             -\tfrac12\log\det B-\tfrac12\log\det K\ge0.
\]
\begin{lemma}[Monotone production]\label{ifc:lem:production}
If $K'\succeq K\succ0$ and $B\succ0$, then
\begin{equation}\label{ifc:eq:production}
 \Sc(B,K')-\Sc(B,K)
 \le\tfrac12\log\frac{\det K'}{\det K}.
\end{equation}
The cache $B$ can be chosen adaptively and need not be a good
approximation to either target.
\end{lemma}
\begin{proof}
On the segment $K(u)=K+u(K'-K)$ the derivative is
\[
 \tr\left[((B+K(u))^{-1}-\tfrac12K(u)^{-1})(K'-K)\right]
 \le\tfrac12\tr(K(u)^{-1}(K'-K)).
\]
Integrate. Positivity of $K'-K$ justifies the inverse-order comparison.
\end{proof}

\begin{theorem}[Cubic maintenance of the lifted overlapping model]
\label{ifc:thm:cubic}
Let $d=q+r$, $p=O(d)$, and suppose \eqref{ifc:eq:lift} is supplied through
its fixed features and its current weights for $T$ states, with
\[
 \lambda I\preceq L_0\preceq L_t\preceq M I,
 \qquad \kappa=M/\lambda.
\]
After $O(d^3)$ initialization there is a randomized solver for one
right-hand side per state, always returning a vector with the requested
certified absolute residual, with expected cost
\begin{equation}\label{ifc:eq:cost}
 \widetilde O\bigl(d^3+Td^2+d^3\log\kappa\bigr).
\end{equation}
The same bound solves $K_ty=b_t$. The suppressed factors include
logarithms of $T,d,\kappa$, and the scaled residual tolerances.
In particular, $T=\widetilde O(d)$ and polynomial spectral ratios and
tolerances give expected $\widetilde O(d^3)$ arithmetic work.
This includes maintaining the constraint block: it requires no fresh
dense factorization of $H+VD_tV^T$ at each state. Supplying the features
and weights, and verifying monotonicity, are part of the input contract.
\end{theorem}
\begin{proof}
Maintain a Cholesky-factored SPD cache $B$, initialized at $L_0$.
A verified direction $w$ with
$(w^TL_tw)/(w^TBw)\notin[1/2,3/2]$ is repaired by
\begin{equation}\label{ifc:eq:bfgs}
 B^+=B-\frac{Bw\,w^TB}{w^TBw}
          +\frac{L_tw\,w^TL_t}{w^TL_tw}.
\end{equation}
Whiten $L_t$ and rotate the direction to the first coordinate. The
transformed cache has the form
$\left(\begin{smallmatrix}b&c^T\\c&D+cc^T/b\end{smallmatrix}\right)$,
and the update changes it to $\diag(1,D)$. Its divergence decrement is
\[
 \log\frac{1+b}{2\sqrt b}
 +\log\left(1+\frac{c^T(D+I)^{-1}c}{b(1+b)}\right)
 \ge\gamma:=\log\frac5{2\sqrt6}>0.0204.
\]
Summing Lemma~\ref{ifc:lem:production} between repairs, and discarding
the nonnegative terminal divergence, gives the pathwise bound
\begin{equation}\label{ifc:eq:repaircount}
 \gamma R\le\tfrac12\log\frac{\det L_T}{\det L_0}
 \le\tfrac d2\log\kappa.
\end{equation}
Any reset to the exact current matrix decreases divergence and preserves
this accounting.

Here is a finite witness test, to specify the cost of finding repairs.
For $B=CC^T$, apply $J=C^{-1}L_tC^{-T}-I$ without forming it.
Take an independent Rademacher vector $\xi$ and normalize $J^k\xi$,
with $(11/12)^{2k}\le1/(16d)$, skipping a zero vector. For the result
$z$, compute the two-dimensional compression to $\operatorname{span}\{z,Jz\}$
and verify its extreme Rayleigh directions in the original quotient.
If $\|J\|>3/4$, a random start has squared overlap at least $1/(2d)$
with an extreme eigenvector with probability at least $1/12$:
apply the second-moment bound
$\mathbb E(v^T\xi)^4\le3$ to $(v^T\xi)^2$.
After the powers, the squared mass on $|\lambda(J)|<11/16$ is at
most $1/8$, so $\|Jz\|\ge(11/16)\sqrt{7/8}>5/8$.
The compression has norm at least $\|Jz\|$ and therefore supplies
a verified quotient outside $[1/2,3/2]$.
Repeating $O(\log(1/\delta))$ independent starts makes a false good
declaration have conditional probability at most $\delta$.
A correct good declaration implies
$B/4\preceq L_t\preceq7B/4$.

For a good declaration run the preconditioned Richardson iteration
$y_{j+1}=y_j+B^{-1}(b-L_ty_j)$ from zero. In whitened coordinates
its error contracts by at least $3/4$ per iteration. The bounds
$\lambda,M$ give an a priori logarithmic cap for the requested
absolute residual, which is checked directly. Failure of the cap
triggers exact elimination and a cache reset. Thus test errors do
not affect correctness.

There are at most $T+R$ tests, with $R$ bounded by
\eqref{ifc:eq:repaircount}. Choose $\delta$ inverse quadratic in a
known upper bound for $T+R+1$. Then the expected number of dense
fallbacks is at most one. Every target product costs $O(d^2)$ because
$L_tv=\diag(H,\alpha I)v+ZD_tZ^Tv$, and $p=O(d)$.
Every cache application and Cholesky rank update also costs $O(d^2)$.
For \eqref{ifc:eq:bfgs}, perform the positive update before the negative
downdate; both the intermediate and final matrices are SPD.
These observations and \eqref{ifc:eq:repaircount} prove \eqref{ifc:eq:cost}.
Forming any fallback target from the features costs $O(d^3)$.

To solve the Schur system, solve $L_t(x,y)^T=(0,b)^T$ and return $y$.
For joint residual $(r_1,r_2)$ and block notation as above,
\[
 K_ty-b=r_2-B^TA^{-1}r_1.
\]
Here $\|B^TA^{-1}\|\le M/\lambda=\kappa$. Asking for joint
residual at most $\varepsilon/(1+\kappa)$ certifies Schur residual
at most $\varepsilon$. This changes only the logarithmic iteration
count and requires no application of $A^{-1}$ by the implementation.
\end{proof}

The theorem includes target applications and constraint-block maintenance
for the stated class.
It removes both the overlap and repeated constraint-factorization costs
in that class. If an epoch of the signing algorithm has such an exact
representation with $d=O(s)$ and the stated logarithmic bounds, its
solve term is cubic in $s$; geometric cleanup then sums these bounds to
cubic total work. An exact representation of the general filtered
Koml\'os factor in this class has not been established.

\subsection{The exact overlap term that cannot be discarded}
For two arbitrary SPD block matrices write their elimination forms as
\[
 L_i=\begin{pmatrix}
 A_i&A_iX_i\\
 X_i^TA_i&S_i+X_i^TA_iX_i
 \end{pmatrix},\qquad i=0,1.
\]
The matrices $X_i=A_i^{-1}(L_i)_{12}$ are the correction maps retained
by elimination. Put
\[
 \overline S=(S_0+S_1)/2,
 \quad W=\tfrac12(X_0-X_1)^T
          (A_0^{-1}+A_1^{-1})^{-1}(X_0-X_1)\succeq0.
\]
\begin{proposition}[Exact block-overlap identity]\label{ifc:prop:overlap}
One has
\begin{equation}\label{ifc:eq:overlap}
 \Sc(L_0,L_1)=\Sc(A_0,A_1)+\Sc(S_0,S_1)
 +\log\frac{\det(\overline S+W)}{\det\overline S}.
\end{equation}
The last term is nonnegative, and vanishes exactly when $X_0=X_1$.
For centered Gaussian densities with covariances $L_0,L_1$, their
Hellinger affinity is $\exp(-\Sc(L_0,L_1)/2)$.
\end{proposition}
\begin{proof}
Block elimination of $(L_0+L_1)/2$ gives leading block
$(A_0+A_1)/2$ and Schur complement $\overline S+W$. The quadratic
identity behind the latter assertion is
\begin{align*}
 &X_0^TA_0X_0+X_1^TA_1X_1\\
 &\quad-(A_0X_0+A_1X_1)^T(A_0+A_1)^{-1}(A_0X_0+A_1X_1)\\
 &=(X_0-X_1)^T(A_0^{-1}+A_1^{-1})^{-1}(X_0-X_1).
\end{align*}
It follows either by completing the square or by minimizing the sum
of the two positive quadratic forms over a common center.
Use $\det L_i=\det A_i\det S_i$ to obtain \eqref{ifc:eq:overlap}.
The parallel sum is SPD, so $W=0$ exactly for equal correction maps.
The affinity formula follows by integrating the square root of the
two Gaussian densities; equivalently use
$\det(L_0^{-1}+L_1^{-1})=\det(L_0+L_1)/(\det L_0\det L_1)$.
\end{proof}

\paragraph{A basic tiny-ridge test.}
Set $A_0=A_1=I_r$, $S_0=S_1=\alpha I_r$, $X_0=0$, and
$X_1=c\sqrt\alpha I_r$. Both leading blocks and Schur complements
are constant and $\det L_0=\det L_1=\alpha^r$, but
\begin{equation}\label{ifc:eq:shear}
 \Sc(L_0,L_1)=r\log(1+c^2/4),\qquad
 \|L_1-L_0\|_F^2=2rc^2\alpha+rc^4\alpha^2.
\end{equation}
For $c=3$, the overlap loss is $r\log(13/4)$ independently of
$\alpha$, while ordinary squared motion tends to zero. Repeated
switches preserve these facts. This rules out a bound on joint overlap
production from these marginal data alone. It is not a lower bound
for every solver, and the displayed matrices are easy to solve.

When the changing correction maps are explicitly cheap to apply,
one may transport the cache by their triangular congruence. This
removes the last term for that representation, but the cost of
constructing, applying, and maintaining the congruences must be counted.
For dense unrelated maps, Schur elimination itself supplies no such
cost bound. Theorem~\ref{ifc:thm:cubic} avoids this issue through the
monotone full lift, rather than dropping the correction term.

\section{Deflation and hereditary elimination}
\label{ifc:sec:tree}

Independent blocks admit independent elimination. A block diagonal matrix
can be solved block by block even if every block changes nonmonotonically.
The cache-repair examples from Section~\ref{c:models} are not lower bounds
for such solvers. More generally, a recursive solver can eliminate the
uncoupled directions locally and retain only the interfaces between
blocks. This is established hierarchical linear algebra
\cite{ulv,hierarchical}; the application here is to identify exactly
which structure survives the constrained, filtered Gram construction.

\subsection{Strong couplings can be deflated without a ridge loss}
Consider an SPD matrix
\[
 K=\begin{pmatrix}A&C\\C^T&D\end{pmatrix},\qquad
 Z=A^{-1/2}CD^{-1/2}.
\]
Positive definiteness implies $\|Z\|<1$. Let
$Z=U\Sigma V^T$ be a singular value decomposition. For $0<\tau<1$,
let $Z_>$ retain precisely the singular pairs with singular value
greater than $\tau$, and define
\[
 B=\begin{pmatrix}
 A&A^{1/2}Z_>D^{1/2}\\
 D^{1/2}Z_>^TA^{1/2}&D
 \end{pmatrix}.
\]
\begin{lemma}[Exact strong-coupling deflation]\label{ifc:lem:deflate}
One has $B\succ0$ and
\begin{equation}\label{ifc:eq:deflate}
 (1-\tau)B\preceq K\preceq(1+\tau)B.
\end{equation}
The generalized eigenvalues are $1$ on retained singular pairs and
$1\pm\sigma_j$ on discarded pairs, together with unit eigenvalues on
unpaired directions. In particular, the preconditioned condition number
is at most $(1+\tau)/(1-\tau)$, independently of the smallest eigenvalues
of $A,D,K$ and of how close a retained singular value is to one.
\end{lemma}
\begin{proof}
Congruence by $\diag(A^{-1/2},D^{-1/2})$, followed by the two singular
bases, splits the pair into independent two-dimensional pencils. A
retained pair has identical blocks
$\left(\begin{smallmatrix}1&\sigma\\\sigma&1\end{smallmatrix}\right)$
in both matrices. A discarded pair has that block in $K$ and $I_2$
in $B$. The asserted eigenvalues and inequalities follow.
\end{proof}
The conclusion is stronger than bounding the discarded block in
absolute norm and dividing by a small ridge. It uses the orthogonality
of retained and discarded singular pairs. It is a local algebraic
statement, not a claim that a fresh dense SVD is cheap or that such
truncations can be composed at every tree level without a further
error analysis. The exact recursive solver below uses supplied exact
factors, so it avoids that latter issue.

\subsection{Relative error through transported factors}

The one-level normalized singular-value truncation is a standard
structured-preconditioning result; see
Agullo--Darve--Giraud--Harness~\cite{ADGH2018} and the formulation in
Feliu-Fab\`a--Ho--Ying~\cite[Proposition 1.1]{FHY2020}.
Its multilevel use needs care. In particular, replacing a principal
diagonal child by a relative approximation while leaving its parent
coupling unchanged need not preserve positive definiteness.

Here is a factor formulation with a direct error composition proof.
For a node $v$, let $K_v$ be its exact original principal matrix and
$D_v=\diag(K_{v_0},K_{v_1})=L_vL_v^T$, with block diagonal $L_v$.
In the normalized matrix $L_v^{-1}K_vL_v^{-T}$, keep the cross singular
values exceeding $\tau_v$ and call the resulting SPD matrix $H_v$.
Define
\[
 F_v=L_vH_v^{1/2}L_v^{-1}.
\]
Then the exact one-level truncated matrix is
$\widehat K_v=F_vD_vF_v^T$ and
\[
 (1-\tau_v)\widehat K_v\preceq K_v
 \preceq(1+\tau_v)\widehat K_v.
\]
Set $B_v=K_v$ at the leaves and recursively set
\begin{equation}\label{ifc:eq:transported-tree}
 B_v=F_v\diag(B_{v_0},B_{v_1})F_v^T.
\end{equation}
This transports the child approximations through the parent factor;
it does not simply replace the diagonal children in $\widehat K_v$.

\begin{proposition}[Composition with exact principal normalization]
Suppose a tree has depth at most $L$ and $\tau_v\le\tau<1$ at every
internal node. Then the matrix in~\eqref{ifc:eq:transported-tree} is SPD and
\[
 (1-\tau)^L B_{\rm root}\preceq K
 \preceq(1+\tau)^L B_{\rm root}.
\]
In particular, $\tau\le1/(4L)$ implies
$\kappa(B_{\rm root}^{-1/2}KB_{\rm root}^{-1/2})<2$.
\end{proposition}
\begin{proof}
If $aB_{v_i}\preceq K_{v_i}\preceq bB_{v_i}$ for both children,
congruence by $F_v$ gives
$aB_v\preceq\widehat K_v\preceq bB_v$.
The one-level comparison therefore replaces $a,b$ by
$(1-\tau_v)a,(1+\tau_v)b$. Induction proves the claim.
Finally
$L\log((1+\tau)/(1-\tau))\le2L\tau/(1-\tau)\le2/3<\log2$.
\end{proof}

\paragraph{Representation and the cost of normalization.}
If $r_v$ cross singular pairs are retained, write
$H_v=I+Y_v(C_v-I)Y_v^T$, where $Y_v^TY_v=I$ and
$C_v$ is a direct sum, up to permutation, of the positive matrices
$\left(\begin{smallmatrix}1&\sigma_j\\\sigma_j&1\end{smallmatrix}\right)$.
With $U_v=L_vY_v$ and $V_v=L_v^{-T}Y_v$,
\[
 F_v^{\pm1}=I+U_v(C_v^{\pm1/2}-I)V_v^T,
 \qquad V_v^TU_v=I.
\]
Thus supplied exact $U_v,V_v,\sigma_j$ give applications of
$F_v^{\pm1}$ and their transposes in $O(n_vr_v)$ operations.
For leaves of size at most $b$, forming the dense leaf factors costs
$O(nb^2)$, and one solve with $B$ costs
\[
 O\left(nb+\sum_v n_vr_v\right).
\]
If $r_v\le r$ on a balanced tree, the sum is $O(nrL)$.

These bounds charge neither exact principal normalization nor the search
for its singular vectors. They must be added separately. In particular,
$V_v=D_v^{-1}U_v$ requires exact solves in the original principal
children, not in their approximations $B_{v_i}$.
Dense Cholesky factorizations of all original principal children already
cost $O(n^3)$ per state, before accounting for the requested accuracy of
the singular decompositions. Exact singular vectors in this algebraic
statement are supplied data; no finite exact-SVD primitive is asserted.
The proved exact-interface solver makes this normalization cheap when
the original matrices already have supplied exact low-rank structure;
it does not make normalized low-rank structure discoverable for free.
Replacing the normalizations by approximate child solves requires a
separate perturbation analysis, especially when retained singular values
approach one.

\paragraph{A four-dimensional check.}
Let $A=\left(\begin{smallmatrix}1&t\\t&1\end{smallmatrix}\right)$ and
\[
 K=\begin{pmatrix}A&(1-\delta)A\\(1-\delta)A&A\end{pmatrix},
 \qquad 0<\delta<t/(1+t),\quad 0<t<1.
\]
Truncating the child cross singular value $t$ replaces $A$ by $I$.
Both parent cross singular values equal $1-\delta$ and can be kept
exactly. Nevertheless, assembling diagonal children $I$ and the
unchanged parent cross matrix gives minimum eigenvalue
$1-(1-\delta)(1+t)<0$.
The transported construction remains SPD and has generalized
eigenvalues $1-t,1-t,1+t,1+t$ in this example.


\subsection{The correct hereditary rank is external coupling rank}
For a coordinate set $J$ define
$r_J(K)=\rank K_{J,J^c}$. This measures interaction with the entire
complement, not just with a sibling in a proposed tree.
\begin{lemma}[Rank survives arbitrary elimination]\label{ifc:lem:heredity}
Let $K\succ0$, let $R$ be a retained coordinate set, and let
$S=K/K_{R^c,R^c}$ be the Schur complement after eliminating $R^c$.
For every original cut $J,J^c$,
\begin{equation}\label{ifc:eq:heredity}
 \rank S_{J\cap R,J^c\cap R}\le r_J(K).
\end{equation}
\end{lemma}
\begin{proof}
For any SPD block matrix, the block inverse formula shows that the
off-diagonal block of the inverse has exactly the same rank as the
original off-diagonal block. Also $S^{-1}=(K^{-1})_{R,R}$. Therefore
\[
 \rank S_{J\cap R,J^c\cap R}
 =\rank(K^{-1})_{J\cap R,J^c\cap R}
 \le\rank(K^{-1})_{J,J^c}=r_J(K).
\]
Empty blocks cause no difficulty.
\end{proof}
Thus genuine non-overlap remains non-overlap under elimination, even
if coordinates are eliminated on both sides of the cut. Low external
rank remains low rank. A bound only on the original sibling block is
insufficient for this conclusion: interactions through ancestors must
also be included. The correction maps still have to be computed, but
their relevant interface dimension is now explicitly bounded.

\subsection{A fully charged recursive solver with no monotonicity assumption}
Fix a balanced binary partition tree for an $s\times s$ SPD matrix.
Dense leaf blocks have orders at most $b$, with
$b\ge\max\{1,r\}$. At every internal node the cross block between
its children is supplied as $UV^T$ of rank at most $r$. This is the
usual exact hierarchically off-diagonal low-rank representation.
The factor matrices, tree, and dense leaves are input data to the next
statement; discovering them is not included for free.

\begin{proposition}[Direct tree cost]\label{ifc:prop:treecost}
One can factor the represented matrix, and solve one right-hand side,
using respectively
\begin{align}
 F(s)&=O\!\left(sb^2+sr^2[1+\log(s/b)]^2\right),\label{ifc:eq:treefactor}\\
 Q(s)&=O\!\left(sb+sr[1+\log(s/b)]\right).\label{ifc:eq:treesolve}
\end{align}
Here $1\le b\le s$; the constants cover trees with leaf sizes within
a fixed factor of $b$. These are exact real-arithmetic bounds with
no spectral-ratio factor. For $T=\widetilde O(s)$ arbitrarily changing
matrices, supplied afresh in this representation, $b,r=\widetilde
O(\sqrt s)$ give total $\widetilde O(s^3)$ solve work. In particular,
no monotonicity or bound on cumulative S-divergence production is needed.
\end{proposition}
\begin{proof}
At a node write
\[
 K_v=D_v+W_vJ_vW_v^T,\quad
 D_v=\diag(K_a,K_b),\quad W_v=\diag(U,V),\quad
 J_v=\begin{pmatrix}0&I\\I&0\end{pmatrix}.
\]
Assuming the children have been factored, solve their systems with the
columns of $W_v$ to obtain $Y_v=D_v^{-1}W_v$, and factor
$H_v=J_v+W_v^TY_v$. This is an at most $2r$ square matrix. It is
invertible because $K_v,D_v,J_v$ are invertible and the determinant
lemma applies. It need not be positive definite; ordinary exact
elimination with row pivoting suffices. The inverse formula is
\begin{equation}\label{ifc:eq:woodburytree}
 K_v^{-1}z=D_v^{-1}z-Y_vH_v^{-1}Y_v^Tz.
\end{equation}
Use the first nonzero pivot in each column for a deterministic pivot rule.

The cost of one node solve, beyond its two child solves, is
$O(|v|r+r^2)$. Since $r\le b$, summing over levels proves
\eqref{ifc:eq:treesolve}. Building $Y_v$ costs at most $r$ right-hand
sides in each child, and building/factoring $H_v$ costs
$O(|v|r^2+r^3)$. At each level the node sizes sum to $s$.
Summing the child-solve costs over all levels gives
$O(sr^2[1+\log(s/b)]^2+srb[1+\log(s/b)])$.
The mixed term is bounded by a constant times the right side of
\eqref{ifc:eq:treefactor}. Dense leaf factorizations cost $O(sb^2)$.
This proves the factorization bound, and summing it over states gives
the final assertion. A zero-rank node simply uses its two child solvers.
\end{proof}
This elementary implementation is in the same family as established
hierarchical direct solvers \cite{ulv,hierarchical}. Its significance
for the present problem is the removal of the monotonicity requirement,
not a new general-purpose linear algebra algorithm. Structured
representation construction is itself a substantive task; randomized
compression results explicitly charge it \cite{compression}.

\subsection{The filter and the exact constraints do not automatically destroy rank}
Here there is a direct connection to the actual Koml\'os operator.
For a fixed cut $J,J^c$, suppose
\[
 \rank M_{J,J^c}\le a,\qquad
 \max\{\rank T_{J,J^c},\rank T_{J^c,J}\}\le c,
 \qquad \rank P_{J,J^c}\le p.
\]
The matrices $M,P$ are symmetric, $T$ need not be symmetric, and
$f$ is a degree-$d$ polynomial.
\begin{lemma}[Closure for the filtered Gram]\label{ifc:lem:filtertree}
For $E=f(M)T$ and $G=\alpha I+EPE^T$,
\begin{equation}\label{ifc:eq:filtertree}
 \rank f(M)_{J,J^c}\le da,\qquad
 \rank G_{J,J^c}\le2(da+c)+p.
\end{equation}
These are exact rank statements, requiring no singular-value gap.
\end{lemma}
\begin{proof}
Put $A=M_{J,J}$ and $C=M_{J,J^c}$. The recurrence for block powers
shows that the columns of $(M^k)_{J,J^c}$ lie in
\[
 \operatorname{span}\{\operatorname{range}C,
                 \operatorname{range}AC,\ldots,
                 \operatorname{range}A^{k-1}C\}.
\]
The span for all powers through degree $d$ has dimension at most $da$.
For two compatible square matrices,
$(XY)_{J,J^c}=X_{J,J}Y_{J,J^c}+X_{J,J^c}Y_{J^c,J^c}$, so the
cross rank of a product is at most the sum of the two cross ranks.
Apply this identity first to $E$, then to $EPE^T$. Symmetry of $f(M)$
controls both orientations; the two orientations of $T$ were assumed.
The ridge contributes no cross block.
\end{proof}
The normalization $f(M/\Lambda)$ changes no rank count. Thus a proved
hereditary interface bound for $M,T,P$ would pass through the actual
finite polynomial filter with only its degree as a multiplier.
Lemma~\ref{ifc:lem:heredity} then explains why exact elimination preserves
the resulting external-rank bounds. This is an algebraic closure
statement; it does not itself construct the factors in
Proposition~\ref{ifc:prop:treecost} or establish a small rank bound for all
states of the walk.

For example, when retained rows are supported inside fixed coordinate
blocks, $M$ is block diagonal plus its rank-one Perron ridge, and $T$
is block diagonal. The row-local constraints give a block diagonal
projector, followed by at most two global equations ($x$ and aggregate
energy). Consequently, across a union-of-blocks cut one may take
$a=1$, $c=0$, and $p\le2$, giving cross rank at most $2d+2$ for $G$.
This checks the closure on the actual formulas rather than on an
unrelated bounded-matrix path. Such block-supported signing inputs
can also be solved independently from the outset; no new discrepancy
theorem for this elementary subclass is claimed.

\subsection{Spatial accounting is not yet temporal accounting}
At a node of a fixed tree, use its two principal child blocks to define
the normalized cross matrix $Z_v$ as above. The determinant identity gives
\[
 I_v:=-\log\det(I-Z_v^TZ_v)
   =\log\det K_{v_L}+\log\det K_{v_R}-\log\det K_v.
\]
If $r_v(\tau)$ singular values exceed $\tau$, then
$\tau^2r_v(\tau)\le\sum_j\sigma_j(Z_v)^2\le I_v$.
Telescoping over the tree proves
\begin{equation}\label{ifc:eq:spatial}
 \tau^2\sum_vr_v(\tau)
 \le\sum_{\ell\text{ leaf}}\log\det K_\ell-\log\det K
 \le s\log\kappa
\end{equation}
whenever $\lambda I\preceq K\preceq\kappa\lambda I$.
Thus the determinant does give an exact accounting for the total
strong-coupling count at one state. It does not bound the time needed
to find those couplings, their squared-rank factorization costs at large
nodes, or their replacements over subsequent states.

A small check prevents overinterpreting this bound. For a half-dimension
$h$, take
\[
 E=\sqrt{3\alpha/2}\,[I_h\ I_h],\qquad K=\alpha I+E^TE.
\]
Across the displayed cut all $h$ canonical correlations equal $3/5$,
although $\|E\|_F^2=3h\alpha$ can be arbitrarily small. The condition
number of $K$ is only four. This is an easy matrix, and permuting into
coordinate pairs decouples it. The example disproves only the claim
that ordinary factor energy forces low interface rank in a prescribed
tree. It does not obstruct adaptive trees, block solvers, or the actual
signing algorithm.

The verified conclusion is therefore sharper than a monotone model:
decoupling and hereditary rank can support a cubic solver despite
arbitrary nonmonotone motion. For unrestricted Koml\'os, the missing
statement is still a cheaply constructible and maintainable structure
for its actual systems. Section~\ref{new:structured} proves construction bounds for explicit structural classes, with all construction costs included.


\section{A local derivative estimate for the filtered factor}\label{app:factor}

The result below concerns the factor occurring in the signing algorithm,
not an arbitrary bounded matrix path. It is a local estimate at the start
of a high-branch accepted direction. It does not prove the signed-production
hypothesis, and it does not assert a finite-step or full-walk variation bound.

Fix the retained profiles and scales. Index medium row-sign pairs by $i$,
write $z_i=p_i/r_i^2$, $\alpha_i=32\exp(2u_i/7)$, and set $\beta=2/7$.
Thus
\[
 W=\sum_i\alpha_i z_i z_i^T,\qquad
 M=W+B+\eta\mathbf1\mathbf1^T,\qquad
 T=\sum_i\alpha_i\beta(z_i^Tv)z_i g_i^T.
\]
Here $Mv=Lv$, $\|v\|=1$, and $v>0$. The fixed-metadata curve is
$x(t)=x+th$. Throughout this section derivatives are evaluated at $t=0$.
The accepted high-branch direction supplies
\[
 |a_i|=|\beta g_i^Th|\le K,\qquad
 b_i=\sum_jz_{ij}h_j^2\le K^2,
 \qquad |L'|\le K/64,
\]
as well as the already proved Perron-response estimates
\[
 \|v'\|\le2K,\qquad |v'_j|\le4Kk_vv_j,
 \qquad k_v=1+\log(1/\min_jv_j).
\]
The row facts used below are $0\le z_{ij}\le1/100$,
$\sum_jz_{ij}\le1$, $1/2\le L\le9/8$, and
$\mathcal U_{jj}\le\rho_0\mathcal D_{jj}$ with
$\rho_0=289/2450$.

\begin{proposition}[Actual-factor derivative at an accepted high state]
With the preceding notation,
\[
 \|T\|_F\le\frac{17}{35}L<1,
 \qquad \|T'[h]\|_F\le\frac{207}{70}Kk_v<3Kk_v.
\]
Let $f$ be any degree-$d\ge1$ polynomial satisfying $0\le f(t)\le1$
on $[0,1]$, held fixed during the differentiation, and define
\[
 E^{\#}(x)=f\!\left(\frac{M(x)}{\max\{1,L(x)\}}\right)T(x).
\]
Then, at every differentiability point, and for either directional
derivative when $L=1$,
\[
 \|(E^{\#})'[h]\|_F\le3K(d^3+k_v).
\]
In particular the local derivative is $\widetilde O(K)$ for the
algorithm's polynomial degree and Perron-coordinate lower bound.
\end{proposition}

\begin{proof}
Let $C$ have columns $\sqrt{\alpha_i}z_i$ and let $R$ have rows
$\sqrt{\alpha_i}\beta(z_i^Tv)g_i^T$. Then $T=CR$,
$\|C\|^2=\|W\|\le L$, and $R^TR=\mathcal U$. Moreover
\[
 \operatorname{tr}\mathcal D
 =2\sum_i\alpha_i(z_i^Tv)^2\sum_jz_{ij}
 \le2v^TWv\le2L.
\]
Therefore
\[
 \|T\|_F^2\le L\operatorname{tr}\mathcal U
 \le2\rho_0L^2=\left(\frac{17L}{35}\right)^2.
\]

Since $\alpha_i'=a_i\alpha_i$ and
$g_i'=-7\operatorname{diag}(z_i)h$, exact differentiation gives
\[
 T'=C\operatorname{diag}(c_i)R-2H_v\operatorname{diag}(h),
 \quad c_i=a_i+\frac{z_i^Tv'}{z_i^Tv},
 \quad H_v=\sum_i\alpha_i(z_i^Tv)z_i z_i^T.
\]
The denominators are positive because $z_i\ne0$ and $v>0$.
Positivity of the $z_i$ and the componentwise Perron bound imply
$|c_i|\le K(1+4k_v)$. Hence
\[
 \|C\operatorname{diag}(c_i)R\|_F
 \le\frac{17L}{35}K(1+4k_v).
\]
For the remaining term write $H_v=C\operatorname{diag}(z_i^Tv)C^T$.
Then
\begin{align*}
 \|H_v\operatorname{diag}(h)\|_F^2
 &\le L\sum_i\alpha_i(z_i^Tv)^2
                  \sum_jz_{ij}^2h_j^2\\
 &\le\frac{LK^2}{100}\sum_i\alpha_i(z_i^Tv)^2
 \le\frac{L^2K^2}{100}.
\end{align*}
Combining the bounds, and using $k_v\ge1$ and $L\le9/8$, yields
\[
 \|T'\|_F
 \le LK\left(\frac{24}{35}+\frac{68}{35}k_v\right)
 \le\frac{207}{70}Kk_v.
\]

For completeness, the needed operator derivative bound for $f$ is
elementary. Expand $f((x+1)/2)=\sum_{k=0}^d c_kT_k(x)$ in Chebyshev
polynomials. The uniform bound on $f$ gives $|c_k|\le2$ for $k\ge1$.
For any self-adjoint contraction $X$, differentiation of the Chebyshev
recurrence gives
\[
 D T_k(X)[H]
 =U_{k-1}(X)H
    +2\sum_{j=1}^{k-1}U_{k-1-j}(X)H T_j(X).
\]
Since $\|T_j(X)\|\le1$ and $\|U_j(X)\|\le j+1$,
$\|D T_k(X)[H]\|\le k^2\|H\|$. Consequently, for $0\preceq A\preceq I$,
\[
 \|Df(A)[H]\|
 \le4\sum_{k=1}^d k^2\|H\|
 \le4d^3\|H\|.
\]
Here $-KW\preceq M'\preceq KW$, so $\|M'\|\le KL$.
For $A=M/\max\{1,L\}$, the scalar normalization and the bound on $L'$
therefore give $\|A'\|\le65K/64$. Using
$\|T\|_F\le153/280$, the filter derivative contributes at most
\[
 4d^3\frac{65K}{64}\frac{153}{280}
 <\frac94d^3K<3d^3K.
\]
The other product-rule term is bounded by $\|T'\|_F<3Kk_v$,
proving the claim.
\end{proof}

\paragraph{Scope of the estimate.}
The weighted energy bound does imply a budget for the formal tangent
increments $d_t(E^{\#})'[h_t]$ over high states. It does not imply the
same budget for $E^{\#}(x_t+d_th_t)-E^{\#}(x_t)$. The accepted leakage
and Perron-response contracts hold at the beginning of a chord; they
have not been proved at all its intermediate points. The finite-step
Perron residual contains an $O(K^2t^2/\eta)$ eigenvector error, so this
distinction cannot be dropped. In addition, even an absolute Frobenius
variation bound would not yet control the ridge-normalized target
variation required by the S-divergence ledger. At low states the
accepted direction has no corresponding high-leakage contract, and the
proposition is not asserted there.


\section{Model behavior and limits of energy charging}\label{c:models}
\subsection{Input and output rotations do not require repeated repairs}
If $P=I$ and $E_t=E_0R_t$ with $R_t$ orthogonal, the dual target $G$ is
constant. Its cache stays exact, and neither cache is repaired. The primal
production telescopes to one terminal S-divergence, at most
$O(s\log\kappa_*)$. If $E_t=U_tE_0$ with $U_t$ orthogonal and $P$ fixed,
the primal target is constant; the same reasoning applies with the two
sides exchanged. These are exact algebraic models, not asserted
trajectories of the full signing algorithm. They explain why a single
primal or single dual cache would be the wrong conjectural object.

\subsection{A nontrivial normal-birth model has the desired correction bound}
Partition $x$ into disjoint blocks $x_i$. Let row $i$ of $E$ be $x_i^\top$
on its block and zero elsewhere. Assume bounded $\|E\|$ and put $a_i=\|x_i\|^2$ and
$P_x=I-xx^\top/\|x\|^2$ when $x\ne0$. Direct multiplication gives
\begin{equation}\label{c:eq:radial}
 EP_xE^\top=\diag(a)-\frac{aa^\top}{\sum_i a_i}.
\end{equation}
At $x=0$, use $P_x=I$ and interpret the displayed matrix as zero.
For nonzero $x$, its quadratic form is
$\sum_i a_i y_i^2-(\sum_i a_i y_i)^2/(\sum_i a_i)$, which equals
$\min_c\sum_i a_i(y_i-c)^2$. Increasing any $a_i$ increases the
expression for every $c$, so the minimum is nondecreasing. This proves
Loewner monotonicity, including zero coordinates by continuity. Under blockwise tangent steps
$x_i^\top h_i=0$, one has
$a_i^+=a_i+\tau^2\|h_i\|^2$. The dual target is therefore Loewner
nondecreasing, even though many constraint singular values can emerge
from zero and $P_x$ itself changes.

Lemma~\ref{c:lem:drift} bounds total dual production by
$\tfrac12\log(\det G_T/\det G_0)=O(s\log\kappa_*)$.
It follows that the number of dual corrections is of that order, regardless
of what the primal cache does. Alternation within a state makes the number
of primal corrections at most the number of dual corrections plus the
number of states. Thus the actual two-cache controller has the desired
$\widetilde O(s)$ correction bound in this model when the number of states is
$\widetilde O(s)$. This proves a cost bound in the model, not
the conjecture \eqref{d:conjecture} for all signing trajectories.

\subsection{What a proof still has to control}
Near a common target, writing $W=K^{-1/2}ZK^{-1/2}$,
\[
 \mathcal{S}(K,K+Z)=\tfrac18\|W\|_F^2+O(\|W\|_F^3).
\]
The expansion follows by subtracting the convergent series for
$\log\det(I+W/2)$ and $\tfrac12\log\det(I+W)$ when $\|W\|<1/2$.
The linear terms cancel and the quadratic coefficient is $1/8$; higher
traces are bounded by a constant times $\|W\|_F^3$. It does not turn the
walk's Euclidean energy bound into a bound on normalized Gram motion.
The production conjecture concerns both changing Gram geometries and is
strictly more than this local expansion.


\subsection{A limitation beyond these protected models}
The preceding cases do not justify a general energy-to-repair estimate.
Proposition~\ref{regmaint:separation}, with no added regularizer, gives
three targets of condition number ten for which an adaptive two-cache
controller can require $\Omega(s^2)$ rank-two repairs over $s$ states.
Their factors can nevertheless have total squared Frobenius motion at
most $9/(8s)$. The same proposition proves that this obstruction persists
under arbitrary selective regularizers within the stated small
effective-dimension allowance. These are algebraic models, not asserted
reachable states of the signing walk.


\subsection{Wide coordinate interfaces in an admissible local state}\label{ifc:sec:wideinterfaces}

The following example rules out deriving narrow coordinate interfaces
from the retained-row invariants alone.  It uses the actual $M,T,P$ and
the canonical filtered factor.  It is an algebraically admissible local
state; no reachability by the specified walk is asserted.  Its systems
nevertheless admit cheap global deflation, so this is not a solver lower
bound.

Let $\mathcal G$ be a connected bipartite $32$-regular graph on $v$
vertices whose nontrivial adjacency eigenvalues have absolute value at
most $16$.  Such graphs exist for arbitrarily large $v$, for example by
taking bipartite double covers of graphs from Friedman's theorem
\cite[Theorem 1]{bordenave}.  Make $B=32$ coordinate copies
of every edge and let $H\in\mathbb R^{n\times v}$ be their unsigned
incidence matrix.  Then, with $k=1024$,
\[
 n=512v,\qquad H^{\mathsf T}H=32(32I+A_{\mathcal G}).
\]
For every vertex include $R=32$ identical input rows, each equal to
$3/(2\sqrt k)$ on the incident coordinates and zero elsewhere.  There
are $m=n/16$ rows.  Each column has squared norm $9/64$, each row has
$q_i=9/4$, and $\max_jp_{ij}/q_i=1/1024$.  Thus all rows are medium,
cleanup deletes nothing, and the prescribed scale is $r_i=3/2$.

Retain all entries with zero offsets and set $x=t\mathbf1$, where
$t=0.9237728688\ldots\in(0,1)$ is the unique solution of
\[
 4e^{-58/7-t^2}\cosh(2t\sqrt k/7)=1.
\]
The left side increases from a value below one to a value above one: its logarithmic derivative is positive by $\tanh(at)\ge t\tanh(a)$ for $0\le t\le1$, with $a=2\sqrt{k}/7$ and $\sqrt{k}\tanh(a)>7$.
All row-sign barriers are
$u_\pm=-29-7t^2/2\pm t\sqrt k$, and
$e^{2u_+/7}+e^{2u_-/7}=1/2$.  Hence $u_+<(7/2)\log(1/2)<-1/100$.
There are no anchors or large or frozen rows; the only constraint is
$x^{\mathsf T}h=0$, so $P=I-\mathbf1\mathbf1^{\mathsf T}/n$.
Direct substitution into the retained-row formulas gives
\begin{equation}
 W=\frac{HH^{\mathsf T}}{2k},\quad
 M=W+\eta\mathbf1\mathbf1^{\mathsf T},\quad L=1+\eta n,
 \qquad T=cW,                                      \label{ifc:interface:actual}
\end{equation}
where
\[
 c=\frac{2}{7\sqrt n}
    \bigl(\sqrt k\tanh(2t\sqrt k/7)-7t\bigr)>0.
\]
Indeed $g_{i\pm}=(\pm\sqrt k-7t)z_i$ and the Perron vector is
$\mathbf1/\sqrt n$.  For $\eta=1/(4096Jn)$,
$\Phi=(\eta n)^3-\eta nt^2<0$ and $L\le9/8$.
The high branch is unambiguous: $L-1=\eta n$ exceeds the eigenvalue
certificate width $\eta^5/2^{60}$.

The nonzero non-Perron eigenvalues of $W$ lie in $[1/4,3/4]$, inside
the passband of the actual filter $f(M/\Lambda)$.  Consequently, for
$E=f(M/\Lambda)T$ and $K_0=EPE^{\mathsf T}$,
\begin{equation}
 K_0=c^2W^2f(M/\Lambda)^2P,\qquad
 \operatorname{range}K_0=U:=\operatorname{range}H\cap\mathbf1^\perp,
 \quad \dim U=v-2.                                \label{ifc:interface:range}
\end{equation}
Its nonzero eigenvalues lie between fixed positive multiples of $1/n$.

\paragraph{Every coordinate tree has a linear external interface.}
Write $r_J=\operatorname{rank}H_J$.  Since $\mathcal G$ is bipartite,
signing its vertex columns turns $H$ into oriented incidence.  Its
graphic-matroid connectivity is
$\lambda(J)=r_J+r_{J^c}-v+2$.
The spectral gap gives $e(A,A^c)\ge8|A|$ for $|A|\le v/2$.
A separator leaving components of size at most $v/2$ must therefore
have size at least $v/16$: choose a union of components with size in
$[v/4,v/2]$ and compare its boundary with $32$ times the separator size.
The usual centroid bag in a tree decomposition is such a separator, so
$\operatorname{tw}(\mathcal G)\ge v/16-1$.
Hlin\v en\'y--Whittle~\cite[Theorems 3.2 and 4.2]{matroidwidth}
give $\operatorname{tw}(M(\mathcal G))=\operatorname{tw}(\mathcal G)$ and
$\operatorname{tw}(M)\le2\operatorname{bw}(M)-2$.
Here a coordinate tree initially has singleton leaves. Deleting duplicate edge coordinates produces $M(\mathcal G)$, and
matroid treewidth does not increase under taking minors (their
Proposition~3.1).  Hence every coordinate tree has a cut with
$\lambda(J)\ge v/32+1/2$.

Factor $K_0=FF^{\mathsf T}$ with $v-2$ columns.  Removing the constant
direction from $\operatorname{range}H$ decreases each restricted rank
by at most one.  At the above cut,
\begin{align*}
 d_J&:=\dim(\operatorname{row}F_J\cap\operatorname{row}F_{J^c})\\
 &=\operatorname{rank}F_J+\operatorname{rank}F_{J^c}-(v-2)
   \ge\lambda(J)-2\ge v/32-3/2.
\end{align*}
Thus $\operatorname{rank}(K_0)_{J,J^c}\ge d_J=\Omega(n)$.
This is coupling to the entire complement.  In a balanced tree of
depth $O(\log n)$, concatenating the ancestor sibling interfaces
forms this external block; at least one sibling interface therefore
has exact rank $\Omega(n/\log n)$. The same conclusion applies to trees with leaf blocks of size $o(n)$: refine each leaf into singletons; a cut inside one such leaf has rank $o(n)$ and cannot be the wide cut above.

\paragraph{The interfaces remain strong at the actual ridge scale.}
For $K=\alpha I+FF^{\mathsf T}$ put
$S_J=F_J^{\mathsf T}F_J$ and $R_J=S_J(S_J+\alpha I)^{-1}$.
The canonical cut-coupling singular values are the square roots of the
eigenvalues of $R_J^{1/2}R_{J^c}R_J^{1/2}$.  Write $F=HZ$, where
$Z$ maps onto the nontrivial vertex space and its squared singular
values lie in $[a/n,b/n]$ for absolute $a,b>0$; this follows explicitly
from~\eqref{ifc:interface:range}.

For any cut, $H_J^{\mathsf T}H_J$ is sign-conjugate to an unweighted
multigraph Laplacian.  Its pseudoinverse trace is at most $v^2/2$:
effective resistances in a component are bounded by spanning-tree
distances, and their sum is the component size times the pseudoinverse
trace.  Therefore at most $\delta v^2/2$ positive eigenvalues are below
$\delta$.  Restricting away the known bipartition kernel and the
constant direction adds at most one to this count by interlacing.
Congruence by $Z$ then shows that $S_J$ has at most
$1+\delta v^2/2$ positive eigenvalues below $a\delta/n$.

Choose $\delta=1/(128v)$.  The large-eigenvalue subspaces on the two
sides intersect in dimension at least $3v/128-4$.
On either such subspace $R_J\succeq(1-\epsilon)I$, where
$\epsilon\le128\alpha nv/a$.  Eigenvalue monotonicity for positive
semidefinite products gives at least $3v/128-4$ canonical singular
values at least $1-\epsilon$.
The prescribed parameters imply
\[
 \alpha=\eta/\rho
 \le\frac{K^2\eta^3}{512C\mu\,256^2}
 \le\frac{q\eta^2}{2^{18}n}
 =\frac{q}{2^{42}J^2n^3},
\]
using $K^2\le128C(\mu/\eta)q/n$.
Here $q=O(\log n)$ because $m=n/16$.
For every fixed threshold less than one and sufficiently large $n$,
every coordinate tree thus has an external cut with $\Omega(n)$ strong
canonical couplings.

\paragraph{A cheap global decomposition is still available.}
The projector onto $U$ is
\[
 \Pi_U=H(H^{\mathsf T}H)^\dagger H^{\mathsf T}
                         -\mathbf1\mathbf1^{\mathsf T}/n.
\]
The only kernel of $H^{\mathsf T}H$ is the known bipartition vector;
on its complement the condition number is at most four.
Thus $\Pi_U$ can be applied to inverse-polynomial accuracy in
$\widetilde O(n)$ using the sparse incidence matrix and a
constant-conditioned vertex-space iteration.
By~\eqref{ifc:interface:range}, $\alpha I+(\gamma/n)\Pi_U$ is a
constant-quality preconditioner for suitable fixed $\gamma>0$.
Its inverse acts explicitly on $U$ and $U^\perp$; the necessary
accuracy, including the amplification by $\alpha^{-1}$, costs
logarithmic factors.  Wide coordinate interfaces therefore coexist
with an inexpensive solver on this family.  A general strategy needs
additional trajectory information or a structured global interface.

Numerical checks use five graphs, full retained-row dimensions
$16384\le n\le524288$, and seventeen cuts.  At $n=524288$ the
copy-halves cut has $1022>\sqrt n$ canonical singular values, each
exceeding $0.9999999999987$, even at the larger ridge $\alpha=n^{-3}$.
The checks use an exact equitable reduction of the actual formulas;
they neither sample signing trajectories nor establish reachability.


\subsection{Two first-order formulations that do not certify a signing}
\label{new:fw}
The following exact examples limit two natural replacements of the whole
walk by a Frank--Wolfe iteration. They do not rule out other first-order
algorithms. A linear minimization oracle is understood to return an exact
optimizer of the current linearization.

\paragraph{A discrepancy penalty minus quadratic energy.}
Let $k$ be a positive even integer, $n=k^2$, and
$A=k^{-1}(1,\ldots,1)$. Its column norms are at most one. Consider
\[
 F(x)=\tfrac12\|Ax\|^2-\|x\|^2,\qquad x\in[-1,1]^n.
\]
At $x=\one$, $\nabla F(x)=-\one$, so the cube's linear minimization
oracle returns $x$ itself. The Frank--Wolfe gap is zero, but
$\|Ax\|_\infty=k=\sqrt n$. Every balanced signing has discrepancy zero
and objective $-n$, compared with $-n/2$ at this stationary vertex.
Moreover, from $x=t\one$, $0<t<1$, the oracle returns $\one$ and exact
line search reaches that vertex in one step. Stationarity for this
nonconvex objective therefore supplies no dimension-free signing bound.

\paragraph{Quadratic energy inside an exactly feasible polytope.}
Fix $c\ge1$ and even $n\ge8c^2$, and put
\[
 A=\tfrac12I+\frac{2c-\tfrac12}{n}\one\one^T,
 \qquad K_c=\{x\in[-1,1]^n:\|Ax\|_\infty\le c\}.
\]
Each column has squared norm
$\tfrac14(1-1/n)+4c^2/n\le3/4$. The matrix is invertible, with
eigenvalues $2c$ on constants and $1/2$ on their orthogonal complement.
At $\bar x=\one/2$, one has $A\bar x=c\one$.
For every $y\in K_c$, summing the upper row constraints gives
$\sum_jy_j\le n/2$. Equality forces every upper constraint to be tight,
and invertibility then forces $y=\bar x$. Thus $\bar x$ is the unique
maximizer of the linearization of $\|x\|^2$ at $\bar x$.
The exact feasible Frank--Wolfe energy iteration is fixed there, although
every coordinate is fractional and its defect is $3n/4$.
A balanced signing $\sigma$ satisfies $A\sigma=\sigma/2$ and belongs to
$K_c$. Exact line search from $t\one$, $0<t<1/2$, reaches the fractional
fixed point in one step. Even a free exact linear oracle does not repair
this formulation; constructing that oracle would require separate cost
accounting as well.

\section{Verification and reproducibility}\label{v:numerics}
\suppressfloats[t]
The accompanying packet contains the standalone source, compiled paper,
scripts, and stored numerical records. Its \texttt{README.txt} maps each
test family to the relevant statement and records how to run the checks.
The proofs above establish the theorems. The component computations check algebra,
local constructions, and implementation formulas. The packet also retains
the previously completed endpoint runs summarized in
Section~\ref{ep:evidence}, together with their certificate and local-state
checks. Those runs use the endpoint wrapper rather than the secant
controller. The optimized-profile parameter checks and matched endpoint experiments
are described in Section~\ref{opt:comparison}. None of these computations
validates the unproved production
hypothesis~\eqref{d:conjecture}.

\begin{table}[t]
\centering\small
\caption{Scope of the component checks retained in the packet.}
\label{new:verification-table}
\begin{tabular}{@{}L{0.25\textwidth}L{0.43\textwidth}L{0.24\textwidth}@{}}
\toprule
Component & Checks retained in the packet & Interpretation\\
\midrule
Secant decrement & 1,200 noncommuting triples, 600 independent congruence
tests, and 320 factor-error comparisons & Local algebra and robustness\\[3pt]
Controller accounting & 56 synthetic paths; 3,360 transitions and
50,690 corrections & Signed ledger identities; dense witness tests\\[3pt]
Actual row formulas & A $10,000$-coordinate symmetry-reduced derivative
test and high-precision block-state replays & Supplied local states;
reachability unproved\\[3pt]
Extra regularization & 240 matrix tests, 4,672 sign patterns, and exact
determinantal subset checks & Curvature allowance and inverse identities\\[3pt]
Schur and tree algebra & Exact rational overlap, elimination, and
deflation tests; independently changing tree systems & Supplied-factor
solvers and error composition\\[3pt]
Constructed interfaces & 480 solve states with 2,272 signed updates;
16 supported-profile cases & Construction, rank guards, and direct solves\\[3pt]
New sampler and Perron options & Scalar and rotated systems, a
near-cancellation trace fixture, positive $2\times2$ Perron matrices &
Fixed-polynomial recurrence, constant tolerance, and gap dependence\\[3pt]
Certified endpoints & Twelve complete low-branch runs through $n=512$;
positive-certificate and protected-row checks & Stored floating-point
traces; scope in Section~\ref{ep:evidence}\\
\bottomrule
\end{tabular}
\end{table}

\paragraph{Secants and regularization.}
The minimum observed robust analytic decrement in the 1,200-case check
was $0.0143576$, above the proved $\gamma_0=0.0121225\ldots$; the largest
Schur-identity error was $6.84\cdot10^{-14}$. Independent congruence tests
included ideal-target condition numbers up to $10^6$.
Controller tests used dense generalized eigendecomposition instead of the
specified randomized witness routine. Their largest signed-ledger residual
was $3.72\cdot10^{-11}$. The rotation and alternating-projector models
exercise the proved examples and counterexamples of Appendix~\ref{c:models}.
They do not sample the signing algorithm's probability law.
The regularization records also include 96 exact rational cases enumerating
720 determinant subsets, 600 floating-point matrix cases, and separate
fixed-reference and moving-reference checks.

\paragraph{Actual local states and the scale of a run from zero.}
The derivative fixture uses unit-column circulant data with an exact
$20$-dimensional symmetry reduction. It has $L=1.005$ and satisfies the
local row barriers and scale invariants. Its observed
$\|T'[h]\|_F=2.86864\cdot10^{-4}$ is below the proved local bound
$0.168393$; central differences agree to $7\cdot10^{-13}$.
This is a supplied state and direction, not a reachability result.

The block replays use $4,096$ columns and $130$ identical normalized rows
per block. Stored runs have $16,384$, $32,768$, or $65,536$ coordinates,
three to eight moves, and 18 to 60 repairs. They use particular balanced
Walsh outcomes, high-precision invariant-coordinate calculations, and a
filter proxy with a certified negligible error. A dense witness replaces
the randomized test; cleanup and phase resets are not executed. The first
three $16,384$-coordinate moves agree at all 45 saved significant digits
when precision is raised from 120 to 150 decimal digits. Endpoint
expectations in these symmetric fixtures do not average over the sampler.

The small prescribed step constants make short runs from zero uninformative
about the high branch. Before a coordinate freezes, if $\|x\|<1$, the row
identities give $u_{i\sigma}\le22-75/r_i$. Column normalization and
$r_i\le3/2$ yield
\[
 L\le\max\{1/128,(256/9)e^{-8}\}+\eta s<0.00956
\]
at the tested dimensions. Orthogonality gives $\|x_T\|^2\le T\tau^2$.
For $s=16,384$, the prescribed $K=1$ and $\tau=2^{-50}$ keep this bound
valid for more than $1.25\cdot10^{30}$ initial steps. Consequently no
short numerical prefix at the original constants tests the difficult
spectral regime of the complete run.

\paragraph{Construction and deflation.}
The Schur tests include 270 noncommuting overlap identities, 270 signed
rank-one updates, 18 exact rational instances, and tiny-ridge fixtures.
Tree checks include exact hereditary rank and normalized deflation,
72 filter-rank cases, and twelve changing systems at each order
$16,32,64,128$. The signed-stream experiments discover their initial
factors from dense matrices and then use their supplied update vectors;
a rank-growth fixture is correctly rejected. Supported-profile checks
construct $M,T,P,E,G$ from row data, including dependent constraint rows.
The largest saved relative Gram error is below $1.6\cdot10^{-12}$ and the
largest solve residual below $2.4\cdot10^{-12}$. Floating implementations
use thin QR and small-core SVD; the arithmetic proofs use exact rank-profile
elimination. The wide-interface fixture uses an equitable reduction of
the actual canonical formulas and makes no reachability claim.

\paragraph{Checks added in this integration.}
The fixed Chebyshev recurrence was compared with its covariance formula
on rotated feasible subspaces at enclosure ratios about $8.8$, $852$,
and $72,932$. Scalar systems verify the one-step case. A trace fixture
with $C=1200$ checks the tolerance when positive and negative curvature
nearly cancel. The acceptance constants $11/28$ and $15/56$ were checked
in rational arithmetic. Power iteration on $M_\eta$ of
Appendix~\ref{new:power}, for $\eta=10^{-2},10^{-3},10^{-4}$, required
562, 5,542, and 55,338 iterations respectively for angle tangent at most
$10^{-6}$, consistent with the proved inverse-gap dependence. The two
first-order obstructions in Appendix~\ref{new:fw} were checked by exact
identities and, for the feasible-polytope example, independent linear
programs. None of these tests changes the stated scope of the theorems.

\begin{thebibliography}{99}
\bibitem{Final37} \emph{Constructive Koml\'os discrepancy below 37.54:
normalized response and joint dimension bounds}, companion addendum,
\nolinkurl{archive/komlos_constructive_final.tex} in this repository,
28 September 2026.

\bibitem{Screened67} \emph{Constructive Koml\'os discrepancy below 67:
coupled spectral signing with input-linear row processing}, screened
finite core, 27 September 2026. Frozen source:
\nolinkurl{archive/komlos_constructive_below67.tex} in this repository.

\bibitem{Below67} \emph{Constructive Koml\'os discrepancy below 67:
a coupled spectral estimate, an optimized coordinate profile, and rounding
from the remaining error allowance}, companion supplement,
\nolinkurl{komlos_constructive_below67.tex}, 28 September 2026.

\bibitem{LiSpectral} X. Li, \emph{Fast Spectral Signing for Vector Balancing},
arXiv:2609.30044v1, 24 September 2026, Lemma 2.4.
\url{https://arxiv.org/html/2609.30044v1}.

\bibitem{CertSteps} \emph{Certified long steps for the retained-row walk},
unpublished implementation note, updated 28 September 2026.

\bibitem{CertAnalysis} \emph{Certified steps for constructive Koml\'os:
analysis and experimental audit}, companion analysis and stored
implementation records, 24 September 2026.

\bibitem{GFL} S. Guo, E. X. Fang, and J. Lu,
\emph{Vector Balancing in Polynomial Time}, arXiv:2609.23540 (2026).
\url{https://arxiv.org/abs/2609.23540}.

\bibitem{Sob} A. Sobczyk,
\emph{Deterministic complexity analysis of Hermitian eigenproblems},
arXiv:2410.21550v2 (2025), Corollaries 2.1 and C.2.
\url{https://arxiv.org/abs/2410.21550}.

\bibitem{DPS} J.-G. Dumas, C. Pernet, and Z. Sultan,
\emph{Simultaneous computation of the row and column rank profiles},
arXiv:1301.4438 (2013), Algorithm 1, Corollary 3, and Theorem 2.
\url{https://arxiv.org/abs/1301.4438}.

\bibitem{MM} E. Dupont, M. Eisenberger, B. Kozlovskii, A. Mehrabian,
F. J. R. Ruiz, A. See, R. Zhou, J. Alman, V. Vassilevska Williams,
and M. Balog, \emph{Improving the matrix multiplication exponent with
modern optimization and AlphaEvolve}, arXiv:2608.16884 (2026).
\url{https://arxiv.org/abs/2608.16884}.

\bibitem{AJSS} A. Ahmadinejad, A. Jambulapati, A. Saberi, and A. Sidford,
\emph{Perron--Frobenius Theory in Nearly Linear Time: Positive Eigenvectors,
M-matrices, Graph Kernels, and Other Applications}, arXiv:1810.02348 (2018),
Theorems 8--9, Section 6.3, and Appendix 10.
\url{https://arxiv.org/abs/1810.02348}.

\bibitem{Sra} S. Sra, \emph{Positive definite matrices and the S-divergence},
arXiv:1110.1773v4 (2013).
\url{https://arxiv.org/abs/1110.1773}.

\bibitem{Gower} R. M. Gower, F. Hanzely, P. Richt\'arik, and S. Stich,
\emph{Accelerated Stochastic Matrix Inversion: General Theory and Speeding
up BFGS Rules for Faster Second-Order Optimization},
arXiv:1802.04079v2 (2018).
\url{https://arxiv.org/abs/1802.04079}.

\bibitem{LSmaint} Y. T. Lee and A. Sidford,
\emph{Efficient Inverse Maintenance and Faster Algorithms for Linear Programming},
arXiv:1503.01752 (2015). \url{https://arxiv.org/abs/1503.01752}.

\bibitem{JLSW} H. Jiang, Y. T. Lee, Z. Song, and S. C.-w. Wong,
\emph{An Improved Cutting Plane Method for Convex Optimization,
Convex-Concave Games and its Applications}, arXiv:2004.04250 (2020).
\url{https://arxiv.org/abs/2004.04250}.

\bibitem{HJSTZ} B. Huang, S. Jiang, Z. Song, R. Tao, and R. Zhang,
\emph{Solving SDP Faster: A Robust IPM Framework and Efficient Implementation},
arXiv:2101.08208 (2021), Theorem 10.8.
\url{https://arxiv.org/abs/2101.08208}.

\bibitem{DMdpp}
M.~Derezi\'nski and M.~W.~Mahoney,
\emph{Determinantal point processes in randomized numerical linear algebra},
arXiv:2005.03185 (2020).
\url{https://arxiv.org/abs/2005.03185}.

\bibitem{DLMdesign}
M.~Derezi\'nski, F.~Liang, and M.~W.~Mahoney,
\emph{Bayesian experimental design using regularized determinantal point processes},
Proceedings of AISTATS, PMLR \textbf{108} (2020), 3197--3207.
\url{https://proceedings.mlr.press/v108/derezinski20a.html}.

\bibitem{Saad} Y. Saad, \emph{Iterative Methods for Sparse Linear Systems},
second edition, SIAM (2003), Sections 6.11.3 and 9.2.
\href{https://www-users.cse.umn.edu/~saad/IterMethBook_2ndEd.pdf}{Author-hosted text}.

\bibitem{GR16} R. M. Gower and P. Richt\'arik,
\emph{Randomized Quasi-Newton Updates are Linearly Convergent Matrix
Inversion Algorithms},
\href{https://arxiv.org/abs/1602.01768}{arXiv:1602.01768} (2016).

\bibitem{ulv} S. Chandrasekaran, M. Gu, and T. Pals,
\emph{A Fast ULV Decomposition Solver for Hierarchically Semiseparable
Representations}, SIAM J. Matrix Anal. Appl. 28 (2006), 603--622.
\href{https://doi.org/10.1137/S0895479803436652}{doi:10.1137/S0895479803436652}.

\bibitem{hierarchical} S. Ambikasaran, M. O'Neil, and K. R. Singh,
\emph{Fast symmetric factorization of hierarchical matrices with applications},
\href{https://arxiv.org/abs/1405.0223}{arXiv:1405.0223v2} (2016).

\bibitem{compression} P.-G. Martinsson,
\emph{Compressing rank-structured matrices via randomized sampling},
\href{https://arxiv.org/abs/1503.07152}{arXiv:1503.07152} (2015).

\bibitem{ADGH2018} E. Agullo, E. Darve, L. Giraud, and Y. Harness,
\emph{Low-rank factorizations in data sparse hierarchical algorithms for
preconditioning symmetric positive definite matrices},
SIAM J. Matrix Anal. Appl. 39 (2018), 1701--1725.

\bibitem{FHY2020} J. Feliu-Fab\`a, K. L. Ho, and L. Ying,
\emph{Recursively preconditioned hierarchical interpolative factorization
for elliptic partial differential equations}, Commun. Math. Sci. 18
(2020), 91--108. \href{https://arxiv.org/abs/1808.01364}{arXiv:1808.01364}.

\bibitem{levitt} J. Levitt and P.-G. Martinsson,
\emph{Linear-Complexity Black-Box Randomized Compression of Rank-Structured
Matrices}, SIAM J. Sci. Comput. 46 (2024), A1747--A1763.
\href{https://arxiv.org/abs/2205.02990}{arXiv:2205.02990v3}.

\bibitem{hss2026} N. Amsel, T. Chen, F. Duman Keles, D. Halikias,
C. Musco, C. Musco, and D. Persson,
\emph{Quasi-optimal Hierarchically Semi-separable Matrix Approximation},
\href{https://arxiv.org/abs/2505.16937}{arXiv:2505.16937v2} (2025).

\bibitem{hodlr2025} T. Chen, F. Duman Keles, D. Halikias,
C. Musco, C. Musco, and D. Persson,
\emph{Near-optimal hierarchical matrix approximation from matrix-vector
products}, SODA (2025), 2656--2692.
\href{https://arxiv.org/abs/2407.04686}{arXiv:2407.04686v2}.

\bibitem{matroidwidth} P. Hlin\v en\'y and G. Whittle,
\emph{Matroid tree-width}, European J. Combin. 27 (2006), 1117--1128.
\href{https://www.fi.muni.cz/~hlineny/papers/matr-tw-final.pdf}{Author manuscript}.

\bibitem{bordenave} C. Bordenave,
\emph{A new proof of Friedman's second eigenvalue theorem and its extension
to random lifts}, \href{https://arxiv.org/abs/1502.04482}{arXiv:1502.04482v4}
(2019).

\bibitem{SaadEig} Y. Saad,
\emph{Numerical Methods for Large Eigenvalue Problems}, second edition,
SIAM (2011), Sections 4.1 and 7.4.
\url{https://www-users.cse.umn.edu/~saad/eig_book_2ndEd.pdf}.

\bibitem{ChebSample} M. Pereira and N. Desassis,
\emph{Efficient simulation of Gaussian Markov random fields by Chebyshev
polynomial approximation}, Spatial Statistics 31 (2019), 100359.
\url{https://arxiv.org/abs/1805.07423}.
\end{thebibliography}
\end{document}
